\documentclass[11pt]{amsart}

\usepackage[T1]{fontenc}
\usepackage[utf8]{inputenc}
\usepackage{lmodern}
\usepackage[a4paper,margin=1.05in]{geometry}
\usepackage{microtype}
\usepackage{amsmath,amssymb,amsthm,mathtools,mathrsfs}
\usepackage{enumitem}
\usepackage{booktabs}
\usepackage[dvipsnames]{xcolor}
\usepackage[colorlinks=true,linkcolor=MidnightBlue,citecolor=ForestGreen,urlcolor=BrickRed,
  pdftitle={Split Support on the Absolute Arithmetic Curve, Version 5},
  pdfauthor={The Clankers},
  pdfsubject={Split-zero geometry, distributive-lattice support, arithmetic Jacobians, Frobenius perfection, finite gerbe anomalies, theta quantization, UHF limits, and nonassociative coefficient algebras}]{hyperref}
\usepackage{array}
\usepackage{longtable}
\usepackage{tikz-cd}

\numberwithin{equation}{section}

\newtheorem{theorem}{Theorem}[section]
\newtheorem{proposition}[theorem]{Proposition}
\newtheorem{corollary}[theorem]{Corollary}
\newtheorem{lemma}[theorem]{Lemma}
\newtheorem{conjecture}[theorem]{Conjecture}
\theoremstyle{definition}
\newtheorem{definition}[theorem]{Definition}
\newtheorem{example}[theorem]{Example}
\newtheorem{remark}[theorem]{Remark}
\newtheorem{warning}[theorem]{Warning}

\newcommand{\B}{\mathbf B}
\newcommand{\Z}{\mathbf Z}
\newcommand{\Q}{\mathbf Q}
\newcommand{\R}{\mathbf R}
\newcommand{\C}{\mathbf C}
\newcommand{\N}{\mathbf N}
\newcommand{\F}{\mathbf F}
\newcommand{\Fone}{\mathbf F_1}
\newcommand{\cA}{\mathcal A}
\newcommand{\cH}{\mathcal H}
\newcommand{\cR}{\mathcal R}
\newcommand{\cL}{\mathcal L}
\newcommand{\Frob}{\operatorname{Frob}}
\newcommand{\dist}{\operatorname{dist}}
\newcommand{\Clop}{\operatorname{Clop}}
\newcommand{\No}{\mathbf{No}}
\newcommand{\Spec}{\operatorname{Spec}}
\newcommand{\Proj}{\operatorname{Proj}}
\newcommand{\supp}{\operatorname{supp}}
\newcommand{\Supp}{\operatorname{Supp}}
\newcommand{\im}{\operatorname{Im}}
\newcommand{\kerp}{\operatorname{Ker}}
\newcommand{\Hom}{\operatorname{Hom}}
\newcommand{\End}{\operatorname{End}}
\newcommand{\Aut}{\operatorname{Aut}}
\newcommand{\GL}{\operatorname{GL}}
\newcommand{\PGL}{\operatorname{PGL}}
\newcommand{\Frac}{\operatorname{Frac}}
\newcommand{\ord}{\operatorname{ord}}
\newcommand{\st}{\operatorname{st}}
\newcommand{\gr}{\operatorname{gr}}
\newcommand{\cO}{\mathcal O}
\newcommand{\cM}{\mathcal M}
\newcommand{\cU}{\mathcal U}
\newcommand{\zero}{\mathbf 0}
\newcommand{\Pum}{\mathbf P_{\mathrm{um}}}
\newcommand{\Sp}{\operatorname{Sp}}
\newcommand{\id}{\operatorname{id}}
\newcommand{\colim}{\operatorname*{colim}}
\newcommand{\liminv}{\operatorname*{lim}}
\newcommand{\eps}{\varepsilon}
\newcommand{\taup}{\tau}
\newcommand{\eR}{e}
\newcommand{\ideal}[1]{\mathfrak #1}
\newcommand{\Picf}{\operatorname{Pic}_{\mathrm f}}
\newcommand{\Picar}{\operatorname{Pic}_{\mathrm{ar}}}
\newcommand{\Jacar}{\operatorname{Jac}_{\mathrm{ar}}}
\newcommand{\Idem}{\operatorname{Idem}}
\newcommand{\Ann}{\operatorname{Ann}}
\newcommand{\Spl}{\operatorname{Spl}}
\newcommand{\Nbar}{\overline{\N}}
\newcommand{\Ext}{\operatorname{Ext}}
\newcommand{\Tor}{\operatorname{Tor}}
\newcommand{\Gal}{\operatorname{Gal}}
\newcommand{\Heis}{\operatorname{Heis}}
\newcommand{\Trg}{\operatorname{Trg}}
\newcommand{\CD}{\operatorname{CD}}
\newcommand{\Mat}{\operatorname{Mat}}
\newcommand{\Vect}{\operatorname{Vect}}
\newcommand{\LagCorr}{\operatorname{LagCorr}}
\newcommand{\UHF}{\mathcal Q_{\Omega}}
\newcommand{\Ad}{\operatorname{Ad}}
\newcommand{\EM}{\operatorname{EM}}

\DeclareMathOperator{\Bool}{Bool}
\DeclareMathOperator{\Mon}{Mon}
\DeclareMathOperator{\Ring}{Ring}
\DeclareMathOperator{\SemiRing}{SemiRing}

\title[Split Support on the Absolute Arithmetic Curve]
{Split Support on the Absolute Arithmetic Curve\\[4pt]
\large Arithmetic Jacobians, Frobenius Perfection, Gerbe Quantization, and Lattice-Valued Support Geometry}
\author{The Clankers}
\date{June 2026\\Pre-publication draft, Version 5}

\begin{document}

\begin{abstract}
Let $G(R)=R\sqcup\{\tau\}$ be the split-zero globalization of a
nonzero commutative ring, with actual semiring zero $\tau$ and distinct
supported ring zero $e=0_R$.  We develop the arithmetic geometry forced
by this distinction in the language of semiring spectra, the
Connes--Consani absolute $\mathbf F_1$-curve, arithmetic Jacobian monoids,
Frobenius-perfect denominator towers, finite gerbe anomalies, theta
quantization, and characteristic-one sites.

For a bounded distributive lattice $L$ we construct
\[
G_L(R)=\{(0,\lambda):\lambda\in L\}\cup
       \{(r,1_L):r\in R\}\subseteq R\times L.
\]
Its universal ring reflection is projection to $R$, its ring-null fibre is
$(L,\vee,\wedge)$, and localization at $(0,\lambda)$ is the principal
interval $\downarrow\lambda$.  Boolean lattices recover independent mixed
support; finite chains recover iterated sub-zero towers and occur as the
support lattices of $\mathbf Z/p^n\mathbf Z$ and
$k[\epsilon]/(\epsilon^n)$.  Ordinary ideal multiplication yields a
second, valuation-sensitive truncated tropical semiring.  The construction
extends to nonassociative coefficients: it preserves the associator and
sends a genuine zero-divisor product to supported zero rather than absence.
For signed graded algebras the basis associator is the normalized
$3$-cocycle obtained from the multiplication $2$-cochain.  In the real
Cayley--Dickson tower, octonionic nonassociativity and sedenionic zero
divisors are therefore distinct support events.

The scalar theory includes
$G(A\times B)\cong G(A)\times_{\mathbf B}G(B)$, the fixed-locus formula
$\operatorname{Fix}_{G(R)}(u)=\{\tau\}\sqcup\operatorname{Ann}_R(u-1)$,
the division trichotomy, Boolean-cubical ring-null halos, mixed
localization, projective stratifications, projective-semimodule
constraints, conservative ideal arithmetic, and residue monads.  The
normed ideal class group and Dedekind zeta function are unchanged, whereas
the support-separated weight $N(\mathfrak a)+1$ is nonmultiplicative.
Hyperreal and surcomplex realizations separate absence, exact zero,
valuation order, and complex leading phase.

For the completed arithmetic curve, the Connes--Consani theorem becomes
\[
\operatorname{Jac}_{\mathrm{ar}}
\cong\operatorname{Pic}_{\mathrm f}\times\mathbf B,
\qquad
\operatorname{Idem}(\operatorname{Jac}_{\mathrm{ar}})
\cong\mathcal P(\mathbb P_f\sqcup\{\infty\}).
\]
The Boolean coordinate is exactly the support truncation of the possibly
zero archimedean seminorm.  The split spectrum resolves the zero
Abel--Jacobi fibre by
$\{0\}\hookrightarrow\mathbf Z\hookrightarrow\mathbf Z[1/p]$.
Rigidified denominator points carry universal Frobenius transitions, and
the unique globally Frobenius-perfect point is $H_\Omega=\mathbf Q$ with
stalk $\mathbf F_1[T^{\mathbf Q_+}]$.  Its carry and character-dual
extensions are
\[
0\to\mathbf Z\to\mathbf Q\to\mathbf Q/\mathbf Z\to0,
\qquad
0\to\widehat{\mathbf Z}\to\mathbb A_f\to\mathbf Q/\mathbf Z\to0.
\]

For every $n$, the phase group
$\Phi_n=n^{-1}\mathbf Z/\mathbf Z$ carries an explicit nondegenerate
homogeneous quadratic refinement.  Eilenberg--Mac Lane delooping and the
Freed--Hopkins--Lurie--Teleman finite path integral produce an invertible
four-dimensional gerbe theory.  Surface transgression gives the finite
Heisenberg extension with twisted group algebra $M_n(\mathbf C)$, and the
winding-$n$ theta sector of the explicit level-one bundle gerbe on $T^3$
is the same irreducible module, uniquely up to scalar.  Frobenius
multiplication is loop-winding dilation, while denominator divisibility is
encoded by cyclotomic-equivariant Lagrangian correspondences.

A strict matrix model has inductive limit the universal UHF algebra
$\mathcal Q_\Omega$, with
\[
K_0(\mathcal Q_\Omega)\cong(\mathbf Q,\mathbf Q_{\ge0},1),
\qquad
\mathcal D_\Omega\cong C(\widehat{\mathbf Z})
\cong C^*(\mathbf Q/\mathbf Z).
\]
The cyclotomic group acts canonically on the correspondence system and
profinite diagonal, but the naive permutation operators fail to commute
with the standard full-matrix embeddings; the defect is an explicit
crossed carry cocycle.  Thus the Morita-level system is canonical, whereas
a strict lift requires coherent carry trivializations.  Split support also
distinguishes an absent anomalous tier from a defined vanishing amplitude.

The reciprocal-root skeleton retains its Beurling realization and the
estimate $x_{\Lambda_Q}\le d_Q^2$; the remaining density statement is the
classical Nyman--B\'aez-Duarte RH-equivalent endpoint.  No
Scholze-perfectoid identification, fully extended infinite-level TQFT, or
proof of the Riemann hypothesis is asserted.
\end{abstract}

\maketitle
\tableofcontents

\section{Introduction and statement of the results}
\label{sec:introduction}

The split-zero globalization of a nonzero commutative ring $R$ is the
commutative semiring
\[
G(R)=R\sqcup\{\tau\},
\]
in which the operations on $R$ are unchanged, $\tau$ is the additive
identity and multiplicative absorber, and
\[
\tau=0_{G(R)}\neq e:=0_R\in R\subseteq G(R).
\]
Thus $\tau$ records absence of support, whereas $e$ records a supported
arithmetic value equal to zero.  The foundational notes
\cite{Globalization,Secondary} establish the ring reflection, the unique
Boolean support character, the ideal and prime spectra, congruences,
localizations, the classification of semimodules by diagrams over
join-semilattices, and the Boolean-cubical support of free semimodules.

The present paper develops the geometry forced by that distinction and then
places it in the language of the Connes--Consani absolute arithmetic curve
\cite{CCAbsolute}.  The result is not a dictionary of suggestive terms.
Every new bridge is formulated as an explicit map, localization,
colimit, line bundle, action groupoid, or sheaf construction, and the
logical dependence on imported results is stated separately in Section
\ref{sec:claim-boundaries}.

The principal results are as follows.  The statements are organized by
logical dependence rather than chronology: the split algebra is developed
first, the absolute and arithmetic-site geometry second, and the analytic
realizations only after the relevant algebraic and categorical structures
have been fixed.

\begin{enumerate}[label=\textup{(\arabic*)},leftmargin=3.2em]
\item \textbf{Synchronized support and fixed loci.}
For rings $A,B$,
\[
G(A\times B)\cong G(A)\times_{\B}G(B),
\]
so mixed support occurs precisely in the Cartesian product outside this
fibre product.  For every supported unit $u$,
$\operatorname{Fix}_{G(R)}(u)=\{\tau\}\sqcup\Ann_R(u-1)$.

\item \textbf{Projective and ideal layers.}
A finitely generated projective $G(R)$-semimodule has finite distributive
support, projective amplitude fibres, and split transport.  Over a Dedekind
domain, the normed ideal class group and Dedekind zeta function are
preserved, but the separated weight $N(\mathfrak a)+1$ is
nonmultiplicative and therefore non-Eulerian.
\item \textbf{Division and absolute infinitesimals.}
For the supported zero and actual semiring zero,
\[
G(R)[e^{-1}]\cong\B,
\qquad
G(R)[\tau^{-1}]\cong\zero.
\]
Ordinary localization away from $0_R$ remains split arithmetic.  The kernel
of the universal ring reflection is $\{\tau,e\}$, and in $G(R)^I$ the
ring-null halo is canonically the Boolean cube $\B^I$.

\item \textbf{Mixed localization and scoped identities.}
For $S=G(K)$ over a field and $x\in S^I$, coordinates with nonzero
amplitude retain $S$, coordinates equal to $e$ become Boolean, and
coordinates equal to $\tau$ become terminal:
\[
S^I[x^{-1}]\cong S^{U(x)}\times\B^{Z(x)}.
\]
Imposing $e=h_i\in K^\times$ only on a face $A\subseteq I$ gives
\[
S^I/(\varepsilon_A=h_A)
\cong\B^A\times S^{I\setminus A}.
\]
This is the precise sense in which an identity may be adjoined locally
without collapsing the complementary geometry.

\item \textbf{Distributive-lattice support and nonassociative coefficients.}
For every bounded distributive lattice $L$ there is a canonical
lattice-split globalization $G_L(R)$ whose ring-null fibre is $L$ and whose
localization at a support idempotent is the corresponding principal
interval.  Boolean lattices recover independent support cubes; chains
recover finite sub-zero towers and the support lattices of truncated
nilpotent thickenings.  The same construction applies to nonassociative
algebras: it preserves the associator and separates intrinsic zero-divisor
cancellation from absence.  Signed graded multiplication is controlled by
a normalized $2$-cochain and its basis associator $3$-cocycle.

\item \textbf{Projective support stratification.}
For a field $K$, unimodular split projective space decomposes over the
nonempty Boolean faces:
\[
\Pum^n(G(K))
\cong
\bigsqcup_{\varnothing\neq A\subseteq\{0,\dots,n\}}
\mathbf P^{|A|-1}(K).
\]
Over $\F_q$ this gives a closed point-count formula and a rational zeta
function.  The invertible matrices are exactly monomial matrices, so
\[
\GL_{n+1}(G(K))\cong(K^\times)^{n+1}\rtimes S_{n+1}.
\]
The global split projective symmetry is therefore toric or $\Fone$-type,
not the full classical projective linear group.

\item \textbf{Leibniz monads and valuation fibres.}
For a residue map $A\twoheadrightarrow k$ with kernel $\mathfrak m$, the
supported fibre over $a\in k$ is the ordinary monad
$\widetilde a+\mathfrak m$, while split support adds a Boolean cube of
possible activations of absent coordinates.  Formal completion commutes
with split globalization.  Hyperreal monads and surcomplex valuation
monads are recovered as special cases.  In $\No[i]$, the complete local
datum is support, surreal valuation order, and complex leading phase;
$e$ is the supported valuation-infinite centre, not the nonzero
infinitesimal $\omega^{-1}$.

\item \textbf{The arithmetic Jacobian is an exact support product.}
Connes--Consani's completed arithmetic Jacobian satisfies
\[
\Jacar\cong\Picf\times\B,
\qquad
\Idem(\Jacar)\cong
\mathcal P(\mathbb P_f\sqcup\{\infty\}).
\]
The finite-prime configuration and the archimedean zero-seminorm bit are
therefore literal support coordinates.  Collapsing the degenerate stratum
gives the reduced pointed split monoid $\Picf^\tau$.

\item \textbf{A split refinement of the absolute zero fibre.}
The split spectrum
\[
P_\tau\subset P_0\subset P_p
\]
admits a point map to the arithmetic site
\[
P_\tau\longmapsto\{0\},
\qquad
P_0\longmapsto\Z,
\qquad
P_p\longmapsto\Z[1/p].
\]
After applying Connes--Consani's map $\kappa$ to arithmetic divisor
classes, the first two points coincide.  Before applying $\kappa$, they
remain nonisomorphic.  This resolves the zero Abel--Jacobi fibre at the
level of point classes and stalks.

\item \textbf{Rigidified denominator points and two denominator laws.}
Although $Q^{-1}\Z\cong\Z$ as ordered groups, the rigidified pairs
$(Q^{-1}\Z,\Z\hookrightarrow Q^{-1}\Z)$ are distinct.  Their morphisms are
classified exactly by divisibility, and the transition
$Q\mid M=mQ$ becomes $\operatorname{Frob}_m$ on the abstract stalk
$\Fone[T]$.  The framed arithmetic divisors satisfy
$H_Q\otimes H_M\cong H_{QM}$, while embedded root refinement satisfies
$H_Q+H_M=H_{\operatorname{lcm}(Q,M)}$.

\item \textbf{Global Frobenius perfection and emergent places.}
The iterated LCM tower has filtered colimit
\[
\varinjlim_jQ_j^{-1}\Z=\Q,
\qquad
\varinjlim_j\Fone[T^{Q_j^{-1}\N}]
=\Fone[T^{\Q_+}].
\]
This is the stalk at the maximal supernatural point
$\prod_pp^\infty$, and every Frobenius is invertible there.  Every finite
stage has no Connes--Consani local place, while the limit has all finite
places and the archimedean place.  The place functor therefore does not
preserve this filtered colimit.

\item \textbf{Cyclotomic carry and adelic phase history.}
The rigidification $\Z\subset\Q$ yields the nonsplit extension
\[
0\to\Z\to\Q\to\Q/\Z\to0,
\]
whose class is the profinite unit
\[
[\mathcal E_\Omega]=1_{\widehat\Z}
\in\Ext^1_{\Z}(\Q/\Z,\Z).
\]
Its finite restrictions are the generators of
$\Ext^1(\Z/Q\Z,\Z)$.  Character duality gives the second nonsplit
extension
\[
0\to\widehat\Z\to\mathbb A_f\to\Q/\Z\to0,
\]
and
$\varprojlim_Q\Hom(Q^{-1}\Z,\Q/\Z)\cong\mathbb A_f$.
Every finite dual is torsion, while the limiting dual is torsion free.
For saturated points $H_S=\Z[S^{-1}]$, relative phase and rooted phase are
complementary primary summands of $\Q/\Z$.  The resulting symmetry group
is $\widehat\Z^\times$, the cyclotomic symmetry group of the
Bost--Connes system.

\item \textbf{Gerbe, theta, and operator-algebraic quantization.}
The phase group $\Phi_n=n^{-1}\Z/\Z$ admits a nondegenerate homogeneous
quadratic refinement $q_n$.  It determines an FHLT finite gerbe anomaly
and, after surface transgression, a Heisenberg extension whose twisted
group algebra is $M_n(\C)$.  The winding-$n$ theta sector of the level-one
$T^3$ bundle gerbe is the corresponding irreducible module.  Frobenius is
loop-winding multiplication.  Lagrangian denominator correspondences
compose exactly.  A strict matrix model has universal UHF limit
$\mathcal Q_\Omega$ with $K_0=\Q$ and diagonal
$C(\widehat\Z)\cong C^*(\Q/\Z)$; the failure of strict cyclotomic
equivariance of the standard embeddings is the explicit carry cocycle.

\item \textbf{Beurling realization at the perfect point.}
The reciprocal monomials $T^{1/n}$ are realized by the Beurling atoms
$F_{1/n}$.  Exact denominator inclusion gives
\[
x_{\Lambda_Q}\leq d_Q^2,
\qquad
\Lambda_Q=\operatorname{lcm}(1,\dots,Q),
\]
which strengthens the projection decrement in \cite{Schultze}.  The
integer Nyman--B\'aez-Duarte criterion becomes a density theorem for the
reciprocal-root skeleton of $\Fone[T^{\Q_+}]$.  This is an exact
reformulation of the RH-level endpoint, not a proof of it.

\item \textbf{Frobenius local systems underlying the helix.}
On the periodic orbit
\[
C_p=\R_+^\times/p^\Z,
\]
the phase $e^{-itu}$ is a section of a flat line with holonomy $p^{-it}$.
The growing helix factor $p^{1/2-it}$ is its half-density twist, and the
critical phasor $p^{-1/2-it}$ is the inverse half-density twist.  The two
chiral phase lines are dual, so their determinant is canonically trivial.
This is the coordinate-free arithmetic-site content of the Lean-verified
multiplier and determinant identities in \cite{Lavery}.

\item \textbf{Boolean place cubes and univalent quotients.}
A local Weil torsor $\cM_H^v$ acquires two fixed support states: absence and
the supported neutral morphism.  Its action groupoid is
\[
[\widetilde\cM_H^v/W_v]
\simeq BW_v\sqcup BW_v\sqcup *.
\]
For $H_p=\Z[1/p]$, the $p$-adic and archimedean directions form a
nine-stratum refinement over the Boolean square $\B^2$.  The classifying
components retain isotropy data that a coarse orbit set would erase.

\item \textbf{The split arithmetic and scaling sites.}
For a semiring object $A$ in a Grothendieck topos, the internal coproduct
$1\amalg A$ is its split globalization.  Applied to the original arithmetic
site it gives
\[
(\widehat{\N^\times},\Spl(\Nbar)),
\]
which separates absence from the supported tropical zero $\infty$ and is
Frobenius-equivariant.  The same construction applies to the scaling-site
sheaf.  It does not identify the spherical coefficient $\Fone$ with the
Boolean semiring.
\end{enumerate}

The results above also clarify the coefficient hierarchy.  The spherical
algebra $\Fone[T^{H_+}]$ records multiplicative generators and Frobenius
roots; the characteristic-one scaling sheaf records valuation profiles;
the Boolean semiring records only support.  These levels are related, but
they are not interchangeable.

The paper is organized so that the early sections develop split algebra,
synchronized products, projective semimodules, ideal arithmetic,
projective geometry, and infinitesimal fibres.  Sections
\ref{sec:absolute-curve}--\ref{sec:lcm-perfection} place these structures on
the absolute arithmetic curve.  Section \ref{sec:perfected-anomaly}
computes the carry, rooted-phase, adelic-history, and place-emergence data
at global Frobenius perfection.  Sections \ref{sec:lattice-support} and
\ref{sec:nonassociative-support} develop higher support fibres and
cohomological associators.  Sections \ref{sec:gerbe-quantization} and
\ref{sec:uhf-history} construct the finite gerbe and theta theories and
their operator-algebraic perfected limit.  Sections \ref{sec:beurling-realization}
and \ref{sec:helix-local-system} give the analytic and
Frobenius-dynamical realizations.  Sections \ref{sec:split-local-moduli}
and \ref{sec:split-scaling-site} construct the place cubes and split
characteristic-one site.

\section{Split-zero algebra and universal maps}

Throughout, a semiring is commutative, unital, and has a specified zero; homomorphisms preserve both $0$ and $1$.  The one-element semiring, in which $0=1$, is denoted by $\zero$ and is allowed.

\begin{definition}[Split-zero globalization]
Let $R$ be a nonzero commutative ring.  Define
\[
S:=G(R):=R\sqcup\{\tau\},
\qquad \tau\notin R,
\]
with operations
\[
r\oplus s=r+s,
\qquad
r\oplus\tau=\tau\oplus r=r,
\qquad
\tau\oplus\tau=\tau,
\]
and
\[
rs=rs,
\qquad
r\tau=\tau r=\tau,
\qquad
\tau^2=\tau.
\]
We write
\[
e:=0_R\in R\subset S.
\]
Thus $0_S=\tau$ and $e\neq\tau$.
\end{definition}

\begin{proposition}[Pair normal form]\label{prop:pair-normal-form}
There is a canonical semiring isomorphism
\[
G(R)\xrightarrow{\sim}
D_R:=\{(0,0)\}\cup(R\times\{1\})
\subseteq R\times\B
\]
given by
\[
\tau\longmapsto(0,0),
\qquad
r\longmapsto(r,1).
\]
Under this isomorphism,
\[
e\longmapsto(0,1).
\]
\end{proposition}

\begin{proof}
The set $D_R$ is closed under the componentwise operations of $R\times\B$.  The displayed map is bijective, and the transported operations are exactly the defining operations of $G(R)$.
\end{proof}

\begin{definition}[Reflection and support]
The ring reflection and support character are
\[
p_R:G(R)\longrightarrow R,
\qquad
p_R(\tau)=0,
\quad
p_R(r)=r,
\]
and
\[
\chi_R:G(R)\longrightarrow\B,
\qquad
\chi_R(\tau)=0,
\quad
\chi_R(r)=1.
\]
\end{definition}

The following two universal properties were proved in \cite{Secondary}; they are recorded because they control every later construction.

\begin{theorem}[Universal ring reflection and Boolean character]\label{thm:universal-maps}
Let $R$ be a nonzero commutative ring.
\begin{enumerate}[label=\textup{(\alph*)}]
\item For every commutative ring $A$, composition with $p_R$ induces a bijection
\[
\Hom_{\Ring}(R,A)
\xrightarrow{\sim}
\Hom_{\SemiRing}(G(R),A).
\]
\item The support map $\chi_R$ is the unique semiring homomorphism $G(R)\to\B$.
\end{enumerate}
\end{theorem}

\begin{proof}
Part (a) is Theorem 2.2 of \cite{Secondary}.  For completeness, if $f:G(R)\to A$ has ring target, then
\[
f(e)=f(e\oplus e)=f(e)+f(e),
\]
so subtraction gives $f(e)=0_A$.  Hence $f|_R$ is a ring homomorphism and $f=(f|_R)\circ p_R$.

For part (b), let $u:G(R)\to\B$.  Since $1\oplus(-1)=e$ in $G(R)$,
\[
u(e)=u(1)+u(-1)=1+u(-1)=1.
\]
If $u(r)=0$ for some supported $r\in R$, then
\[
1=u(e)=u(re)=u(r)u(e)=0,
\]
a contradiction.  Hence every supported element maps to $1$, while $\tau$ maps to $0$.
\end{proof}


\section{Synchronized products, independent mixed support, and fixed loci}
\label{sec:synchronized-products}

The product of two split globalizations and the globalization of a product
ring are different objects.  The difference is exactly the availability of
mixed support.

\begin{theorem}[Boolean fibre-product theorem]
\label{thm:boolean-fibre-product}
For nonzero commutative rings $A$ and $B$, the map
\[
\Phi:G(A\times B)\longrightarrow G(A)\times_{\B}G(B)
\]
defined by
\[
\Phi(\tau)=(\tau,\tau),
\qquad
\Phi(a,b)=(a,b)
\]
is a canonical semiring isomorphism.  The fibre product is taken with
respect to the support characters $\chi_A$ and $\chi_B$.
\end{theorem}

\begin{proof}
A pair $(x,y)\in G(A)\times G(B)$ lies in the fibre product precisely when
$\chi_A(x)=\chi_B(y)$.  If this common value is $0$, then
$(x,y)=(\tau,\tau)$.  If it is $1$, then $x=a\in A$ and $y=b\in B$.
Consequently the fibre product is the disjoint union
\[
\{(\tau,\tau)\}\sqcup(A\times B),
\]
and the displayed map is bijective.  Addition and multiplication are
componentwise on both sides and agree with the defining operations of
$G(A\times B)$, so $\Phi$ is a semiring isomorphism.
\end{proof}

\begin{corollary}[Exact source of mixed support]
\label{cor:mixed-support-source}
Inside the independent product $G(A)\times G(B)$ one has the four support
faces
\[
\{(\tau,\tau)\},\qquad A\times\{\tau\},\qquad
\{\tau\}\times B,\qquad A\times B.
\]
The image of $G(A\times B)$ consists exactly of the first and fourth
faces.  Hence the two middle faces are the genuinely mixed-support states.
Equivalently, the support square
\[
G(A)\times G(B)\longrightarrow\B^2
\]
restricts on $G(A\times B)$ to the diagonal copy of $\B$.
\end{corollary}

\begin{proof}
This is the support-mask decomposition of the Cartesian product together
with Theorem~\ref{thm:boolean-fibre-product}.
\end{proof}

The support maps fit into the commutative pullback diagram
\[
\begin{tikzcd}[column sep=large,row sep=large]
G(A\times B) \arrow[r,hook] \arrow[d,"\chi_{A\times B}"']
& G(A)\times G(B) \arrow[d,"\chi_A\times\chi_B"]\\
\B \arrow[r,"\Delta"'] & \B^2,
\end{tikzcd}
\]
where $\Delta(\epsilon)=(\epsilon,\epsilon)$.  The two off-diagonal
Boolean vertices are exactly the mixed-support faces omitted by the upper
left object.

\begin{theorem}[Fixed locus of a supported unit]
\label{thm:fixed-locus-unit}
Let $R$ be a commutative ring and let $u\in R^\times$.  Then
\[
\operatorname{Fix}_{G(R)}(u)
:=\{x\in G(R):ux=x\}
=\{\tau\}\sqcup\Ann_R(u-1).
\]
In particular, if $K$ is a field, then
\[
\operatorname{Fix}_{G(K)}(u)=
\begin{cases}
G(K),&u=1,\\
\{\tau,e\}\cong\B,&u\ne1.
\end{cases}
\]
\end{theorem}

\begin{proof}
The element $\tau$ is fixed because it is multiplicatively absorbing.  A
supported element $r\in R$ is fixed precisely when $ur=r$, equivalently
$(u-1)r=0$.  This proves the first formula.  Over a field, the annihilator
of $u-1$ is all of $K$ when $u=1$ and is $\{0\}$ otherwise.
\end{proof}

\begin{corollary}[The Boolean involution locus]
\label{cor:boolean-involution-locus}
If $R$ is an integral domain of characteristic different from $2$, then
\[
\operatorname{Fix}_{G(R)}(-1)=\{\tau,e\}\cong\B.
\]
\end{corollary}

\begin{proof}
Apply Theorem~\ref{thm:fixed-locus-unit} and use
$\Ann_R(-2)=\{0\}$.
\end{proof}

\section{Idempotent localization and the division trichotomy}

We use localization only through its universal property.  If $A$ is a semiring and $M\subseteq A$ is multiplicatively closed, a localization of $A$ at $M$ is a semiring $M^{-1}A$ with a map $\lambda:A\to M^{-1}A$ such that every $\lambda(m)$ is invertible and every map from $A$ that inverts $M$ factors uniquely through $\lambda$.

\begin{theorem}[Localization at an idempotent]\label{thm:idempotent-localization}
Let $A$ be a commutative semiring and let $a\in A$ satisfy $a^2=a$.  Then
\[
A[a^{-1}]\cong aA,
\]
where $aA=\{ax:x\in A\}$ is regarded as a semiring with zero $a0_A$ and multiplicative identity $a$.
\end{theorem}

\begin{proof}
The subset $aA$ is closed under addition and multiplication because
\[
ax+ay=a(x+y),
\qquad
(ax)(ay)=a^2xy=axy.
\]
The element $a$ is its multiplicative identity.  Define
\[
\pi_a:A\longrightarrow aA,
\qquad
x\longmapsto ax.
\]
Then
\[
\pi_a(a)=a^2=a=1_{aA},
\]
so $a$ becomes invertible.

Let $f:A\to T$ be a semiring homomorphism such that $f(a)$ is invertible.  Since $f(a)$ is idempotent,
\[
f(a)^2=f(a).
\]
Multiplication by $f(a)^{-1}$ gives $f(a)=1_T$.  Define
\[
\overline f:aA\longrightarrow T,
\qquad
\overline f(ax)=f(x).
\]
If $ax=ay$, then
\[
f(a)f(x)=f(a)f(y),
\]
and $f(a)=1_T$, so $f(x)=f(y)$.  Thus $\overline f$ is well defined.  It is a semiring homomorphism and satisfies $f=\overline f\circ\pi_a$.  Uniqueness follows because every element of $aA$ has the form $\pi_a(x)$.
\end{proof}

\begin{corollary}[The two zero localizations]\label{cor:two-zero-localizations}
For every nonzero commutative ring $R$,
\[
G(R)[e^{-1}]\cong\{\tau,e\}\cong\B,
\]
and
\[
G(R)[\tau^{-1}]\cong\{\tau\}\cong\zero.
\]
Under the first isomorphism, the localization map is the support character:
\[
\frac{x}{e}=\chi_R(x).
\]
\end{corollary}

\begin{proof}
Both $e$ and $\tau$ are idempotent.  Moreover,
\[
eG(R)=\{\tau,e\},
\qquad
\tau G(R)=\{\tau\}.
\]
In $eG(R)$ the element $e$ is the identity and $\tau$ is the zero, so $eG(R)\cong\B$.  The localization map is $x\mapsto ex$, which sends $\tau$ to $\tau$ and every supported element to $e$; this is exactly $\chi_R$ after the identification $\tau\leftrightarrow0$, $e\leftrightarrow1$.
\end{proof}

\begin{theorem}[Compatibility with ordinary ring localization]\label{thm:ring-localization}
Let $M\subseteq R$ be a multiplicatively closed set with $0\notin M$.  Then the canonical map $R\to M^{-1}R$ induces a canonical semiring isomorphism
\[
M^{-1}G(R)\xrightarrow{\sim}G(M^{-1}R).
\]
\end{theorem}

\begin{proof}
Let
\[
\lambda:G(R)\longrightarrow G(M^{-1}R)
\]
be defined by $\lambda(\tau)=\tau$ and $\lambda(r)=r/1$.  Every $m\in M$ maps to the supported unit $m/1$.

Let $f:G(R)\to T$ be a semiring homomorphism such that every $f(m)$ is invertible.  Define
\[
\overline f(\tau):=0_T,
\qquad
\overline f(r/m):=f(r)f(m)^{-1}.
\]
If $r/m=s/n$ in $M^{-1}R$, then there exists $u\in M$ with
\[
u(nr-ms)=0
\]
in $R$, equivalently $unr=ums$.  Applying $f$ gives
\[
f(u)f(n)f(r)=f(u)f(m)f(s).
\]
Since $f(u)$, $f(m)$, and $f(n)$ are invertible, one obtains
\[
f(r)f(m)^{-1}=f(s)f(n)^{-1}.
\]
Thus $\overline f$ is well defined.  The usual fraction calculations show that it is a semiring homomorphism.  It is the unique map satisfying $\overline f\circ\lambda=f$, because the supported fractions and $\tau$ exhaust $G(M^{-1}R)$.
\end{proof}

\begin{corollary}[Classical fraction field versus Boolean fraction residue]\label{cor:fraction-comparison}
Let $R$ be an integral domain with fraction field $K$.
\begin{enumerate}[label=\textup{(\alph*)}]
\item Localizing at the classically nonzero elements gives
\[
G(R)[(R\setminus\{0\})^{-1}]
\cong
G(K).
\]
\item Localizing at all semiring-nonzero elements, namely at $R=G(R)\setminus\{\tau\}$, gives
\[
G(R)[R^{-1}]
\cong
\B.
\]
\end{enumerate}
\end{corollary}

\begin{proof}
Part (a) is Theorem \ref{thm:ring-localization}.  For part (b), the denominator set contains $e$, so the localization factors through $G(R)[e^{-1}]\cong\B$.  In $\B$, every supported element of $G(R)$ already maps to $1$, hence is invertible.  Therefore $\B$ has the universal property of the localization at all supported elements.
\end{proof}

\begin{remark}[Element localization versus localization at a prime]\label{rem:prime-v-element}
For $R=\Z$, let
\[
P_0=\{\tau,0\}.
\]
The stalk localization at $P_0$ inverts the complement $\Z\setminus\{0\}$ and is therefore $G(\Q)$.  Principal localization at the element $e=0$ instead gives $\B$.  Thus ``localizing at the generic zero prime'' and ``placing the zero element in the denominator set'' are distinct operations.
\end{remark}

\begin{theorem}[Strict identification with a nonzero amplitude]\label{thm:strict-identification-collapse}
Let $K$ be a field, let $S=G(K)$, and let $h\in K^\times$.  Let $\sim$ be the smallest semiring congruence on $S$ satisfying
\[
e\sim h.
\]
Then
\[
S/{\sim}\cong\B.
\]
\end{theorem}

\begin{proof}
In the quotient, the class of $h$ is invertible, so the class of $e$ is invertible.  Since $e^2=e$, its class is an invertible idempotent and therefore equals $1$.  For every supported $x\in K$,
\[
ex=e
\]
in $S$, hence in the quotient
\[
x=1x=ex=e=1.
\]
Thus all supported elements lie in one congruence class.  The support character $\chi_K:S\to\B$ satisfies $\chi_K(e)=\chi_K(h)=1$, so it coequalizes the imposed relation and distinguishes $\tau$ from the supported class.  Hence the quotient has exactly two elements and is Boolean.
\end{proof}

\begin{theorem}[Division trichotomy]\label{thm:division-trichotomy}
Let $R$ be an integral domain with fraction field $K$.  The three universal division operations are
\[
\boxed{
\begin{aligned}
\text{invert }R\setminus\{0\}
&\quad\Longrightarrow\quad G(K),\\
\text{invert }e=0_R
&\quad\Longrightarrow\quad \B,\\
\text{invert }\tau=0_{G(R)}
&\quad\Longrightarrow\quad \zero.
\end{aligned}}
\]
The first preserves amplitude and support, the second preserves only support, and the third preserves no distinction.
\end{theorem}

\begin{proof}
This is the conjunction of Corollaries \ref{cor:two-zero-localizations} and \ref{cor:fraction-comparison}.
\end{proof}

\section{Ring-null infinitesimals and Boolean cubes}

The preceding localization theorem suggests an intrinsic infinitesimal notion which is independent of order, topology, or valuation.

\begin{definition}[Absolute ring-null halo]\label{def:absolute-halo}
Let $A$ be a commutative semiring admitting a universal ring reflection
\[
\rho_A:A\longrightarrow A^{\mathrm{ring}}.
\]
Define
\[
\mathfrak I_{\mathrm{abs}}(A):=\rho_A^{-1}(0).
\]
An element of $\mathfrak I_{\mathrm{abs}}(A)\setminus\{0_A\}$ is called an \emph{absolute ring-null infinitesimal}.
\end{definition}

\begin{proposition}[Intrinsic characterization]\label{prop:absolute-characterization}
For $x\in A$,
\[
x\in\mathfrak I_{\mathrm{abs}}(A)
\]
if and only if every semiring homomorphism from $A$ to a commutative ring sends $x$ to zero.
\end{proposition}

\begin{proof}
Every map from $A$ to a ring factors uniquely through $\rho_A$.  Hence $\rho_A(x)=0$ implies that every ring-valued map kills $x$.  Conversely, apply the assumed condition to the ring-valued map $\rho_A$ itself.
\end{proof}

\begin{corollary}[The scalar split-zero infinitesimal]\label{cor:scalar-absolute-halo}
For $S=G(R)$,
\[
\mathfrak I_{\mathrm{abs}}(S)=\{\tau,e\}.
\]
Thus $e$ is the unique nonzero absolute ring-null infinitesimal of $S$.
\end{corollary}

\begin{proof}
The universal ring reflection is $p_R$ by Theorem \ref{thm:universal-maps}, and
\[
p_R^{-1}(0_R)=\{\tau,e\}.
\]
The semiring zero is $\tau$, so $e$ is the unique nonzero member of this fibre.
\end{proof}

Let $I$ be a finite set and write
\[
S^I=\prod_{i\in I}S.
\]
For $A\subseteq I$, define the supported-zero mask
\[
(\eps_A)_i=
\begin{cases}
e,&i\in A,\\
\tau,&i\notin A.
\end{cases}
\]

\begin{theorem}[Finite-product ring reflection]\label{thm:product-ring-reflection}
The coordinatewise map
\[
p_I:=p_R^I:S^I\longrightarrow R^I
\]
is the universal ring reflection of $S^I$.  Its kernel is
\[
p_I^{-1}(0)=\{\eps_A:A\subseteq I\}.
\]
Moreover,
\[
\eps_A+\eps_B=\eps_{A\cup B},
\qquad
\eps_A\eps_B=\eps_{A\cap B}.
\]
Consequently,
\[
\mathfrak I_{\mathrm{abs}}(S^I)
\cong\B^I
\]
as semirings.
\end{theorem}

\begin{proof}
Let $f:S^I\to T$ be a semiring homomorphism with ring target.  Each $\eps_A$ is additively idempotent:
\[
\eps_A+\eps_A=\eps_A.
\]
Therefore
\[
f(\eps_A)+f(\eps_A)=f(\eps_A),
\]
and subtraction in $T$ gives $f(\eps_A)=0$.

For $r=(r_i)_{i\in I}\in R^I$, let $\widetilde r\in S^I$ be its fully supported lift, so $(\widetilde r)_i=r_i\in R\subset S$.  Define
\[
g:R^I\longrightarrow T,
\qquad
g(r):=f(\widetilde r).
\]
Because sums and products of fully supported lifts remain fully supported,
\[
\widetilde{r+s}=\widetilde r+\widetilde s,
\qquad
\widetilde{rs}=\widetilde r\,\widetilde s,
\]
where the operations on $R^I$ are coordinatewise.  Hence $g$ is a ring homomorphism.

If $x\in S^I$ and $r=p_I(x)$, then $x$ and $\widetilde r$ differ only in coordinates where $r_i=0$: in such a coordinate $x_i$ is either $\tau$ or $e$, while $(\widetilde r)_i=e$.  Therefore there is a mask $\eps_A$ such that
\[
\widetilde r=x+\eps_A.
\]
Applying $f$ gives
\[
f(\widetilde r)=f(x)+f(\eps_A)=f(x).
\]
Thus $f=g\circ p_I$.  Uniqueness follows from surjectivity of $p_I$.

The kernel description is immediate coordinatewise.  The union and intersection formulas follow from
\[
e+e=e,
\quad e+\tau=e,
\quad \tau+\tau=\tau,
\]
and
\[
e^2=e,
\quad e\tau=\tau,
\quad \tau^2=\tau.
\]
These are exactly the join and meet operations on $\B^I$.
\end{proof}

\begin{example}[The mixed zero $(\tau,e)$]
In $S^2$ the four absolute infinitesimals are
\[
(\tau,\tau),
\quad
(e,\tau),
\quad
(\tau,e),
\quad
(e,e).
\]
The element $(\tau,e)$ means that the first coordinate is absent while the second coordinate is supported but arithmetically zero.  For $a\in R$,
\[
(a,\tau)+(-a,\tau)=(e,\tau),
\]
so cancellation produces supported zero rather than absence.
\end{example}

\subsection{Derivations and the failure of tangent displacement}

\begin{definition}[Semiring derivation]
Let $A$ be a semiring and let $M$ be an $A$-semimodule.  A derivation is a map $d:A\to M$ satisfying
\[
d(0)=0,
\qquad
d(1)=0,
\qquad
d(x+y)=d(x)+d(y),
\]
and
\[
d(xy)=x\,d(y)+y\,d(x).
\]
\end{definition}

\begin{proposition}[The supported zero is derivationally constant]\label{prop:derivation-kills-e}
Every derivation $d:G(R)\to M$ satisfies
\[
d(e)=0.
\]
Every derivation $\B\to M$ is zero.
\end{proposition}

\begin{proof}
Since $1+e=1$ in $G(R)$,
\[
0=d(1)=d(1+e)=d(1)+d(e)=d(e).
\]
For $\B$, the only elements are $0$ and $1$, both of which are killed by the defining axioms of a derivation.
\end{proof}

\begin{remark}
The identity $x+e=x$ for every supported $x\in R$ shows that $e$ does not move a supported point.  It is therefore an infinitesimal only in the ring-null support sense of Definition \ref{def:absolute-halo}; it is not a first-order tangent parameter.
\end{remark}

The closest ring-valued multiplicative model of the relation $e^2=e$ makes this distinction exact.

\begin{theorem}[Idempotent probes are finite-difference probes]\label{thm:idempotent-probe}
Let $A$ be a commutative ring and set
\[
C:=A[q]/(q^2-q).
\]
Then evaluation at $q=0$ and $q=1$ induces a canonical ring isomorphism
\[
C\xrightarrow{\sim}A\times A.
\]
For every polynomial $f\in A[T]$ and every $x,h\in A$,
\[
f(x+qh)
=
f(x)+q\bigl(f(x+h)-f(x)\bigr)
\]
in $C$.  Moreover,
\[
\Omega^1_{C/A}=0.
\]
\end{theorem}

\begin{proof}
The ideals $(q)$ and $(q-1)$ of $A[q]$ are comaximal because
\[
q-(q-1)=1.
\]
Their intersection equals their product $(q(q-1))=(q^2-q)$.  The Chinese remainder theorem gives
\[
A[q]/(q^2-q)
\cong
A[q]/(q)\times A[q]/(q-1)
\cong
A\times A.
\]
Under this isomorphism, $q\mapsto(0,1)$ and $x+qh\mapsto(x,x+h)$.  Both sides of the polynomial identity map to
\[
\bigl(f(x),f(x+h)\bigr),
\]
so they are equal.

For relative K\"ahler differentials, the relation $q^2-q=0$ gives
\[
(2q-1)dq=0.
\]
But
\[
(2q-1)^2=4q^2-4q+1=1
\]
in $C$, so $2q-1$ is invertible and $dq=0$.  Since $C$ is generated over $A$ by $q$, this implies $\Omega^1_{C/A}=0$.
\end{proof}

\begin{corollary}[Idempotence versus nilpotence]
An idempotent parameter records two endpoint values and exact finite differences.  A nilpotent parameter $\epsilon^2=0$ records first-order tangent data.  These structures cannot be identified in a nonzero ring.
\end{corollary}

\begin{proof}
The finite-difference statement is Theorem \ref{thm:idempotent-probe}.  If a ring element satisfies both $q^2=q$ and $q^2=0$, then $q=0$.
\end{proof}

\section{Mixed-support localization and scoped identities}

\begin{lemma}[Localization of a finite product]\label{lem:product-localization}
Let $A_i$ be commutative semirings indexed by a finite set $I$, let
\[
A:=\prod_{i\in I}A_i,
\qquad
x=(x_i)_{i\in I}\in A.
\]
Then
\[
A[x^{-1}]
\cong
\prod_{i\in I}A_i[x_i^{-1}].
\]
\end{lemma}

\begin{proof}
Use the standard fraction construction for localization.  Every element of $A[x^{-1}]$ is represented by a fraction
\[
\frac{y}{x^n},
\qquad y\in A,
\quad n\ge0.
\]
Define
\[
\Phi\left(\frac{y}{x^n}\right)
:=
\left(\frac{y_i}{x_i^n}\right)_{i\in I}.
\]
If two fractions are equal in $A[x^{-1}]$, their coordinate fractions are equal, so $\Phi$ is well defined.

For surjectivity, choose fractions $y_i/x_i^{n_i}$ in the factors and put $N=\max_i n_i$.  Replacing $y_i$ by $y_i x_i^{N-n_i}$ gives a common exponent $N$, hence a preimage in $A[x^{-1}]$.

For injectivity, suppose the coordinate fractions of $y/x^n$ and $z/x^m$ are equal.  For each $i$ there exists $k_i\ge0$ such that
\[
x_i^{k_i+m}y_i=x_i^{k_i+n}z_i.
\]
Let $k=\max_i k_i$.  Multiplying each equality by a further power of $x_i$ yields
\[
x_i^{k+m}y_i=x_i^{k+n}z_i
\]
for every $i$.  Therefore
\[
x^{k+m}y=x^{k+n}z
\]
in $A$, which is precisely the equality of the original fractions in $A[x^{-1}]$.
\end{proof}

Let $K$ be a field, $S=G(K)$, and $I$ finite.  For $x\in S^I$, define
\[
U(x):=\{i\in I:x_i\in K^\times\},
\]
\[
Z(x):=\{i\in I:x_i=e\},
\]
and
\[
N(x):=\{i\in I:x_i=\tau\}.
\]
These form a partition of $I$.

\begin{theorem}[Mixed-support localization normal form]\label{thm:mixed-localization}
For every $x\in S^I$,
\[
S^I[x^{-1}]
\cong
S^{U(x)}\times\B^{Z(x)}.
\]
The coordinates in $N(x)$ contribute one-element factors and hence disappear.  Under the displayed isomorphism,
\[
\frac{y}{x}
\longmapsto
\left(
\left(\frac{y_i}{x_i}\right)_{i\in U(x)},
\left(\chi_K(y_i)\right)_{i\in Z(x)}
\right).
\]
\end{theorem}

\begin{proof}
By Lemma \ref{lem:product-localization},
\[
S^I[x^{-1}]
\cong
\prod_{i\in I}S[x_i^{-1}].
\]
If $i\in U(x)$, then $x_i$ is already a unit of $S$, so $S[x_i^{-1}]\cong S$.  If $i\in Z(x)$, Corollary \ref{cor:two-zero-localizations} gives $S[e^{-1}]\cong\B$.  If $i\in N(x)$, the same corollary gives $S[\tau^{-1}]\cong\zero$.  Since a finite product with a one-element factor is canonically isomorphic to the product of the remaining factors, the stated formula follows.  The explicit fraction formula is the coordinatewise description of these three identifications.
\end{proof}

\begin{example}
For $x=(\tau,e)\in S^2$,
\[
S^2[x^{-1}]\cong\B,
\]
and
\[
\frac{(a,b)}{(\tau,e)}=\chi_K(b).
\]
The first coordinate is erased; the second records only support.
\end{example}

\begin{example}
For $x=(h,e)\in S^2$ with $h\in K^\times$,
\[
S^2[x^{-1}]\cong S\times\B,
\]
and
\[
\frac{(a,b)}{(h,e)}=\bigl(ah^{-1},\chi_K(b)\bigr).
\]
This is ordinary division in one coordinate and Boolean division by supported zero in the other.
\end{example}

We next formalize the effect of adjoining an equality only on selected support coordinates.

For $i\in I$ and $s\in S$, write $r_i(s)\in S^I$ for the vector with entry $s$ at $i$ and $\tau$ elsewhere.  Fix a subset $A\subseteq I$ and units $h_i\in K^\times$ for $i\in A$.  Define
\[
\eps_A:=\sum_{i\in A}r_i(e),
\qquad
h_A:=\sum_{i\in A}r_i(h_i).
\]

\begin{theorem}[Scoped identity quotient]\label{thm:scoped-identity}
Let $\sim_A$ be the smallest semiring congruence on $S^I$ satisfying
\[
\eps_A\sim_A h_A.
\]
Then
\[
S^I/{\sim_A}
\cong
\B^A\times S^{I\setminus A}.
\]
The quotient map is
\[
q_A(x)
=
\left(
(\chi_K(x_i))_{i\in A},
(x_j)_{j\notin A}
\right).
\]
\end{theorem}

\begin{proof}
The map $q_A$ is a semiring homomorphism and satisfies
\[
q_A(\eps_A)=q_A(h_A),
\]
so it factors through $S^I/{\sim_A}$.

Let $f:S^I\to T$ be any semiring homomorphism satisfying $f(\eps_A)=f(h_A)$.  For $i\in A$, multiply the relation by $r_i(1)$ to obtain
\[
f(r_i(e))=f(r_i(h_i)).
\]
Multiplying again by $r_i(h_i^{-1})$ gives
\[
f(r_i(e))=f(r_i(1)).
\]
For every supported $s\in K$,
\[
r_i(s)=r_i(s)r_i(1),
\]
and therefore
\[
f(r_i(s))
=f(r_i(s))f(r_i(1))
=f(r_i(s))f(r_i(e))
=f(r_i(se))
=f(r_i(e)).
\]
Thus on coordinate $i\in A$, the value of $f$ depends only on whether the coordinate is $\tau$ or supported.  On coordinates outside $A$, no relation has been imposed.

Every vector $x\in S^I$ decomposes as
\[
x=\sum_{i\in I}r_i(x_i).
\]
Hence $f(x)$ depends only on $q_A(x)$, which gives a unique factorization of $f$ through $q_A$.  This is the universal property of the quotient.
\end{proof}

\begin{corollary}[Modal equality without arithmetic equality]\label{cor:modal-equality}
For every $h\in K^\times$,
\[
\chi_K(e)=\chi_K(h)=1.
\]
Thus $e$ and $h$ are equal after support reflection.  Their strict arithmetic identification gives the Boolean quotient by Theorem \ref{thm:strict-identification-collapse}; their identification on selected coordinates gives the partial Booleanization of Theorem \ref{thm:scoped-identity}.
\end{corollary}


\section{Distributive-lattice support and higher zero fibres}
\label{sec:lattice-support}

The scalar split-zero semiring and the Boolean cube are the first two
instances of a construction controlled by a bounded distributive lattice.
This section gives the construction in a form which simultaneously covers
independent mixed support, iterated sub-zero layers, and nilpotent ideal
filtrations.

Let $L$ be a bounded distributive lattice with bottom $0_L$, top $1_L$,
join $\vee$, and meet $\wedge$.  Regard $L$ as an idempotent semiring by
\[
\lambda\oplus\mu:=\lambda\vee\mu,
\qquad
\lambda\odot\mu:=\lambda\wedge\mu.
\]
Its additive zero is $0_L$ and its multiplicative unit is $1_L$.

\begin{definition}[Lattice-split globalization]
\label{def:lattice-split}
For a commutative ring $R$, define
\[
G_L(R):=
\{(0,\lambda):\lambda\in L\}
\cup
\{(r,1_L):r\in R\}
\subseteq R\times L,
\]
with the componentwise operations
\begin{align*}
(r,\lambda)+(s,\mu)&=(r+s,\lambda\vee\mu),\\
(r,\lambda)(s,\mu)&=(rs,\lambda\wedge\mu).
\end{align*}
Write
\[
\tau_L:=(0,0_L),
\qquad
e_L:=(0,1_L).
\]
\end{definition}

The point $\tau_L$ is actual absence, whereas $e_L$ is the supported ring
zero.  A nonzero amplitude can occur only over maximal support $1_L$.

\begin{theorem}[Lattice normal form and universal maps]
\label{thm:lattice-normal-form}
The following statements hold.
\begin{enumerate}[label=\textup{(\arabic*)}]
\item $G_L(R)$ is a commutative unital semiring with zero $\tau_L$ and
unit $(1_R,1_L)$.
\item Projection to amplitude
\[
p_L:G_L(R)\longrightarrow R,
\qquad
p_L(r,\lambda)=r,
\]
is the universal morphism from $G_L(R)$ to a ring.
\item Projection to support
\[
\chi_L:G_L(R)\longrightarrow L,
\qquad
\chi_L(r,\lambda)=\lambda,
\]
is a semiring homomorphism.
\item The ring-null fibre is canonically the lattice semiring:
\[
p_L^{-1}(0_R)
=\{(0,\lambda):\lambda\in L\}
\cong (L,\vee,\wedge).
\]
\end{enumerate}
\end{theorem}

\begin{proof}
The subset in Definition~\ref{def:lattice-split} is closed under both
componentwise operations.  If one summand has nonzero amplitude, its
support is $1_L$, so the sum again has support $1_L$.  If one factor has
support below $1_L$, its amplitude is zero, so the product has zero
amplitude and support below or equal to that factor.  Associativity and
commutativity follow componentwise.  Distributivity follows from
\[
\lambda\wedge(\mu\vee\nu)
=(\lambda\wedge\mu)\vee(\lambda\wedge\nu)
\]
and the distributive law in $R$.  The stated zero and unit are immediate.

The maps $p_L$ and $\chi_L$ preserve the operations.  Let
$f:G_L(R)\to A$ be a semiring homomorphism into a ring $A$.  For every
$\lambda\in L$, put $z_\lambda=(0,\lambda)$.  Since
\[
e_L+z_\lambda=e_L,
\]
additive cancellation in $A$ gives $f(z_\lambda)=0$.  Hence $f$ is
determined by its restriction to the supported copy of $R$, and that
restriction is a ring homomorphism.  Conversely, every ring homomorphism
$R\to A$ extends uniquely by sending all $z_\lambda$ to zero.  This is the
universal property of $p_L$.  The description of the zero fibre is
immediate, and its inherited operations are join and meet.
\end{proof}

\begin{theorem}[Localization at a support stratum]
\label{thm:lattice-localization}
For $\lambda\in L$, let $z_\lambda=(0,\lambda)$.  Then
$z_\lambda^2=z_\lambda$ and
\[
G_L(R)[z_\lambda^{-1}]
\cong
\downarrow\lambda
:=\{\mu\in L:\mu\leq\lambda\},
\]
where $\downarrow\lambda$ is equipped with join and meet, zero $0_L$, and
unit $\lambda$.  In particular,
\[
G_L(R)[e_L^{-1}]\cong L,
\qquad
G_L(R)[\tau_L^{-1}]\cong\zero.
\]
\end{theorem}

\begin{proof}
The element $z_\lambda$ is multiplicatively idempotent.  By
Theorem~\ref{thm:idempotent-localization}, localization at it is the
principal semiring $z_\lambda G_L(R)$.  Multiplication gives
\[
z_\lambda(r,1_L)=z_\lambda,
\qquad
z_\lambda z_\mu=z_{\lambda\wedge\mu}.
\]
Therefore
\[
z_\lambda G_L(R)=\{z_\nu:\nu\leq\lambda\},
\]
and the induced operations are join and meet on the principal interval.
The endpoint cases are $\lambda=1_L$ and $\lambda=0_L$.
\end{proof}

\begin{example}[Boolean cubes]
Let $L=\mathcal P(I)$, ordered by inclusion.  Then the ring-null fibre of
$G_L(R)$ is the Boolean cube of support masks, addition is union, and
multiplication is intersection.  Localization at a mask $A\subseteq I$
retains exactly the subcube $\mathcal P(A)$.
\end{example}

\begin{example}[The iterated sub-zero chain]
Let
\[
C_n=\{0<1<\cdots<n\}.
\]
In $G_{C_n}(R)$ define
\[
\tau_i:=(0,n-i),
\qquad 1\leq i\leq n.
\]
Then
\[
\tau_i+\tau_j=\tau_{\min(i,j)},
\qquad
\tau_i\tau_j=\tau_{\max(i,j)},
\]
while $r+\tau_i=r$ and $r\tau_i=\tau_i$ for $r\in R$.  Thus
$G_{C_n}(R)$ is exactly the $n$-fold sub-zero globalization
\[
R\sqcup\{\tau_1,\dots,\tau_n\}.
\]
The case $n=1$ is the ordinary split-zero construction.
\end{example}

\begin{theorem}[Nilpotent-thickening realization of the support chain]
\label{thm:nilpotent-chain-realization}
Let $n\geq1$.
\begin{enumerate}[label=\textup{(\arabic*)}]
\item For every prime $p$, let
\[
B_{p,n}=\Z/p^n\Z,
\qquad
I_a=p^aB_{p,n}
\quad(0\leq a\leq n).
\]
Then
\[
I_a+I_b=I_{\min(a,b)},
\qquad
I_a\cap I_b=I_{\max(a,b)}.
\]
The map
\[
C_n\longrightarrow\{I_0,\dots,I_n\},
\qquad
j\longmapsto I_{n-j},
\]
is a semiring isomorphism from $(C_n,\vee,\wedge)$ to the ideal chain
with addition given by ideal sum and multiplication by intersection.
\item For every field $k$, the same statement holds for
\[
B_{\epsilon,n}=k[\epsilon]/(\epsilon^n),
\qquad
J_a=\epsilon^aB_{\epsilon,n}.
\]
\end{enumerate}
Consequently, the zero fibre of the $n$-fold sub-zero globalization is
canonically isomorphic, after choosing its length, to the support semiring
of either a $p^n$-thickening or an infinitesimal $\epsilon^n$-thickening.
\end{theorem}

\begin{proof}
Both rings are principal ideal rings with the displayed chains.  If
$a\leq b$, then $I_b\subseteq I_a$ and $J_b\subseteq J_a$.  Hence the
sum of two ideals is the larger one and their intersection is the smaller
one, which gives the formulas.  Under $j\mapsto I_{n-j}$ or
$j\mapsto J_{n-j}$, join in $C_n$ becomes ideal sum and meet becomes
intersection.
\end{proof}

\begin{proposition}[The valuation-sensitive shadow]
\label{prop:truncated-valuation-semiring}
On the same ideal sets, use ordinary ideal multiplication instead of
intersection.  Then
\[
I_aI_b=I_{\min(n,a+b)},
\qquad
J_aJ_b=J_{\min(n,a+b)}.
\]
Thus both ideal semirings are isomorphic to
\[
V_n=\{0,1,\dots,n\},
\qquad
a\oplus b=\min(a,b),
\qquad
a\odot b=\min(n,a+b),
\]
with additive zero $n$ and multiplicative unit $0$.  The lattice semiring
of Theorem~\ref{thm:nilpotent-chain-realization} records support strata;
$V_n$ records order of vanishing.  They agree on addition and differ on
multiplication.
\end{proposition}

\begin{proof}
The product of principal ideals generated by $p^a,p^b$ is generated by
$p^{a+b}$, with $p^n=0$ in $B_{p,n}$.  The identical calculation holds
for powers of $\epsilon$.  The displayed formulas give the claimed
semiring isomorphisms.
\end{proof}

\begin{remark}
The preceding theorems identify support and valuation semirings, not the
ambient rings.  A characteristic-zero infinitesimal thickening and a
positive-characteristic $p$-power thickening need not be isomorphic.  What
is common is the finite filtration geometry seen by their ideal chains.
\end{remark}

\section{Nonassociative coefficients and cocycle-controlled associators}
\label{sec:nonassociative-support}

The support construction does not require associative coefficients.  This
provides a precise interface with Cayley--Dickson algebras and with
nonassociative algebras internal to tensor categories.

\begin{definition}[Lattice-supported algebra]
Let $k$ be a commutative ring and let $A$ be a unital, not necessarily
associative, $k$-algebra.  Define
\[
G_L(A)=
\{(0,\lambda):\lambda\in L\}
\cup
\{(a,1_L):a\in A\}
\]
with componentwise addition and
\[
(a,\lambda)(b,\mu)=(ab,\lambda\wedge\mu).
\]
\end{definition}

\begin{theorem}[Separation of support and associator]
\label{thm:support-associator-separation}
For $x=(a,\lambda)$, $y=(b,\mu)$, and $z=(c,\nu)$ in $G_L(A)$,
\begin{align*}
(xy)z&=((ab)c,\lambda\wedge\mu\wedge\nu),\\
x(yz)&=(a(bc),\lambda\wedge\mu\wedge\nu).
\end{align*}
Therefore:
\begin{enumerate}[label=\textup{(\arabic*)}]
\item $G_L(A)$ is associative if and only if $A$ is associative;
\item alternativity, flexibility, and power associativity pass from $A$
to $G_L(A)$ and are reflected by the fully supported copy of $A$;
\item if $a,b\neq0$ and $ab=0$ in $A$, then
\[
(a,1_L)(b,1_L)=e_L=(0,1_L)\neq\tau_L;
\]
thus intrinsic zero-divisor cancellation is supported, whereas multiplication
by an absent factor produces $\tau_L$.
\end{enumerate}
\end{theorem}

\begin{proof}
The two displayed formulas are immediate from componentwise multiplication
and associativity of the meet operation.  Every polynomial identity in the
multiplication having the same multiset of variables on both sides has the
same support on both sides, so any such identity valid in $A$ remains valid
in $G_L(A)$.  The fully supported embedding $a\mapsto(a,1_L)$ reflects
failures of these identities.  The final assertion follows directly from
the multiplication law and from $e_L\neq\tau_L$.
\end{proof}

\begin{definition}[Twisted graded coefficient algebra]
Let $\Gamma$ be a group, let $k$ be a field, and let
\[
F:\Gamma\times\Gamma\longrightarrow k^\times
\]
be a normalized $2$-cochain.  On the vector space with basis
$\{u_g:g\in\Gamma\}$ define
\[
u_gu_h=F(g,h)u_{gh}.
\]
Set
\[
\phi_F(g,h,k)
=
\frac{F(g,h)F(gh,k)}{F(h,k)F(g,hk)}.
\]
\end{definition}

\begin{theorem}[Associator cocycle and internal associativity]
\label{thm:associator-cocycle}
The function $\phi_F$ is a normalized $3$-cocycle; with the displayed
convention it is the coboundary of $F$.  On homogeneous basis vectors,
\[
(u_gu_h)u_k
=\phi_F(g,h,k)\,u_g(u_hu_k).
\]
Consequently, the twisted algebra is associative precisely when
$\phi_F=1$.  For arbitrary $F$ it is an associative algebra object in the
monoidal category of $\Gamma$-graded vector spaces whose associator on
degrees $(g,h,k)$ is multiplication by $\phi_F(g,h,k)$.

Moreover, the lattice-supported algebra $G_L(k_F[\Gamma])$ has the same
basis associator cocycle $\phi_F$; the support meet contributes no
additional associator.
\end{theorem}

\begin{proof}
The displayed associator identity follows by comparing
\[
F(g,h)F(gh,k)u_{ghk}
\quad\text{and}\quad
F(h,k)F(g,hk)u_{ghk}.
\]
The pentagon identity for the monoidal associator is exactly the
$3$-cocycle equation for $\phi_F$; equivalently, $\phi_F$ is the
cochain coboundary of $F$ in the stated convention.  The algebra-object
associativity diagram is the displayed identity.  In the
lattice-supported extension, both parenthesizations have support
$\lambda\wedge\mu\wedge\nu$, so the scalar associator remains
$\phi_F$.
\end{proof}

\begin{corollary}[A three-axis coefficient package]
\label{cor:support-phase-associator-package}
The object
\[
\mathscr A(L,\Gamma,F):=G_L(k_F[\Gamma])
\]
separates three kinds of information:
\[
\begin{array}{c|c}
L&\text{support and definedness},\\
\Gamma&\text{phase or grading sector},\\
\phi_F\in Z^3(\Gamma;k^\times)&\text{pointwise associator data}.
\end{array}
\]
Ring reflection forgets only the support lattice, localization restricts
support to a principal interval, and neither operation changes the
supported associator cochain.
\end{corollary}

\begin{remark}[Cohomology class versus chosen multiplication]
Because a global multiplication cochain $F$ has been chosen,
$\phi_F$ is exact in ordinary group cohomology.  Its pointwise values can
nevertheless be nontrivial and encode the failure of the chosen basis
multiplication to associate.  A genuinely nonzero class in
$H^3(\Gamma;k^\times)$ arises when the monoidal category carries a
background associator which cannot be globally strictified by such an
$F$.
\end{remark}

\begin{theorem}[Cayley--Dickson support threshold]
\label{thm:cayley-dickson-support-threshold}
Let $\CD_r$ denote the standard real Cayley--Dickson algebra of dimension
$2^r$, so that
\[
\CD_0=\R,
\quad
\CD_1=\C,
\quad
\CD_2=\mathbb H,
\quad
\CD_3=\mathbb O,
\quad
\CD_4=\mathbb S
\]
with $\mathbb S$ the sedenions.  Then:
\begin{enumerate}[label=\textup{(\arabic*)}]
\item for $r\leq3$, no two nonzero fully supported elements of
$G_L(\CD_r)$ multiply to $e_L$;
\item for every $r\geq4$, there exist nonzero fully supported elements
$x,y\in G_L(\CD_r)$ with $xy=e_L$;
\item in a signed graded Cayley--Dickson basis indexed by
$(\Z/2\Z)^r$, the support completion preserves the basis associator
$3$-cocycle.
\end{enumerate}
Thus the octonionic loss of associativity and the sedenionic appearance of
zero divisors are distinct events, and split support records the second
without conflating it with absence.
\end{theorem}

\begin{proof}
The algebras $\R,\C,\mathbb H,\mathbb O$ are division algebras, so the
first assertion follows from
Theorem~\ref{thm:support-associator-separation}.  The sedenions possess
nontrivial zero divisors, and every later Cayley--Dickson algebra contains
the preceding one through the standard embedding $a\mapsto(a,0)$.  Hence
sedenion zero divisors give the second assertion at every later stage.
A signed graded basis has multiplication controlled by a
$\{\pm1\}$-valued $2$-cochain, so the third assertion is
Theorem~\ref{thm:associator-cocycle}.
\end{proof}

\begin{remark}
Theorem~\ref{thm:cayley-dickson-support-threshold} does not identify every
supported zero with a Cayley--Dickson zero divisor.  It gives an exact
criterion for when an intrinsic coefficient zero-divisor pair becomes a
supported-zero product in the split completion.  Combinatorial support
failure and coefficient-algebra zero-divisibility remain independent
invariants.
\end{remark}

\section{The split spectrum and the \texorpdfstring{$\mathbb F_1$}{F1} support skeleton}

\begin{theorem}[Ideals and prime ideals]\label{thm:prime-classification}
Let $R$ be a nonzero commutative ring and $S=G(R)$.
\begin{enumerate}[label=\textup{(\alph*)}]
\item Every ideal of $S$ is either $\{\tau\}$ or
\[
I_J:=J\cup\{\tau\}
\]
for a unique ideal $J\triangleleft R$.
\item The prime ideals are
\[
P_\tau:=\{\tau\}
\]
and
\[
P_{\mathfrak p}:=\mathfrak p\cup\{\tau\},
\qquad
\mathfrak p\in\Spec R.
\]
\end{enumerate}
\end{theorem}

\begin{proof}
Every ideal contains $\tau$.  If $I\cap R=\varnothing$, then $I=\{\tau\}$.  Otherwise $J:=I\cap R$ is closed under addition, multiplication by arbitrary elements of $R$, and additive inverses because $-a=(-1)a$.  Hence $J\triangleleft R$ and $I=J\cup\{\tau\}$.  Conversely every such set is an ideal.

The ideal $P_\tau$ is prime because a product in $S$ equals $\tau$ only when one factor equals $\tau$.  For $I_J$, the quotient by the corresponding ring-reflection congruence has supported part $R/J$; equivalently, for supported $a,b$,
\[
ab\in I_J
\iff
ab\in J.
\]
Thus $I_J$ is prime exactly when $J$ is prime.
\end{proof}

\begin{theorem}[Generic-point suspension]\label{thm:generic-suspension}
The map
\[
i:\Spec R\longrightarrow\Spec G(R),
\qquad
\mathfrak p\longmapsto P_{\mathfrak p},
\]
is a homeomorphism onto the closed subspace
\[
V(e)=\{P_{\mathfrak p}:\mathfrak p\in\Spec R\}.
\]
Furthermore,
\[
D(e)=\{P_\tau\},
\]
and $P_\tau$ is open, dense, and contained in every prime ideal.  Hence $\Spec G(R)$ is obtained from $\Spec R$ by adjoining a new open dense generic point below every classical prime in the specialization order.
\end{theorem}

\begin{proof}
Since $e=0_R$ belongs to every ring prime $\mathfrak p$ but does not belong to $P_\tau$, the displayed descriptions of $V(e)$ and $D(e)$ follow from Theorem \ref{thm:prime-classification}.  For an ideal $J\triangleleft R$,
\[
i(V_R(J))
=
V_S(I_J)\cap V_S(e),
\]
so $i$ identifies the Zariski closed sets and is a homeomorphism onto $V(e)$.

The containment $P_\tau\subseteq P$ holds for every prime $P$ because every ideal contains the semiring zero $\tau$.  Therefore the closure of $\{P_\tau\}$ is all of $\Spec S$.  Since $D(e)=\{P_\tau\}$, this point is also open.
\end{proof}

\begin{corollary}[Local stalks]\label{cor:split-stalks}
For $\mathfrak p\in\Spec R$,
\[
G(R)_{P_\tau}\cong\B,
\qquad
G(R)_{P_{\mathfrak p}}\cong G(R_{\mathfrak p}).
\]
\end{corollary}

\begin{proof}
The complement of $P_\tau$ is the whole supported part $R$, which contains $e$, so Corollary \ref{cor:fraction-comparison} gives $\B$.  The complement of $P_{\mathfrak p}$ is $R\setminus\mathfrak p$, which does not contain $0$; Theorem \ref{thm:ring-localization} gives $G(R_{\mathfrak p})$.
\end{proof}

\begin{example}[The shifted generic point of $\Spec\Z$]
For $S=G(\Z)$,
\[
P_\tau=\{\tau\}
\subsetneq
P_0=\{\tau,0\}
\subsetneq
P_p=p\Z\cup\{\tau\}
\]
for every rational prime $p$.  Thus the classical zero ideal $(0)$ becomes the nonzero prime $P_0$, while $P_\tau$ is the new zero ideal and new generic point.
\end{example}

\subsection{The exact blueprint presentation}

Let $M_S$ be the multiplicative monoid of $S$, including its absorbing element $\tau$.  Let $\N_0[M_S]$ denote the contracted monoid semiring in which the basis element $[\tau]$ is identified with the additive zero.

\begin{theorem}[Multiplicative skeleton plus preaddition]\label{thm:blueprint-presentation}
The canonical map
\[
\N_0[M_S]\longrightarrow S,
\qquad
[x]\longmapsto x,
\]
induces an isomorphism
\[
\N_0[M_S]\Big/\Big\langle [a]+[b]\sim[a+b]
\;:\;a,b\in R\Big\rangle
\xrightarrow{\sim}
G(R).
\]
Consequently the associated blueprint consists of the multiplicative monoid $M_S$ together with precisely the additive relations inherited from $R$ and the contraction of $\tau$ to additive zero.
\end{theorem}

\begin{proof}
The canonical map is surjective.  The imposed relations hold in $G(R)$, so the map factors through the quotient.

Every element of the quotient is represented by a finite sum of monoid generators.  Terms $[\tau]$ vanish in the contracted monoid semiring.  If at least one supported term remains, repeated use of
\[
[a]+[b]\sim[a+b]
\]
reduces the sum to a single generator $[r]$ with $r\in R$.  If no supported term remains, the sum is zero, represented by $[\tau]$.  Thus every congruence class has a representative in $M_S$.  Two such representatives have the same image in $G(R)$ only when they are equal, because the reduction relations reproduce the actual ring sum and never identify $\tau$ with a supported element.  Hence the induced map is bijective.
\end{proof}

\begin{remark}[The role of $\mathbb F_1$]
The monoid $M_S$ is the multiplicative $\mathbb F_1$ shadow in the sense of monoidal geometry, while Theorem \ref{thm:blueprint-presentation} gives the exact blueprint that restores additive cancellation data; see \cite{Deitmar,Lorscheid}.  The Boolean quotient $\chi_R:S\to\B$ is the pure support completion.  The projective geometry below makes this skeleton visible as a finite Boolean face poset whose fibres carry ordinary projective geometry.
\end{remark}



\section{Conservative ideal arithmetic and the separated zeta sheet}
\label{sec:ideal-zeta}

Let $R$ be a commutative ring.  For an ideal $J\triangleleft R$, write
\[
I_J:=J\sqcup\{\tau\}\subseteq G(R).
\]
The baseline ideal classification in \cite{Globalization,Secondary} shows
that these are exactly the nontrivial supported ideals.

\begin{theorem}[Ideal multiplication is classical]
\label{thm:ideal-multiplication-classical}
For ideals $J,K\triangleleft R$,
\[
I_JI_K=I_{JK}.
\]
\end{theorem}

\begin{proof}
Every product of an element of $I_J$ with an element of $I_K$ is either
$\tau$ or a supported product $jk$, so $I_JI_K\subseteq I_{JK}$.  Conversely,
$\tau\in I_JI_K$, and every element of $JK$ is a finite sum of products
$j_ak_a$; the ideal product $I_JI_K$ contains each such product and is
closed under addition.  Hence $I_{JK}\subseteq I_JI_K$.
\end{proof}

Let $A=\mathcal O_K$ be the ring of integers of a number field.  Define the
\emph{normed split ideal class group} to be the group of nonzero ideals
$I_J$ modulo the principal supported ideals $I_{(a)}$.

\begin{corollary}[Class group and Dedekind zeta preservation]
\label{cor:dedekind-preservation}
The map $J\mapsto I_J$ induces canonical isomorphisms
\[
\operatorname{Cl}^{\mathrm{sp}}(G(A))\cong\operatorname{Cl}(A)
\]
and, on the ordinary quotient-norm sheet,
\[
\zeta^{\mathrm{ar}}_{G(A)}(s)
:=\sum_{0\ne J\triangleleft A}N(J)^{-s}
=\zeta_K(s).
\]
\end{corollary}

\begin{proof}
Theorem~\ref{thm:ideal-multiplication-classical} makes $J\mapsto I_J$ a
multiplicative bijection between nonzero ideals of $A$ and nontrivial
finite-norm supported ideals of $G(A)$; it carries $(a)$ to $I_{(a)}$.
This proves the class-group assertion.  The ordinary congruence quotient
collapsing $\tau$ with the supported zero is $A/J$, so its cardinality is
$N(J)$, giving the zeta identity term by term.
\end{proof}

The split quotient retaining the distinction between $\tau$ and $0_{A/J}$
is $G(A/J)$ and has cardinality $N(J)+1$.

\begin{theorem}[The separated weight is nonmultiplicative]
\label{thm:separated-weight}
For nonzero finite-norm ideals $J,K$ put
\[
w(J):=N(J)+1.
\]
Then
\[
w(JK)=N(J)N(K)+1
\]
while
\[
w(J)w(K)=N(J)N(K)+N(J)+N(K)+1.
\]
In particular $w$ is not multiplicative.
\end{theorem}

\begin{proof}
Use multiplicativity of the ordinary ideal norm and expand the product.
Since $N(J),N(K)>0$, the two expressions differ.
\end{proof}

\begin{corollary}[The support-separated sheet is non-Eulerian]
\label{cor:separated-non-eulerian}
The Dirichlet series
\[
\zeta^{\mathrm{sep}}_{G(A)}(s)
:=\sum_{0\ne J\triangleleft A}(N(J)+1)^{-s}
\]
has no ordinary Euler product indexed by the prime ideals of $A$.
For $A=\Z$,
\[
\zeta^{\mathrm{sep}}_{G(\Z)}(s)=\sum_{n\geq1}(n+1)^{-s}=\zeta(s)-1.
\]
\end{corollary}

\begin{proof}
An ordinary ideal Euler product gives multiplicative Dirichlet
coefficients on coprime ideals.  Theorem~\ref{thm:separated-weight} rules
this out.  The final identity is the change of index $m=n+1$.
\end{proof}

\begin{remark}
The two series answer different questions.  The normed sheet counts the
classical arithmetic quotient and preserves the Euler product.  The
separated sheet counts the extra unsupported class as a genuine point and
therefore destroys multiplicativity.  The absence of an Euler product is
not a defect in the calculation; it is the exact analytic cost of retaining
split support in quotient cardinalities.
\end{remark}

\section{Projective semimodules and finite distributive support}
\label{sec:projective-semimodules}

The semimodule equivalence of \cite{Secondary} sends a $G(R)$-semimodule
$M$ to its support join-semilattice $L_M=eM$, its $R$-module fibres
$M_\ell$, and the transport maps $T_{\ell,\ell'}$ for
$\ell\leq\ell'$.  The following result is the projective consequence that
can be obtained without imposing an unproved normal form on all transport
data.

\begin{theorem}[Necessary structure of a finitely generated projective]
\label{thm:projective-support-structure}
Let $P$ be a finitely generated projective $G(R)$-semimodule.  Then:
\begin{enumerate}[label=\textup{(\roman*)}]
\item $L_P=eP$ is a finite distributive lattice;
\item every fibre $P_\ell$ is a finitely generated projective $R$-module;
\item every transport map
\[
T_{\ell,\ell'}:P_\ell\longrightarrow P_{\ell'}
\qquad(\ell\leq\ell')
\]
is a split monomorphism of $R$-modules.
\end{enumerate}
If $R$ is a field, the function
$\ell\mapsto\dim_R P_\ell$ is therefore order-preserving.
\end{theorem}

\begin{proof}
Choose $G(R)$-linear maps
\[
i:P\longrightarrow G(R)^n,
\qquad
r:G(R)^n\longrightarrow P,
\qquad
r i=\id_P.
\]
Applying multiplication by $e$ gives join-homomorphisms
\[
i_s:L_P\longrightarrow\mathcal P([n]),
\qquad
r_s:\mathcal P([n])\longrightarrow L_P,
\qquad
r_s i_s=\id_{L_P}.
\]
Thus $L_P$ is a finite join-retract of a Boolean lattice.  For
$x,y\in L_P$, define
\[
x\wedge y:=r_s\bigl(i_s(x)\cap i_s(y)\bigr).
\]
If $z\leq x,y$, then $i_s(z)\subseteq i_s(x)\cap i_s(y)$, hence
$z\leq x\wedge y$; conversely $x\wedge y\leq x,y$ by monotonicity of
$r_s$.  This is the meet in $L_P$.  Since $i_s$ and $r_s$ preserve joins,
distributivity follows from distributivity in the Boolean lattice:
\[
\begin{aligned}
x\wedge(y\vee z)
&=r_s\bigl(i_s(x)\cap(i_s(y)\cup i_s(z))\bigr)\\
&=(x\wedge y)\vee(x\wedge z).
\end{aligned}
\]
This proves (i).

Fix $\ell\in L_P$ and write $I=i_s(\ell)$.  The maps $i$ and $r$
restrict to $R$-linear maps
\[
P_\ell\longrightarrow R^I\longrightarrow P_\ell
\]
whose composite is the identity.  Hence $P_\ell$ is a direct summand of a
finite free $R$-module, proving (ii).

For $\ell\leq\ell'$, put $I=i_s(\ell)$ and $I'=i_s(\ell')$; then
$I\subseteq I'$.  Let
\[
t_{I,I'}:R^I\hookrightarrow R^{I'}
\]
be extension by zero and let $\pi_{I',I}$ be coordinate projection.  Define
\[
s_{\ell',\ell}
:=r_\ell\circ\pi_{I',I}\circ i_{\ell'}:P_{\ell'}\to P_\ell.
\]
Naturality of $i$ with respect to transport gives
$i_{\ell'}T_{\ell,\ell'}=t_{I,I'}i_\ell$.  Therefore
\[
s_{\ell',\ell}T_{\ell,\ell'}
=r_\ell\pi_{I',I}t_{I,I'}i_\ell
=r_\ell i_\ell
=\id_{P_\ell}.
\]
This proves (iii).  The dimension assertion is immediate.
\end{proof}

\begin{remark}[Precision boundary for the decorated-poset claim]
\label{rem:projective-classification-boundary}
The intake corpus proposes the stronger normal form
$P_U\cong\bigoplus_{j\in U}Q_j$ over an ideal lattice
$\operatorname{Idl}(J)$.  Theorem~\ref{thm:projective-support-structure}
proves the distributive support, projective fibres, and split transports
needed for such a statement.  A complete classification still requires a
coherent-splitting theorem showing that the successive complements can be
chosen compatibly around every distributive square.  That compatibility is
not supplied by the abbreviated proof in the intake, and the stronger
``if and only if'' classification is therefore not used in this paper.
\end{remark}

\section{Unimodular split projective geometry}

For a semiring $A$ and $n\ge0$, a vector
\[
x=(x_0,\dots,x_n)\in A^{n+1}
\]
is called \emph{unimodular} if the ideal generated by its coordinates is the unit ideal.  Equivalently, there exist $a_0,\dots,a_n\in A$ such that
\[
a_0x_0+\cdots+a_nx_n=1.
\]

\begin{definition}[Unimodular projective point set]\label{def:unimodular-projective}
Define
\[
\Pum^n(A)
:=
\{x\in A^{n+1}:x\text{ is unimodular}\}/A^\times,
\]
where $A^\times$ acts by simultaneous scalar multiplication.
\end{definition}

This is the usual set of projective points over a field.  Over a general semiring it is deliberately denoted by $\Pum$ rather than identified without qualification with a scheme-theoretic $\Proj$.

Let $K$ be a field and set $S=G(K)$.  For $x\in S^{n+1}$ define
\[
\supp(x):=\{i:x_i\in K\}.
\]
This records supported coordinates, including coordinates equal to $e=0_K$.

\begin{lemma}[Unimodularity criterion]\label{lem:unimodularity-criterion}
A vector $x\in S^{n+1}$ is unimodular if and only if at least one coordinate belongs to $K^\times$.
\end{lemma}

\begin{proof}
If $x_i\in K^\times$, then $x_i^{-1}x_i=1$, so the coordinates generate the unit ideal.

Conversely, suppose every coordinate belongs to $\{\tau,e\}$.  Multiplication of such a coordinate by any scalar remains in $\{\tau,e\}$, and a finite sum of elements of $\{\tau,e\}$ again lies in $\{\tau,e\}$.  It therefore cannot equal $1$.  Hence the vector is not unimodular.
\end{proof}

The support character induces a map
\[
\Supp:\Pum^n(S)\longrightarrow\Pum^n(\B),
\qquad
[x_0:\cdots:x_n]
\longmapsto
[\chi_K(x_0):\cdots:\chi_K(x_n)].
\]
Since $\B^\times=\{1\}$,
\[
\Pum^n(\B)=\B^{n+1}\setminus\{0\}
\]
is canonically the set of nonempty subsets of $[n]_0=\{0,\dots,n\}$.

\begin{theorem}[Boolean-face decomposition]\label{thm:projective-face-decomposition}
For every field $K$ and $n\ge0$, there is a canonical decomposition
\[
\boxed{
\Pum^n(G(K))
\cong
\bigsqcup_{\varnothing\neq A\subseteq[n]_0}
\mathbf P^{|A|-1}(K).}
\]
Under this identification, the fibre of $\Supp$ over $A$ is $\mathbf P^{|A|-1}(K)$.
\end{theorem}

\begin{proof}
Fix a nonempty subset $A\subseteq[n]_0$.  A vector of support exactly $A$ has the form
\[
x_i\in K\quad(i\in A),
\qquad
x_i=\tau\quad(i\notin A).
\]
By Lemma \ref{lem:unimodularity-criterion}, it is unimodular exactly when the amplitude vector $(x_i)_{i\in A}\in K^A$ is nonzero.  The unit group of $S$ is $K^\times$, and simultaneous multiplication by a supported unit acts on the amplitude vector by the usual scalar action.  Hence the set of projective classes of support $A$ is
\[
(K^A\setminus\{0\})/K^\times
=\mathbf P^{|A|-1}(K).
\]
Distinct support masks cannot be related by multiplication by a supported unit, so the union is disjoint.  The description of the support fibre is immediate.
\end{proof}

\begin{corollary}[Amplitude and support projections]\label{cor:projective-projections}
There are canonical maps
\[
\Pum^n(G(K))\xrightarrow{\;p_*\;}\mathbf P^n(K)
\]
and
\[
\Pum^n(G(K))\xrightarrow{\;\Supp\;}\Pum^n(\B).
\]
The first forgets whether a zero coordinate was absent or supported; the second forgets all amplitudes and retains only the support face.
\end{corollary}

\begin{proof}
A unimodular split vector has at least one nonzero amplitude coordinate, so its coordinatewise ring reflection is a nonzero vector in $K^{n+1}$.  This defines $p_*$.  The support map was defined above.  Both are invariant under multiplication by $K^\times$.
\end{proof}

\subsection{The split projective line}

For $n=1$, Theorem \ref{thm:projective-face-decomposition} gives
\[
\Pum^1(G(K))
\cong
\mathbf P^1(K)\sqcup\{0_\tau,\infty_\tau\},
\]
where
\[
0_\tau:=[\tau:1],
\qquad
\infty_\tau:=[1:\tau].
\]
Inside the full-support stratum $\mathbf P^1(K)$ lie the supported zero and supported infinity
\[
0_e:=[e:1],
\qquad
\infty_e:=[1:e].
\]

\begin{proposition}[Four zero-pole states]\label{prop:four-zero-pole-states}
The ring-reflection map satisfies
\[
p_*(0_\tau)=p_*(0_e)=0,
\qquad
p_*(\infty_\tau)=p_*(\infty_e)=\infty.
\]
The support map satisfies
\[
\Supp(0_\tau)=[0:1],
\qquad
\Supp(\infty_\tau)=[1:0],
\]
while every point of the full-support stratum, including $0_e$ and $\infty_e$, maps to
\[
[1:1].
\]
\end{proposition}

\begin{proof}
This follows directly from
\[
p_K(\tau)=p_K(e)=0,
\qquad
\chi_K(\tau)=0,
\quad
\chi_K(e)=1.
\]
\end{proof}

Thus the Boolean projective line
\[
\Pum^1(\B)=\{[0:1],[1:1],[1:0]\}
\]
is a three-point incidence skeleton: absent zero, full-support torus, and absent infinity.

\section{Finite-field counts and split projective symmetry}

\subsection{Point counts and the support-projective zeta function}

\begin{theorem}[Finite-field point count]\label{thm:finite-field-point-count}
For every prime power $q$ and every $n\ge0$,
\[
\#\Pum^n(G(\F_q))
=
\sum_{r=1}^{n+1}
\binom{n+1}{r}\frac{q^r-1}{q-1}
=
\frac{(q+1)^{n+1}-2^{n+1}}{q-1}.
\]
\end{theorem}

\begin{proof}
There are $\binom{n+1}{r}$ support faces of cardinality $r$, and the fibre over each such face is $\mathbf P^{r-1}(\F_q)$, which has $(q^r-1)/(q-1)$ points.  This gives the first formula.  For the second,
\[
\sum_{r=1}^{n+1}\binom{n+1}{r}(q^r-1)
=
\bigl((q+1)^{n+1}-1\bigr)-\bigl(2^{n+1}-1\bigr),
\]
which equals $(q+1)^{n+1}-2^{n+1}$.
\end{proof}

\begin{definition}[Support-projective zeta function]\label{def:support-zeta}
Define
\[
Z_{\mathrm{sp}}(\Pum^n/\F_q,T)
:=
\exp\left(
\sum_{m\ge1}
\#\Pum^n(G(\F_{q^m}))\frac{T^m}{m}
\right).
\]
This is a point-counting zeta function for the split projective functor; no assertion is made here that $\Pum^n(G(\F_q))$ is the rational-point set of an ordinary scheme.
\end{definition}

\begin{theorem}[Closed zeta product]\label{thm:support-zeta-product}
Put
\[
c_{n,j}:=\sum_{r=j+1}^{n+1}\binom{n+1}{r}
\qquad(0\le j\le n).
\]
Then
\[
\#\Pum^n(G(\F_q))
=
\sum_{j=0}^{n}c_{n,j}q^j,
\]
and
\[
\boxed{
Z_{\mathrm{sp}}(\Pum^n/\F_q,T)
=
\prod_{j=0}^{n}(1-q^jT)^{-c_{n,j}}.}
\]
In particular,
\[
Z_{\mathrm{sp}}(\Pum^1/\F_q,T)
=
\frac{1}{(1-T)^3(1-qT)}.
\]
\end{theorem}

\begin{proof}
Since
\[
\frac{q^r-1}{q-1}=1+q+\cdots+q^{r-1},
\]
the coefficient of $q^j$ in the point count of Theorem \ref{thm:finite-field-point-count} is the sum of $\binom{n+1}{r}$ over $r\ge j+1$, which is $c_{n,j}$.

Substituting $q^m$ for $q$ gives
\[
\#\Pum^n(G(\F_{q^m}))
=
\sum_{j=0}^{n}c_{n,j}q^{mj}.
\]
Therefore
\[
\begin{aligned}
\log Z_{\mathrm{sp}}
&=
\sum_{m\ge1}\sum_{j=0}^{n}c_{n,j}q^{mj}\frac{T^m}{m}\\
&=
\sum_{j=0}^{n}c_{n,j}
\sum_{m\ge1}\frac{(q^jT)^m}{m}\\
&=
-\sum_{j=0}^{n}c_{n,j}\log(1-q^jT),
\end{aligned}
\]
which exponentiates to the product formula.  For $n=1$,
\[
c_{1,0}=\binom21+\binom22=3,
\qquad
c_{1,1}=\binom22=1.
\]
\end{proof}

\begin{remark}
The ordinary projective-line zeta function is
\[
Z(\mathbf P^1/\F_q,T)=\frac{1}{(1-T)(1-qT)}.
\]
The split support completion contributes two additional degree-zero factors, corresponding to $0_\tau$ and $\infty_\tau$.
\end{remark}

\subsection{Invertible matrices}

\begin{lemma}[Boolean invertible matrices]\label{lem:boolean-gl}
Every invertible $m\times m$ matrix over $\B$ is a permutation matrix.  Hence
\[
\GL_m(\B)\cong S_m.
\]
\end{lemma}

\begin{proof}
The free $\B$-semimodule $\B^m$ is the Boolean lattice of subsets of $\{1,\dots,m\}$, with addition given by union.  A $\B$-linear automorphism preserves joins and zero, hence is an automorphism of this finite Boolean lattice.  The atoms are exactly the standard basis vectors.  Every lattice automorphism permutes the atoms, so the corresponding matrix is a permutation matrix.  Conversely every permutation matrix is invertible.
\end{proof}

\begin{theorem}[Monomial general linear group]\label{thm:monomial-gl}
Let $K$ be a field and $S=G(K)$.  Then every matrix in $\GL_m(S)$ is monomial with one entry from $K^\times$ in each row and column and $\tau$ in every other position.  Consequently,
\[
\boxed{
\GL_m(G(K))\cong(K^\times)^m\rtimes S_m.}
\]
\end{theorem}

\begin{proof}
Let $U\in\GL_m(S)$ and let $V$ be its inverse.  Applying the support character entrywise gives
\[
\chi(U)\chi(V)=I_m,
\qquad
\chi(V)\chi(U)=I_m
\]
over $\B$.  By Lemma \ref{lem:boolean-gl}, $\chi(U)$ is a permutation matrix.  Therefore $U$ has exactly one supported entry in each row and column and has $\tau$ elsewhere.

Applying the ring reflection entrywise gives
\[
p(U)p(V)=I_m
\]
over $K$.  Thus $p(U)$ is an invertible monomial matrix, so each of its monomial entries is nonzero.  The corresponding supported entries of $U$ therefore lie in $K^\times$.

Conversely, a monomial matrix with entries in $K^\times$ has the evident monomial inverse.  The semidirect-product description is the standard decomposition into diagonal and permutation matrices.
\end{proof}

\begin{corollary}[Split projective linear group]\label{cor:split-pgl}
The group of projective transformations induced by invertible split matrices is
\[
\operatorname{PGL}^{\mathrm{lin}}_m(G(K))
\cong
\bigl((K^\times)^m/K^\times\bigr)\rtimes S_m,
\]
where the denominator is the diagonal scalar subgroup.  For $m=2$,
\[
\operatorname{PGL}^{\mathrm{lin}}_2(G(K))
\cong
K^\times\rtimes C_2.
\]
On the full-support affine coordinate $z$, the transformations are exactly
\[
z\longmapsto az
\qquad\text{and}\qquad
z\longmapsto\frac{a}{z},
\qquad a\in K^\times.
\]
\end{corollary}

\begin{proof}
Projectivization quotients $\GL_m(G(K))$ by the central diagonal copy of $K^\times$.  For $m=2$, a diagonal monomial matrix acts by scaling and an antidiagonal monomial matrix acts by scaling followed by inversion.
\end{proof}

\begin{corollary}[Obstruction to global translation]\label{cor:no-global-translation}
For $b\neq0$, the classical translation
\[
z\longmapsto z+b
\]
does not arise from a global split-projective linear automorphism of $\Pum^1(G(K))$.  More generally, a classical M\"obius transformation extends linearly to the mixed-support projective line only when its matrix is monomial after replacing structural zeros by $\tau$.
\end{corollary}

\begin{proof}
A matrix for translation is
\[
\begin{pmatrix}1&b\\0&1\end{pmatrix}.
\]
Its support pattern is not a permutation matrix, regardless of whether the classical zero entry is represented by $e$ or by $\tau$.  Theorem \ref{thm:monomial-gl} therefore excludes invertibility over $G(K)$.  The same argument applies to a general M\"obius matrix.
\end{proof}

\begin{remark}[Support-homogeneous affine maps]
Corollary \ref{cor:no-global-translation} does not say that affine maps disappear.  It says that they cannot be represented by the naive homogeneous coordinate $(x,1)$ on all support fibres.  The support-idempotent homogenization of \cite{Homogenization} replaces $(x,1)$ by $(x,h)$, where $h$ is the support idempotent, and represents
\[
x\longmapsto ax+b
\]
by
\[
(x,h)\longmapsto(ax+bh,h).
\]
This is the correct fibrewise extension across the Boolean support diagram.
\end{remark}

\section{Residue monads and Boolean-cubical infinitesimal geometry}

The support residue $\B$ is not itself Leibniz's infinitesimal monad.  A monad requires a residue map with a nontrivial kernel of infinitesimal amplitudes.  Split zero supplies the support geometry in which such kernels are organized.

\begin{definition}[Split residue datum]\label{def:split-residue-datum}
A \emph{split residue datum} is a surjective ring homomorphism
\[
q:A\twoheadrightarrow k
\]
with kernel
\[
\mathfrak m:=\ker q.
\]
It induces a semiring homomorphism
\[
G(q):G(A)\longrightarrow G(k)
\]
by
\[
G(q)(\tau)=\tau,
\qquad
G(q)(a)=q(a)
\quad(a\in A).
\]
\end{definition}

\begin{definition}[Supported monad]
For a supported element $a\in k\subset G(k)$ and a chosen lift $\widetilde a\in A$, define
\[
\mu_q(a):=G(q)^{-1}(a).
\]
For the absent state define
\[
\mu_q(\tau):=G(q)^{-1}(\tau).
\]
\end{definition}

\begin{theorem}[Monad fibre formula]\label{thm:monad-fibre}
Let $q:A\twoheadrightarrow k$ be a split residue datum.  Then
\[
\mu_q(a)=\widetilde a+\mathfrak m
\qquad(a\in k),
\]
and
\[
\mu_q(\tau)=\{\tau\}.
\]
In particular,
\[
\mu_q(e_k)=\mathfrak m,
\]
where $e_k=0_k$ is the supported zero in $G(k)$.
\end{theorem}

\begin{proof}
A supported element $x\in A\subset G(A)$ maps to $a$ exactly when $q(x)=a$, which is equivalent to $x\in\widetilde a+\mathfrak m$.  No supported element maps to $\tau$, because $G(q)$ preserves support.  Hence the only preimage of $\tau$ is $\tau$ itself.
\end{proof}

\begin{corollary}[Centre and halo]\label{cor:centre-and-halo}
The supported zero $e_k$ is the centre of the zero monad, the ideal $\mathfrak m$ is its infinitesimal halo, and $\tau$ is the absent state outside that halo.
\end{corollary}

\begin{proof}
The fibre over $e_k$ is $\mathfrak m$ by Theorem \ref{thm:monad-fibre}.  Its distinguished zero element is $0_A=e_A$, which maps to $e_k$.  The absent state has the singleton fibre $\{\tau\}$.
\end{proof}

The higher-dimensional construction is cubical.  Let $I$ be finite and let
\[
a=(a_i)_{i\in I}\in G(k)^I.
\]
Write
\[
J:=\supp(a)=\{i:a_i\in k\}.
\]
For $i\in J$, choose a lift $\widetilde a_i\in A$.

\begin{definition}[Support star and cubical monad]\label{def:cubical-monad}
The \emph{support star} of $a$ is
\[
\operatorname{Star}(a)
:=
\left\{
 b\in G(k)^I:
 b_i=a_i\ (i\in J),
\quad
b_i\in\{\tau,e_k\}\ (i\notin J)
\right\}.
\]
The \emph{upward cubical monad} of $a$ is
\[
\mu_q^\square(a)
:=
(G(q)^I)^{-1}(\operatorname{Star}(a)).
\]
\end{definition}

\begin{theorem}[Boolean-cubical monad decomposition]\label{thm:cubical-monad}
There is a canonical disjoint decomposition
\[
\boxed{
\mu_q^\square(a)
=
\bigsqcup_{L\subseteq I\setminus J}
\left(
\prod_{i\in J}(\widetilde a_i+\mathfrak m)
\right)
\times
\mathfrak m^L
\times
\{\tau\}^{I\setminus(J\cup L)}.}
\]
The subset $L$ records exactly the formerly absent coordinates that have been activated with residue zero.
\end{theorem}

\begin{proof}
Every $b\in\operatorname{Star}(a)$ is determined by the subset
\[
L:=\{i\in I\setminus J:b_i=e_k\}.
\]
For $i\in J$, Theorem \ref{thm:monad-fibre} gives the fibre $\widetilde a_i+\mathfrak m$.  For $i\in L$, the fibre of $e_k$ is $\mathfrak m$.  For the remaining coordinates, the fibre of $\tau$ is $\{\tau\}$.  Different subsets $L$ produce disjoint support patterns, so the union is disjoint.
\end{proof}

\begin{example}[The local path from absence to infinitesimal amplitude]
Take one coordinate and the absent base point $a=\tau$.  Then
\[
\mu_q^\square(\tau)
=
\{\tau\}\sqcup\mathfrak m.
\]
Inside the second component lie
\[
e_A=0_A
\]
and all nonzero infinitesimals $h\in\mathfrak m\setminus\{0\}$.  Thus the exact local sequence is
\[
\tau
\longrightarrow
 e_A
\longrightarrow
 h\in\mathfrak m\setminus\{0\},
\]
meaning absence, supported exact zero, and supported infinitesimal amplitude.
\end{example}

\subsection{First-order cubical tangent data}

Suppose now that $k$ is a field.  The ordinary first-order tangent space of the residue datum is the $k$-vector space
\[
T_q:=\mathfrak m/\mathfrak m^2.
\]
The cubical monad carries one copy of this vector space for every active coordinate.

\begin{theorem}[Boolean cube of first-order tangent fibres]\label{thm:first-order-cubical}
Let $q:A\twoheadrightarrow k$ be a surjection onto a field, let $a\in G(k)^I$ have support $J$, and fix a stratum indexed by $L\subseteq I\setminus J$ in Theorem \ref{thm:cubical-monad}.  Modulo coordinatewise congruence by $\mathfrak m^2$, that stratum is an affine torsor under
\[
T_q^{J\cup L}.
\]
As $L$ varies, the first-order split neighbourhood is a Boolean cube of affine tangent spaces indexed by $\mathcal P(I\setminus J)$.
\end{theorem}

\begin{proof}
For each $i\in J$, the coset $\widetilde a_i+\mathfrak m$ modulo $\mathfrak m^2$ is an affine torsor under $\mathfrak m/\mathfrak m^2$.  For each newly activated coordinate $i\in L$, the fibre $\mathfrak m$ modulo $\mathfrak m^2$ is the vector space $T_q$ itself.  Coordinates outside $J\cup L$ are fixed at $\tau$.  Taking products gives an affine torsor under $T_q^{J\cup L}$.  The indexing subsets $L$ form the Boolean cube of possible support activations.
\end{proof}

\begin{remark}
Theorem \ref{thm:first-order-cubical} separates two independent infinitesimal directions:
\[
\text{Boolean activation between support faces}
\]
and
\[
\text{linear tangent displacement inside an active face}.
\]
The element $e$ controls the first, while $\mathfrak m/\mathfrak m^2$ controls the second.
\end{remark}

\subsection{Formal completion}

Let $A$ be a ring and $I\triangleleft A$ an ideal.  The quotient maps $A/I^{r+1}\to A/I^r$ induce an inverse system of split-zero semirings.

\begin{theorem}[Split zero commutes with adic completion]\label{thm:completion-commutes}
There is a canonical semiring isomorphism
\[
\varprojlim_{r\ge1}G(A/I^r)
\xrightarrow{\sim}
G\left(\varprojlim_{r\ge1}A/I^r\right).
\]
In particular, if $A$ is $I$-adically complete, then
\[
\varprojlim_{r\ge1}G(A/I^r)
\cong
G(A).
\]
\end{theorem}

\begin{proof}
An element of the inverse limit on the left is a coherent sequence $(x_r)_{r\ge1}$.  Every transition map sends $\tau$ to $\tau$ and sends supported elements to supported elements.  Hence either all $x_r$ equal $\tau$, or every $x_r$ is supported.  In the first case the sequence corresponds to the external zero of the right-hand side.  In the second case the sequence is exactly a coherent sequence in the ring quotients $A/I^r$, hence an element of $\varprojlim A/I^r$, regarded as supported.  This gives a bijection, and the operations are preserved coordinatewise.
\end{proof}

\begin{definition}[Split formal monad]
The split formal neighbourhood associated with $(A,I)$ is the pro-semiring
\[
\operatorname{Spf}_{\mathrm{sp}}(A,I)
:=
\{G(A/I^r)\}_{r\ge1}.
\]
Its supported sector is the ordinary formal neighbourhood; its support masks form the Boolean incidence skeleton.
\end{definition}

\section{Hyperreal and surcomplex realizations}

\subsection{Finite hyperreals}

Let ${}^{\ast}\R$ be a nonstandard elementary extension of $\R$.  Let
\[
\cO_{\mathrm{hyp}}
:=
\{x\in{}^{\ast}\R:|x|<n\text{ for some }n\in\N\}
\]
be the ring of finite hyperreals and let
\[
\mathfrak m_{\mathrm{hyp}}
:=
\{x\in{}^{\ast}\R:|x|<1/n\text{ for every }n\ge1\}
\]
be its ideal of infinitesimals.  The standard-part map is a surjective ring homomorphism
\[
\st:\cO_{\mathrm{hyp}}\twoheadrightarrow\R
\]
with kernel $\mathfrak m_{\mathrm{hyp}}$; see \cite{Robinson,Goldblatt}.

\begin{corollary}[Split hyperreal monads]\label{cor:hyperreal-monads}
For every $a\in\R$,
\[
\mu_{\mathrm{hyp}}(a)
=
a+\mathfrak m_{\mathrm{hyp}}.
\]
For a mixed-support point $a\in G(\R)^I$, the upward monad is the Boolean cube of ordinary hyperreal monads described by Theorem \ref{thm:cubical-monad}.
\end{corollary}

\begin{proof}
Apply Theorems \ref{thm:monad-fibre} and \ref{thm:cubical-monad} to the standard-part map.
\end{proof}

\subsection{Surreal and surcomplex normal forms}

We work in von Neumann--Bernays--G\"odel set theory so that the proper class $\No$ may be treated as a class-sized ordered field.  Every support occurring in a Conway normal form is a set.  We use Conway's normal-form theorem and the standard monomial map
\[
x\longmapsto\omega^x
\]
from \cite{Conway}.  The uploaded expository notes \cite{Simons} give a concise account of birthdays, ordinal embeddings, and the first infinitesimal $\omega^{-1}$.

Set
\[
K:=\No[i]
=
\{x+iy:x,y\in\No,\ i^2=-1\}.
\]
Conway proved that $K$ is algebraically closed \cite{Conway}.

\begin{theorem}[Surcomplex Conway normal form]\label{thm:surcomplex-normal-form}
Every nonzero $z\in K$ has a unique expansion
\[
z=
\sum_{\rho<\lambda}
 c_\rho\omega^{\gamma_\rho},
\]
where
\[
c_\rho\in\C^\times,
\qquad
\gamma_0>\gamma_1>\cdots,
\]
and the support $\{\gamma_\rho:\rho<\lambda\}$ is reverse well ordered.  Equivalently, $K$ is the class-sized Hahn field with coefficient field $\C$ and value group $\No$, with set-sized reverse-well-ordered support.
\end{theorem}

\begin{proof}
Write $z=x+iy$ with $x,y\in\No$.  By Conway's normal-form theorem,
\[
x=\sum_{\gamma\in A}a_\gamma\omega^\gamma,
\qquad
 y=\sum_{\gamma\in B}b_\gamma\omega^\gamma,
\]
where $A$ and $B$ are reverse well ordered and the coefficients are real.  The finite union $A\cup B$ is reverse well ordered.  Extending missing coefficients by zero gives
\[
z=
\sum_{\gamma\in A\cup B}
(a_\gamma+ib_\gamma)\omega^\gamma.
\]
Delete zero coefficients.  Uniqueness follows by comparing real and imaginary parts and using uniqueness of the real Conway normal form.
\end{proof}

For nonzero $z$, let $\ell(z)$ be its leading exponent and define the natural valuation
\[
v(z):=-\ell(z),
\qquad
v(0):=\infty.
\]

\begin{proposition}[Valuation laws]\label{prop:surcomplex-valuation}
For $z,w\in K$,
\[
v(zw)=v(z)+v(w),
\]
and
\[
v(z+w)\ge\min\{v(z),v(w)\},
\]
with equality unless the leading terms cancel.  If $z=x+iy$, then
\[
v(z)=\min\{v(x),v(y)\}.
\]
\end{proposition}

\begin{proof}
If
\[
z=c\omega^\alpha+\text{lower terms},
\qquad
w=d\omega^\beta+\text{lower terms},
\]
then
\[
zw=cd\,\omega^{\alpha+\beta}+\text{lower terms},
\]
and $cd\neq0$.  Thus $\ell(zw)=\alpha+\beta$, proving multiplicativity of $v$.  The addition inequality is the standard leading-term argument.

For $z=x+iy$, let the leading real exponents of $x$ and $y$ be $\alpha$ and $\beta$.  If $\alpha\neq\beta$, the larger exponent survives in $z$.  If $\alpha=\beta$, the leading complex coefficient is $a+ib$ with $a,b\in\R$ not both zero, so there is no cancellation between real and imaginary components.  Hence $\ell(z)=\max\{\alpha,\beta\}$, which is equivalent to the displayed valuation formula.
\end{proof}

Define the finite valuation ring and its maximal ideal by
\[
\cO_K:=\{z\in K:v(z)\ge0\},
\qquad
\mathfrak m_K:=\{z\in K:v(z)>0\}.
\]

\begin{theorem}[Surcomplex standard part]\label{thm:surcomplex-standard-part}
The coefficient of $\omega^0$ defines a surjective ring homomorphism
\[
\st:\cO_K\twoheadrightarrow\C
\]
with kernel $\mathfrak m_K$.  Hence
\[
\cO_K/\mathfrak m_K\cong\C.
\]
\end{theorem}

\begin{proof}
An element of $\cO_K$ has no positive exponent in its normal form.  Define $\st(z)$ to be the coefficient of $\omega^0$, taking that coefficient to be zero if the exponent $0$ is absent.  Addition preserves constant coefficients.  In a product of two finite series, an exponent sum can equal zero only when both exponents are zero, because all exponents are nonpositive.  Therefore the constant coefficient of a product is the product of the constant coefficients, so $\st$ is a ring homomorphism.  It is surjective because every $c\in\C$ is the constant series $c\omega^0$.  Its kernel consists exactly of series whose leading exponent is negative, equivalently those of positive valuation.
\end{proof}

\begin{corollary}[The surcomplex Leibniz monad]\label{cor:surcomplex-monad}
For every $c\in\C$,
\[
\mu_{\mathrm{sur}}(c)
=
c+\mathfrak m_K.
\]
In particular, the zero monad is
\[
\mu_{\mathrm{sur}}(0)=\mathfrak m_K.
\]
After split-zero globalization, the upward mixed-support monad is the Boolean cube of these surcomplex fibres.
\end{corollary}

\begin{proof}
Apply Theorems \ref{thm:monad-fibre} and \ref{thm:cubical-monad} to the standard-part map of Theorem \ref{thm:surcomplex-standard-part}.
\end{proof}

\subsection{Ordinal scales and reciprocal infinitesimals}

Every ordinal embeds canonically in $\No$.  Inside the ordered field $\No$, its arithmetic is the natural or Hessenberg arithmetic induced by Conway normal forms, not the noncommutative ordinal addition and multiplication.

\begin{proposition}[Reciprocals of infinite ordinals]\label{prop:ordinal-reciprocals}
If $\alpha$ is an infinite ordinal, regarded as a positive surreal number, then
\[
\alpha^{-1}
\]
is a positive surreal infinitesimal.  The map
\[
\alpha\longmapsto\alpha^{-1}
\]
is order reversing on positive ordinals.  Since $\omega$ is the least infinite ordinal, $\omega^{-1}$ is the largest reciprocal of an infinite ordinal, although it is not the largest positive surreal infinitesimal.
\end{proposition}

\begin{proof}
For every positive integer $n$,
\[
\alpha>n.
\]
Inversion in an ordered field reverses inequalities among positive elements, so
\[
0<\alpha^{-1}<n^{-1}.
\]
This is the definition of a positive infinitesimal.  Order reversal is the same argument.  The final assertion follows from minimality of $\omega$ among infinite ordinals.  There is no largest positive infinitesimal because $2\epsilon$ is infinitesimal and larger than $\epsilon$ whenever $\epsilon>0$ is infinitesimal.
\end{proof}

\begin{theorem}[Reciprocal ordinals are coinitial]\label{thm:reciprocal-coinitial}
For every positive surreal infinitesimal $\epsilon$, there exists an infinite ordinal $\alpha$ such that
\[
0<\alpha^{-1}<\epsilon.
\]
\end{theorem}

\begin{proof}
The reciprocal $x:=\epsilon^{-1}$ is positive infinite.  The ordinals are cofinal in $\No$; equivalently, every surreal number is bounded above by an ordinal \cite[Chap.~5]{Conway}.  Choose an ordinal $\alpha>x$.  Then $\alpha$ is infinite and inversion gives
\[
0<\alpha^{-1}<x^{-1}=\epsilon.
\]
\end{proof}

\begin{lemma}[Geometric inversion in the Hahn field]\label{lem:hahn-geometric}
If $u\in K$ satisfies $v(u)>0$, then the family $\{(-u)^n:n\ge0\}$ is Hahn-summable and
\[
(1+u)^{-1}=\sum_{n=0}^{\infty}(-u)^n.
\]
\end{lemma}

\begin{proof}
Since $v(u^n)=nv(u)$ tends strictly upward, for every valuation bound only finitely many terms have valuation below that bound.  This is precisely the Hahn summability condition.  Multiplying partial sums gives
\[
(1+u)\sum_{n=0}^{N}(-u)^n
=1+(-1)^Nu^{N+1}.
\]
The remainder tends to arbitrarily high valuation, so the Hahn sum is an inverse of $1+u$.
\end{proof}

\begin{theorem}[Reciprocal from Cantor normal form]\label{thm:ordinal-inverse-expansion}
Let an infinite ordinal have Cantor normal form
\[
\alpha
=
n_0\omega^{\beta_0}
+n_1\omega^{\beta_1}
+\cdots
+n_m\omega^{\beta_m},
\]
where
\[
\beta_0>\beta_1>\cdots>\beta_m,
\qquad
n_j\in\N_{>0}.
\]
Set
\[
u:=
\sum_{j=1}^{m}
\frac{n_j}{n_0}\omega^{\beta_j-\beta_0}.
\]
Then $u$ is infinitesimal and
\[
\boxed{
\alpha^{-1}
=
\frac1{n_0}\omega^{-\beta_0}
\sum_{k=0}^{\infty}(-u)^k.}
\]
Its leading term is
\[
\frac1{n_0}\omega^{-\beta_0},
\]
so
\[
v(\alpha^{-1})=\beta_0.
\]
\end{theorem}

\begin{proof}
Inside $\No$, Cantor normal form is Conway normal form for ordinals.  Factor
\[
\alpha
=
n_0\omega^{\beta_0}(1+u).
\]
Each exponent $\beta_j-\beta_0$ is negative, so $v(u)>0$.  Lemma \ref{lem:hahn-geometric} gives the inverse of $1+u$, and multiplication by $(n_0\omega^{\beta_0})^{-1}$ gives the formula.  Since the geometric factor has leading term $1$, the leading term of $\alpha^{-1}$ is the displayed monomial.
\end{proof}

\begin{example}
For $\alpha=\omega+1$,
\[
\frac1{\omega+1}
=
\omega^{-1}
-\omega^{-2}
+\omega^{-3}
-\omega^{-4}
+\cdots.
\]
This is an exact surreal equality.
\end{example}

\subsection{Complex phase at every infinitesimal order}

For $\lambda\in\No$, define
\[
F_{\ge\lambda}:=\{z\in K:v(z)\ge\lambda\},
\qquad
F_{>\lambda}:=\{z\in K:v(z)>\lambda\}.
\]

\begin{theorem}[Complex graded layers]\label{thm:complex-graded-layers}
For every $\lambda\in\No$, coefficient extraction at the monomial $\omega^{-\lambda}$ induces a canonical $\C$-linear isomorphism
\[
F_{\ge\lambda}/F_{>\lambda}
\xrightarrow{\sim}
\C.
\]
Multiplication of initial terms gives bilinear maps
\[
\gr_\lambda(K)\times\gr_\mu(K)
\longrightarrow
\gr_{\lambda+\mu}(K)
\]
corresponding to multiplication in $\C$.
\end{theorem}

\begin{proof}
For $z\in F_{\ge\lambda}$, let $c_{-\lambda}(z)$ be the coefficient of $\omega^{-\lambda}$, taking it to be zero when that monomial is absent.  This is a $\C$-linear map to $\C$.  Its kernel consists exactly of elements with valuation strictly greater than $\lambda$, namely $F_{>\lambda}$.  It is surjective because $c\omega^{-\lambda}$ maps to $c$.  The multiplication statement follows from
\[
(c\omega^{-\lambda})(d\omega^{-\mu})
=cd\,\omega^{-(\lambda+\mu)}.
\]
\end{proof}

\begin{definition}[Support--order--phase initial datum]\label{def:support-order-phase}
For $x\in G(K)$ define
\[
\operatorname{in}_{\mathrm{sp}}(x)
=
\begin{cases}
(0,\bot,0),&x=\tau,\\
(1,\infty,0),&x=e,\\
(1,v(x),c),&x\in K^\times,
\end{cases}
\]
where $c\in\C^\times$ is the leading coefficient of $x$.  Here $\bot$ denotes absence and $\infty$ denotes supported exact zero.
\end{definition}

\begin{remark}
The three coordinates have independent meanings:
\[
\text{support}\in\B,
\qquad
\text{order}\in\No\cup\{\infty\},
\qquad
\text{phase}\in\C.
\]
The split-zero residue retains the first, the valuation retains the second, and the associated graded retains the third.
\end{remark}

\begin{corollary}[Mixed surreal division]\label{cor:mixed-surreal-division}
Let $\alpha>0$ be an ordinal and let $S=G(\No[i])$.  Then
\[
S^2[(\omega^\alpha,e)^{-1}]
\cong
S\times\B,
\]
and
\[
\frac{(a,b)}{(\omega^\alpha,e)}
=
\bigl(a\omega^{-\alpha},\chi(b)\bigr).
\]
Similarly,
\[
\frac{(a,b)}{(\omega^{-\alpha},e)}
=
\bigl(a\omega^{\alpha},\chi(b)\bigr).
\]
\end{corollary}

\begin{proof}
Every nonzero surreal or surcomplex number is a unit of $K$.  Apply Theorem \ref{thm:mixed-localization}.
\end{proof}

\begin{corollary}[Strictly identifying $e$ with a surreal infinitesimal]
For every nonzero $h\in\No[i]$, including $h=\omega^{-\alpha}$,
\[
G(\No[i])/(e=h)\cong\B.
\]
On a selected coordinate face $A$, the corresponding quotient is the partial Booleanization of Theorem \ref{thm:scoped-identity}.
\end{corollary}

\begin{proof}
Apply Theorems \ref{thm:strict-identification-collapse} and \ref{thm:scoped-identity}.
\end{proof}

\subsection{Differential separation of \texorpdfstring{$e$ and $\omega^{-1}$}{e and omega inverse}}

Berarducci and Mantova construct a simplest surreal derivation
\[
\partial:\No\longrightarrow\No
\]
with constant field $\R$, strong additivity, compatibility with exponentiation, and
\[
\partial(\omega)=1;
\]
see \cite{BerarducciMantova}.  Extend it $\C$-linearly to $K=\No[i]$ by $\partial(i)=0$.

\begin{proposition}[Differential obstruction to identification]\label{prop:differential-obstruction}
One has
\[
\partial(\omega^{-1})=-\omega^{-2}\neq0,
\]
whereas every semiring derivation on $G(K)$ kills the supported zero $e$.  Hence no derivation-preserving identification can send $e$ to $\omega^{-1}$.
\end{proposition}

\begin{proof}
Differentiate $\omega\omega^{-1}=1$:
\[
0
=
\partial(1)
=
\partial(\omega)\omega^{-1}
+
\omega\partial(\omega^{-1})
=
\omega^{-1}
+
\omega\partial(\omega^{-1}).
\]
Multiplication by $\omega^{-1}$ gives
\[
\partial(\omega^{-1})=-\omega^{-2}.
\]
The statement for $e$ is Proposition \ref{prop:derivation-kills-e}.
\end{proof}

\section{Projective zero, projective infinity, and the surcomplex circle}

\subsection{Inversion and the two poles}

Let $K$ be a field and $S=G(K)$.  Consider the matrix
\[
J:=
\begin{pmatrix}
\tau&1\\
1&\tau
\end{pmatrix}.
\]

\begin{proposition}[Global split inversion]\label{prop:global-inversion}
The matrix $J$ belongs to $\GL_2(S)$ and satisfies $J^2=I$.  Its action on $\Pum^1(S)$ is
\[
[x:y]\longmapsto[y:x].
\]
In particular,
\[
0_\tau\longleftrightarrow\infty_\tau,
\qquad
0_e\longleftrightarrow\infty_e.
\]
\end{proposition}

\begin{proof}
A direct multiplication gives
\[
J^2
=
\begin{pmatrix}
\tau^2+1&\tau+\tau\\
\tau+\tau&1+\tau^2
\end{pmatrix}
=
\begin{pmatrix}
1&\tau\\
\tau&1
\end{pmatrix}
=I.
\]
The coordinate action and the stated exchanges are immediate.
\end{proof}

\begin{theorem}[Equivalence of zero and infinity monads]\label{thm:zero-infinity-monads}
Let $q:A\twoheadrightarrow k$ be a split residue datum and let $S_A=G(A)$.  In the affine chart $[x:1]$, the upward zero monad is
\[
\{[\tau:1]\}
\sqcup
\{[h:1]:h\in\mathfrak m\}.
\]
Projective inversion restricts to a bijection onto the upward infinity monad
\[
\{[1:\tau]\}
\sqcup
\{[1:h]:h\in\mathfrak m\}.
\]
It sends
\[
[\tau:1]\mapsto[1:\tau],
\qquad
[e:1]\mapsto[1:e],
\]
and, for $h\neq0$,
\[
[h:1]\mapsto[1:h]=[h^{-1}:1]
\]
whenever $A$ is contained in a field in which $h$ is invertible.
\end{theorem}

\begin{proof}
The description of the upward zero monad is Theorem \ref{thm:cubical-monad} in one coordinate, written in the affine projective chart.  Proposition \ref{prop:global-inversion} swaps the homogeneous coordinates and gives the displayed image.  If $h$ is nonzero in an ambient field, multiplication by the unit $h^{-1}$ identifies $[1:h]$ with $[h^{-1}:1]$.
\end{proof}

\begin{corollary}[No equation $x=\infty$ is required]\label{cor:no-x-infinity}
The supported projective infinity
\[
\infty_e=[1:e]
\]
and the absent-denominator infinity
\[
\infty_\tau=[1:\tau]
\]
are points of $\Pum^1(G(K))$.  Infinite field elements form the punctured valuative monad of $\infty_e$ under inversion; they are not obtained by adjoining an arithmetic element satisfying $x=\infty$.
\end{corollary}

\begin{proof}
The two points are distinct by their support masks.  The punctured-monad assertion is Theorem \ref{thm:zero-infinity-monads} applied to a valued field: nonzero infinitesimals invert to infinite elements.
\end{proof}

\subsection{The Cayley circle}

Let $F$ be a real closed field, let $F[i]$ be its algebraic closure, and define
\[
U_F:=\{z\in F[i]:z\overline z=1\}.
\]

\begin{theorem}[Algebraic Cayley circle]\label{thm:cayley-circle}
The Cayley transform
\[
C:\mathbf P^1(F)\longrightarrow U_F
\]
defined by
\[
C(x)=\frac{x-i}{x+i}
\qquad(x\in F),
\qquad
C(\infty)=1,
\]
is a bijection.  Its inverse is
\[
C^{-1}(z)=i\frac{1+z}{1-z}
\qquad(z\neq1),
\qquad
C^{-1}(1)=\infty.
\]
In particular,
\[
\mathbf P^1(\No)
\cong
\{z\in\No[i]:z\overline z=1\}.
\]
\end{theorem}

\begin{proof}
For $x\in F$,
\[
C(x)\overline{C(x)}
=
\frac{x-i}{x+i}\frac{x+i}{x-i}=1.
\]
Conversely, let $z\in U_F$ with $z\neq1$.  Since $\overline z=z^{-1}$,
\[
\overline{\left(i\frac{1+z}{1-z}\right)}
=
-i\frac{1+z^{-1}}{1-z^{-1}}
=
-i\frac{z+1}{z-1}
=
i\frac{1+z}{1-z}.
\]
Thus $C^{-1}(z)$ is fixed by conjugation and lies in $F$.  Direct substitution shows that the formulas are inverse.  The final statement uses that $\No$ is real closed.
\end{proof}

\begin{proposition}[Cayley obstruction on mixed support]\label{prop:cayley-obstruction}
The Cayley transform acts on the full-support classical stratum but is not induced by a global split-projective linear automorphism of $\Pum^1(G(F[i]))$.
\end{proposition}

\begin{proof}
A matrix for the Cayley transform is
\[
\begin{pmatrix}
1&-i\\
1&i
\end{pmatrix}.
\]
Every entry is supported, so its Boolean support matrix is
\[
\begin{pmatrix}
1&1\\
1&1
\end{pmatrix},
\]
which is not invertible over $\B$.  Theorem \ref{thm:monomial-gl} therefore excludes invertibility over the split semiring.  Over the field $F[i]$, the determinant is nonzero and the classical action is valid on the full-support stratum.
\end{proof}

\subsection{Placing \texorpdfstring{$\omega$}{omega} at a projective pole}

In the classical surcomplex projective line, the M\"obius transformation
\[
u=\frac{x}{\omega-x}
\]
is induced by
\[
[X:Y]\longmapsto[X:\omega Y-X].
\]
It sends $x=0$ to $u=0$ and $x=\omega$ to $u=\infty$.  This realizes $\omega$ as a projective pole rather than as a new arithmetic zero.

\begin{proposition}[The $\omega$-pole chart is not globally split-linear]\label{prop:omega-pole-obstruction}
The transformation $x\mapsto x/(\omega-x)$ is a classical projective automorphism of $\mathbf P^1(\No[i])$, but it does not extend as a global linear automorphism of $\Pum^1(G(\No[i]))$.
\end{proposition}

\begin{proof}
Its classical determinant is $\omega\neq0$, so it is a field-projective automorphism.  Any representing split matrix has a support pattern with at least three supported entries, hence its Boolean support matrix is not a permutation matrix.  Theorem \ref{thm:monomial-gl} excludes global split invertibility.
\end{proof}

\begin{remark}[Birthdays]
In Conway's construction, $0=\{\mid\}$ is born on day $0$, while $\omega^{-1}$ is born at a transfinite stage; see \cite{Conway,Simons}.  The external state $\tau$ is not an earlier surreal number.  One may assign it a separate support grade $\bot$ below all birthdays, but this is an enrichment of the birthday filtration, not a modification of Conway's recursive order.
\end{remark}

\section{Limits, poles, and removable indeterminacies}

The split formalism does not replace all limit theories by a single operation.  It separates a limit problem into an amplitude condition inside a monad and a support condition across the Boolean incidence skeleton.

\begin{definition}[Residue-monadic limit]\label{def:monadic-limit}
Let $q:A\twoheadrightarrow k$ be a split residue datum, let $a,L\in k$, choose a lift $\widetilde a\in A$, and let
\[
f:(\widetilde a+\mathfrak m)\setminus\{\widetilde a\}\longrightarrow A
\]
be a function.  We say that $f$ has \emph{residue-monadic limit} $L$ at $a$ if
\[
q(f(\widetilde a+h))=L
\]
for every nonzero $h\in\mathfrak m$ for which $f$ is defined.
\end{definition}

For finite hyperreals and an internal extension of a real function, this is the standard nonstandard criterion for a limit; see \cite{Robinson,Goldblatt}.  For the surcomplex valuation ring it states that all infinitesimal perturbations have the same complex standard part.

\begin{proposition}[Projective reciprocal extension]\label{prop:reciprocal-extension}
For every field $K$, the map
\[
\mathcal R:G(K)\longrightarrow\Pum^1(G(K)),
\qquad
x\longmapsto[1:x],
\]
extends ordinary reciprocal projectively.  It satisfies
\[
\mathcal R(h)=[h^{-1}:1]
\qquad(h\in K^\times),
\]
\[
\mathcal R(e)=\infty_e,
\qquad
\mathcal R(\tau)=\infty_\tau.
\]
\end{proposition}

\begin{proof}
If $h\in K^\times$, simultaneous multiplication by $h^{-1}$ gives
\[
[1:h]=[h^{-1}:1].
\]
The two zero cases are the definitions of $\infty_e$ and $\infty_\tau$.
\end{proof}

Thus a pole at supported zero and an absent denominator have distinct projective values.

\begin{theorem}[Resolution of a common-factor indeterminacy]\label{thm:common-factor-resolution}
Let $K$ be a field, let $a\in K$, and let $P,Q\in K[T]$.  Suppose
\[
P(T)=(T-a)^rP_0(T),
\qquad
Q(T)=(T-a)^rQ_0(T)
\]
for some $r\ge1$, with
\[
Q_0(a)\neq0.
\]
Then the rational projective map
\[
\Phi(x):=[P(x):Q(x)]
\]
agrees for $x\neq a$ with
\[
\widetilde\Phi(x):=[P_0(x):Q_0(x)],
\]
and $\widetilde\Phi$ extends to $x=a$ with value
\[
\widetilde\Phi(a)=[P_0(a):Q_0(a)].
\]
In $G(K)$, the unresolved pair at $x=a$ is
\[
[e:e],
\]
which has full support but is non-unimodular.  It is distinct from the absent pair
\[
[\tau:\tau].
\]
\end{theorem}

\begin{proof}
For $x\neq a$, the scalar $(x-a)^r$ is a nonzero unit of $K$, so
\[
[P(x):Q(x)]
=
[(x-a)^rP_0(x):(x-a)^rQ_0(x)]
=
[P_0(x):Q_0(x)].
\]
At $x=a$, the second coordinate $Q_0(a)$ is a unit, so $[P_0(a):Q_0(a)]$ is unimodular and defines a projective point.  Before cancellation, both $P(a)$ and $Q(a)$ are the supported ring zero $e$, giving the full-support but non-unimodular pair $(e,e)$.  The pair $(\tau,\tau)$ has empty support, so the two are distinct in $G(K)^2$.
\end{proof}

\begin{example}[The form $0/0$ at $x=1$]
For
\[
f(x)=\frac{x^2-1}{x-1},
\]
Theorem \ref{thm:common-factor-resolution} gives
\[
[x^2-1:x-1]=[x+1:1]
\]
away from $x=1$, and the resolved value is
\[
[2:1].
\]
If $h$ is a nonzero hyperreal or surcomplex infinitesimal, then
\[
f(1+h)=\frac{2h+h^2}{h}=2+h,
\]
so its standard part is $2$.  The Newtonian limit, the Leibniz infinitesimal computation, and the projective common-factor resolution therefore give the same amplitude result, while split support additionally distinguishes $(e,e)$ from $(\tau,\tau)$.
\end{example}

\begin{theorem}[Two-component split limit datum]\label{thm:two-component-limit}
For a map defined on a mixed-support punctured neighbourhood, a split limit consists of two independent conditions:
\begin{enumerate}[label=\textup{(\alph*)}]
\item constancy of the residue or standard part on each active infinitesimal fibre;
\item compatibility of the induced support map on the Boolean star of the base point.
\end{enumerate}
Neither condition implies the other.
\end{theorem}

\begin{proof}
The cubical decomposition of Theorem \ref{thm:cubical-monad} separates a neighbourhood into disjoint strata indexed by support activation masks.  Within each stratum, the residue map forgets infinitesimal amplitudes and tests condition (a).  The support character forgets all amplitudes and tests condition (b).  A function can preserve residues while changing a support mask, or preserve support while changing residues, so the conditions are independent.
\end{proof}

\section{Identity levels and univalence}

There are three mathematically distinct senses in which two infinitesimal objects may be called identical.

\begin{definition}[Three identity levels]\label{def:identity-levels}
Let $S=G(K)$.
\begin{enumerate}[label=\textup{(\roman*)}]
\item \emph{Arithmetic identity} is equality in $S$ or in a specified semiring quotient.
\item \emph{Support-modal identity} is equality after applying $\chi_K:S\to\B$.
\item \emph{Structural identity} is isomorphism or equivalence of pointed monads, formal neighbourhoods, or projective local geometries.
\end{enumerate}
\end{definition}

\begin{proposition}[Arithmetic and modal identities]\label{prop:arithmetic-modal}
For every $h\in K^\times$,
\[
e\neq h
\]
in $S$, but
\[
\chi_K(e)=\chi_K(h)=1.
\]
Freely adjoining the arithmetic identity $e=h$ yields $\B$, while adjoining it only on coordinates in $A$ yields
\[
\B^A\times S^{I\setminus A}.
\]
\end{proposition}

\begin{proof}
The inequality is set-theoretic because $e=0_K$ and $h\neq0$.  The support equality follows from the definition of $\chi_K$.  The quotient statements are Theorems \ref{thm:strict-identification-collapse} and \ref{thm:scoped-identity}.
\end{proof}

\begin{proposition}[Structural identity of zero and infinity monads]\label{prop:structural-identity}
Projective inversion gives an isomorphism of pointed split monads
\[
(\mu^\square(0_e),0_e)
\xrightarrow{\sim}
(\mu^\square(\infty_e),\infty_e)
\]
and likewise for the unsupported points $0_\tau$ and $\infty_\tau$.
\end{proposition}

\begin{proof}
This is Theorem \ref{thm:zero-infinity-monads}, with the distinguished centres included in the data.
\end{proof}

\begin{theorem}[Univalent interpretation]\label{thm:univalent-interpretation}
Assume a univalent universe containing the type of pointed split monads, with equivalences as its notion of structure-preserving equivalence.  Then the equivalence of Proposition \ref{prop:structural-identity} determines an identity path
\[
(\mu^\square(0_e),0_e)
=
(\mu^\square(\infty_e),\infty_e)
\]
in that universe.  This path does not induce an arithmetic equality
\[
e=\omega^{-1}
\]
in $G(\No[i])$.
\end{theorem}

\begin{proof}
The univalence axiom identifies the identity type between objects of a universe with the type of equivalences between them; see \cite{HoTT}.  Apply it to the projective inversion equivalence.  Arithmetic equality is a different identity type, namely equality of elements of the semiring.  By Theorem \ref{thm:strict-identification-collapse}, imposing $e=\omega^{-1}$ as an arithmetic relation produces the Boolean quotient and therefore cannot be obtained merely from the equivalence of pointed monads.
\end{proof}

\begin{remark}
Univalence therefore validates a genuine identity, but at the level at which the equivalence exists.  It does not erase invariants that belong to a stricter structure.  In particular,
\[
e^2=e,
\qquad
(\omega^{-1})^2=\omega^{-2}\neq\omega^{-1},
\]
so no multiplicative equivalence of pointed amplitude elements can identify them.
\end{remark}

\section{The absolute arithmetic curve and the split Abel--Jacobi fibre}
\label{sec:absolute-curve}

Three coefficient geometries must be distinguished.
\begin{enumerate}[label=\textup{(\roman*)}]
\item The original arithmetic site of Connes--Consani is
\[
\mathfrak A=(\widehat{\N^\times},\Nbar),
\qquad
\Nbar=(\N\cup\{\infty\},\min,+),
\]
a semiringed topos of characteristic one \cite{CCArithmeticSite}.
\item The absolute $\Fone$-site is
\[
\operatorname{Pic}_*
 =\bigl(\widehat{\N_0^\times},\cO_{\Fone}\bigr),
\qquad
\cO_{\Fone}=\Fone[T],
\]
where $\Fone[T]$ is a spherical algebra rather than an ordinary semiring
object \cite{CCAbsolute}.
\item Extension of scalars of the arithmetic site produces the scaling
site, whose structure sheaf is idempotent and encodes valuation profiles
\cite{CCScalingSite}.
\end{enumerate}
The point and stalk results below concern the absolute $\Fone$-site.  The
sheafwise split construction of Section~\ref{sec:split-scaling-site}
applies directly to the ordinary semiring objects $\Nbar$ and the scaling
sheaf; it does not silently identify spherical $\Fone$-algebras with
Boolean semirings.  The facts recalled in the next theorem are imported
from \cite{CCAbsolute,CCArithmeticSite}; they fix the precise interface
with split support.

\begin{theorem}[Connes--Consani point and stalk package]
\label{thm:cc-package}
The following hold.
\begin{enumerate}[label=\textup{(\roman*)}]
\item A point of $\widehat{\N_0^\times}$ is represented by a totally ordered
rank-one abelian group $H$ isomorphic to a subgroup of $\Q$, together with
the distinguished trivial point $\{0\}$.
\item The stalk of $\cO_{\Fone}$ at $H$ is
\[
A_H=\Fone[T^{H_+}].
\]
\item There is a geometric morphism
\[
\Theta:\Spec\Z\longrightarrow\widehat{\N_0^\times}
\]
whose map on points satisfies
\[
\Theta(\eta)=[\{0\}],
\qquad
\Theta(p)=[\Z[1/p]]
\]
for the generic point $\eta$ and every rational prime $p$.
\item There is a surjective map on point classes
\[
\kappa:\operatorname{Pt}(\widehat{\N_0^\times})
\longrightarrow \operatorname{Pic}(\Spec\Z)
\]
such that
\[
\kappa\circ\Theta=\theta,
\]
where $\theta$ is the arithmetic Abel--Jacobi map.  The two point objects
$\{0\}$ and $\Z$ are nonisomorphic, but both are sent by $\kappa$ to the
zero divisor class.
\end{enumerate}
\end{theorem}

Let
\[
S=G(\Z)=\Z\sqcup\{\tau\},
\qquad e=0_{\Z}\in S.
\]
Its prime spectrum is
\[
P_\tau=\{\tau\},
\qquad
P_0=\{\tau,e\},
\qquad
P_p=p\Z\cup\{\tau\}
\]
for rational primes $p$.  Recall that
\[
P_\tau\subsetneq P_0\subsetneq P_p.
\]

\begin{proposition}[Classical collapse of the split spectrum]
\label{prop:classical-collapse}
The assignment
\[
q:\Spec S\longrightarrow\Spec\Z,
\qquad
q(P_\tau)=q(P_0)=\eta,
\qquad
q(P_p)=p,
\]
is continuous.
\end{proposition}

\begin{proof}
For every nonzero $n\in\Z$, the usual basic open $D_{\Z}(n)$ has inverse
image
\[
q^{-1}(D_{\Z}(n))=D_S(n).
\]
Indeed, $P_\tau$ and $P_0$ both map to the generic point and neither
contains a nonzero integer, while $P_p$ belongs to $D_S(n)$ precisely when
$p\nmid n$.  The opens $D_{\Z}(n)$ with $n\neq0$, together with the empty
open, form a basis.  Hence $q$ is continuous.
\end{proof}

\begin{definition}[Split absolute point map]
\label{def:split-absolute-point-map}
Define a map on point classes
\[
\widetilde\Theta_{\mathrm{pt}}:\Spec S
\longrightarrow \operatorname{Pt}(\widehat{\N_0^\times})
\]
by
\[
\widetilde\Theta_{\mathrm{pt}}(P_\tau)=\{0\},
\qquad
\widetilde\Theta_{\mathrm{pt}}(P_0)=\Z,
\qquad
\widetilde\Theta_{\mathrm{pt}}(P_p)=\Z[1/p].
\]
\end{definition}

\begin{theorem}[Resolution of the zero Abel--Jacobi fibre]
\label{thm:abel-jacobi-resolution}
At the level of point classes one has a commutative diagram
\[
\begin{tikzcd}[column sep=large,row sep=large]
\Spec G(\Z)
  \arrow[r,"\widetilde\Theta_{\mathrm{pt}}"]
  \arrow[d,"q"']
& \operatorname{Pt}(\widehat{\N_0^\times})
  \arrow[d,"\kappa"]\\
\Spec\Z \arrow[r,"\theta"']
& \operatorname{Pic}(\Spec\Z).
\end{tikzcd}
\]
Moreover,
\[
\widetilde\Theta_{\mathrm{pt}}(P_\tau)
 \not\cong
\widetilde\Theta_{\mathrm{pt}}(P_0),
\]
whereas
\[
\kappa\bigl(\widetilde\Theta_{\mathrm{pt}}(P_\tau)\bigr)
=
\kappa\bigl(\widetilde\Theta_{\mathrm{pt}}(P_0)\bigr).
\]
Thus the split spectrum separates two absolute point objects lying over the
classical zero divisor class.
\end{theorem}

\begin{proof}
At $P_p$, both composites give the divisor class represented by
$\Z[1/p]$.  At $P_\tau$ and $P_0$, the lower composite is $\theta(\eta)$.
By Theorem \ref{thm:cc-package},
\[
\kappa(\{0\})=[\Z]=\kappa(\Z)=\theta(\eta).
\]
This proves commutativity.  The trivial group $\{0\}$ and the infinite
cyclic ordered group $\Z$ are nonisomorphic, proving the strict
separation.
\end{proof}

\begin{corollary}[The resolved stalk chain]
\label{cor:resolved-stalk-chain}
The point chain
\[
P_\tau\subset P_0\subset P_p
\]
is sent to the ordered-group chain
\[
\{0\}\hookrightarrow\Z\hookrightarrow\Z[1/p]
\]
and to the spherical-algebra chain
\[
\Fone
\longrightarrow
\Fone[T]
\longrightarrow
\Fone[T^{\Z[1/p]_+}].
\]
The last term is perfect for the $p$-th Frobenius and generally imperfect
for the other prime Frobenii.
\end{corollary}

\begin{proof}
The ordered-group arrows are the unique map from the trivial group and the
standard inclusion $\Z\subseteq\Z[1/p]$.  Applying
$H\mapsto\Fone[T^{H_+}]$ gives the displayed algebra maps.  Multiplication
by $p$ is bijective on $\Z[1/p]_+$, with inverse multiplication by $1/p$.
\end{proof}

\begin{remark}[Point-level status]
Theorem \ref{thm:abel-jacobi-resolution} is a theorem about point classes,
specialization, and the associated stalk diagram.  It does not by itself
construct a geometric morphism
$\Spec G(\Z)\to\widehat{\N_0^\times}$.  Such a morphism would require a
separate flatness and continuity analysis in the category of points of the
topos.  Keeping this distinction explicit prevents a pointwise dictionary
from being mistaken for a topos-level theorem.
\end{remark}


\section{The arithmetic Jacobian as an unreduced split-support monoid}
\label{sec:arithmetic-jacobian-split}

We now use the completed arithmetic Picard and Jacobian monoids of
Connes--Consani \cite{CCJacobian}.  To avoid notational ambiguity, let
$\Picf$ denote the monoid of isomorphism classes of torsion-free rank-one
abelian groups, let $\Picar$ denote the monoid of such groups equipped with
a possibly zero archimedean seminorm, and let
\[
\Jacar:=\Picar/\R_+^\times.
\]

\begin{theorem}[Connes--Consani configuration and Jacobian package]
\label{thm:cc-jacobian-package}
The following statements hold.
\begin{enumerate}[label=\textup{(\roman*)}]
\item The map
\[
\Theta_{\mathrm{conf}}:\mathcal P(\mathbb P_f)\longrightarrow\Picf,
\qquad
S\longmapsto\Z[S^{-1}],
\]
where $\mathbb P_f$ is the set of finite rational primes, is an isomorphism
from the union semilattice onto $\Idem(\Picf)$.
\item There is a canonical monoid isomorphism
\[
\Jacar\cong\Picf\times\{0,\infty\},
\]
where $0$ denotes a nondegenerate metric orbit, $\infty$ the zero-seminorm
fixed stratum, and the second factor is multiplied by
$0+0=0$ and $0+\infty=\infty+\infty=\infty$.
\item The extended Abel--Jacobi map sends
\[
\eta\longmapsto(\Z,0),
\qquad
p\longmapsto(\Z[1/p],0),
\qquad
\infty\longmapsto(\Z,\infty).
\]
\end{enumerate}
\end{theorem}

\begin{proof}
These are Theorems 2.15--2.16, Definition 4.1 and the structural
identification of Section 4, and the extended Abel--Jacobi formula of
\cite{CCJacobian}.  They are collected here with notation separating the
finite Picard monoid from its completed metric version.
\end{proof}

Identify $\{0,\infty\}$ with the multiplicative Boolean monoid by
\[
0\longmapsto1,
\qquad
\infty\longmapsto0.
\]
The rule for the archimedean coordinate then becomes Boolean
multiplication.

\begin{theorem}[Exact split-support form of the arithmetic Jacobian]
\label{thm:jacobian-boolean-product}
There is a canonical commutative-monoid isomorphism
\[
\boxed{\Jacar\cong\Picf\times\B.}
\]
Consequently,
\[
\boxed{
\Idem(\Jacar)
\cong
\mathcal P(\mathbb P_f)\times\B
\cong
\mathcal P(\mathbb P_f\sqcup\{\infty\})
}
\]
as join-semilattices.  Under the last isomorphism,
$(S,1)$ corresponds to $S$ and $(S,0)$ to $S\cup\{\infty\}$.
\end{theorem}

\begin{proof}
The first assertion is Theorem~\ref{thm:cc-jacobian-package}(ii) after the
stated identification of the two-element factors.  An element of a product
monoid is idempotent exactly when both components are idempotent.  Every
Boolean element is idempotent, and
Theorem~\ref{thm:cc-jacobian-package}(i) identifies the idempotents of
$\Picf$ with prime configurations.  Finally, Boolean multiplication in the
second coordinate becomes union of the distinguished symbol $\infty$ under
$0\leftrightarrow\{\infty\}$ and $1\leftrightarrow\varnothing$.
\end{proof}

\begin{proposition}[Boolean truncation of a degenerate Arakelov metric]
\label{prop:arakelov-support-truncation}
Let
\[
\sigma:\R_{\geq0}\longrightarrow\B,
\qquad
\sigma(t)=
\begin{cases}
0,&t=0,\\
1,&t>0.
\end{cases}
\]
Then $\sigma$ is a semiring homomorphism from
$(\R_{\geq0},+,\cdot)$ to $(\B,\vee,\wedge)$.  If $\ell$ is a real
one-dimensional vector space with a possibly zero seminorm $\|\cdot\|$,
the value
\[
\epsilon(\ell,\|\cdot\|):=\sigma(\|v\|),
\qquad 0\neq v\in\ell,
\]
is independent of $v$, invariant under positive rescaling of the
seminorm, and satisfies
\[
\epsilon(\ell\otimes m)=\epsilon(\ell)\epsilon(m).
\]
Under Theorem~\ref{thm:jacobian-boolean-product}, this is precisely the
projection $\Jacar\to\B$.
\end{proposition}

\begin{proof}
For nonnegative reals, $x+y=0$ exactly when $x=y=0$, and $xy>0$ exactly
when both factors are positive; hence $\sigma$ preserves addition and
multiplication.  A seminorm on a one-dimensional real vector space is
either identically zero or positive on every nonzero vector, so
$\epsilon$ is well defined.  Positive rescaling does not change its
support.  The tensor-product seminorm is zero exactly when at least one
factor seminorm is zero, proving multiplicativity.  Connes--Consani's two
archimedean orbits are precisely the nondegenerate and zero-seminorm
orbits, which are identified with $1$ and $0$ in
Theorem~\ref{thm:jacobian-boolean-product}.
\end{proof}

\begin{remark}
Thus split support is the exact Boolean shadow of the boundary degeneracy
of an Arakelov metric.  It does not retain the logarithmic metric,
Arakelov degree, or height data on the nondegenerate stratum.
\end{remark}

\begin{corollary}[The Abel--Jacobi support vertices]
\label{cor:abel-jacobi-support-vertices}
Under the idempotent identification of
Theorem~\ref{thm:jacobian-boolean-product}, the extended Abel--Jacobi map
sends
\[
\eta\longmapsto\varnothing,
\qquad
p\longmapsto\{p\},
\qquad
\infty\longmapsto\{\infty\}.
\]
For any finite set $J$ of places, its configuration submonoid is the
literal Boolean cube $\mathcal P(J)$.
\end{corollary}

\begin{proof}
The group $\Z$ is the empty localization, $\Z[1/p]$ is the singleton
finite-prime localization, and the zero seminorm sets precisely the
archimedean Boolean coordinate to $0$.
\end{proof}

Let $\pi_{\infty}(S)=(S\cap\mathbb P_f,\epsilon_S)$, where $\epsilon_S=0$ exactly when $\infty\in S$.  The two support coordinates and the ambient Jacobian are summarized by
\[
\begin{tikzcd}[column sep=large,row sep=large]
\mathcal P(\mathbb P_f\sqcup\{\infty\})
  \arrow[r,"\sim"] \arrow[d,"\pi_{\infty}"']
& \Idem(\Jacar) \arrow[d,hook]\\
\mathcal P(\mathbb P_f)\times\B
  \arrow[r,"\sim"']
& \Picf\times\B\cong\Jacar.
\end{tikzcd}
\]

The preceding product retains the group label even on the degenerate
stratum.  This is an unreduced split completion, whereas $G(R)$ has only
one new unsupported zero.

\begin{definition}[Unreduced and reduced multiplicative split]
\label{def:unreduced-multiplicative-split}
For a commutative monoid $M$, define
\[
\operatorname{USpl}(M):=M\times\B.
\]
Define the pointed split monoid
\[
M^\tau:=M\sqcup\{\tau\},
\]
where $\tau$ is absorbing.  Let
\[
c_M:\operatorname{USpl}(M)\to M^\tau,
\qquad
c_M(m,1)=m,
\quad
c_M(m,0)=\tau.
\]
\end{definition}

\begin{proposition}[Reduction of the degenerate stratum]
\label{prop:jacobian-reduction}
Let $\sim_0$ be the Rees congruence on $\operatorname{USpl}(M)$ whose only
non-singleton class is $M\times\{0\}$.  The map $c_M$ is a surjective
monoid homomorphism with kernel congruence $\sim_0$.  Hence
\[
\operatorname{USpl}(M)/{\sim_0}\cong M^\tau.
\]
For $M=\Picf$, the arithmetic Jacobian is the unreduced split completion
and its quotient by the whole zero-seminorm stratum is the reduced pointed
split monoid $\Picf^\tau$.
\end{proposition}

\begin{proof}
For $(m,\epsilon),(n,\delta)\in M\times\B$,
\[
c_M((m,\epsilon)(n,\delta))
=c_M(mn,\epsilon\delta).
\]
If $\epsilon\delta=1$, both factors have Boolean value $1$ and the image is
$mn$; if $\epsilon\delta=0$, at least one image is $\tau$ and both sides
are $\tau$.  Thus $c_M$ is multiplicative.  The remaining assertions are
immediate from its fibres.
\end{proof}

\begin{remark}[What is exact and what is analogy]
Theorem~\ref{thm:jacobian-boolean-product} is an exact restatement of the
Connes--Consani Jacobian theorem, not a dictionary.  The comparison with
$G(R)$ occurs after forgetting addition and then applying the explicit
quotient $c_{\Picf}$.  No claim is made that the full arithmetic Jacobian
is itself a split-zero semiring.
\end{remark}

\section{Rigidified denominator points and universal Frobenius}
\label{sec:rigidified-denominators}

Finite denominator refinement is invisible on isomorphism classes of
arithmetic-site points.  It becomes visible after retaining the embedding
of the arithmetic lattice.

\begin{definition}[Rigidified arithmetic-site point]
\label{def:rigidified-point}
A rigidified point is a pair $(H,\iota)$ consisting of an ordered rank-one
subgroup $H\subseteq\Q$ and an order-preserving embedding
\[
\iota:\Z\hookrightarrow H.
\]
A morphism
\[
f:(H,\iota_H)\longrightarrow(K,\iota_K)
\]
is an injective order-preserving homomorphism satisfying
$f\circ\iota_H=\iota_K$.

For $Q\geq1$, put
\[
H_Q=\frac1Q\Z,
\qquad
\mathfrak x_Q=(H_Q,\iota_Q),
\qquad
\iota_Q(n)=n.
\]
\end{definition}

\begin{theorem}[Divisibility classification]
\label{thm:rigidified-divisibility}
For positive integers $Q,M$,
\[
\operatorname{Hom}(\mathfrak x_Q,\mathfrak x_M)
=
\begin{cases}
\{\text{one morphism}\},&Q\mid M,\\
\varnothing,&Q\nmid M.
\end{cases}
\]
Consequently,
\[
\mathfrak x_Q\cong\mathfrak x_M
\quad\Longleftrightarrow\quad
Q=M.
\]
\end{theorem}

\begin{proof}
Let $f:\mathfrak x_Q\to\mathfrak x_M$ be a morphism.  Rigidification gives
$f(1)=1$.  For every $a\in\Z$,
\[
Qf(a/Q)=f(a)=a.
\]
Since $H_M\subseteq\Q$ is torsion free, $f(a/Q)=a/Q$.  Hence $f$ is forced
to be the inclusion $H_Q\hookrightarrow H_M$.  Such an inclusion exists
if and only if $1/Q\in(1/M)\Z$, equivalently $Q\mid M$.  The isomorphism
criterion follows by applying this condition in both directions.
\end{proof}

\begin{corollary}[The framed fibre over the arithmetic point $\Z$]
\label{cor:framed-fibre}
After forgetting the rigidification, every $\mathfrak x_Q$ represents the
same arithmetic-site point $[\Z]$.  The full divisibility category
$(\N^\times,\mid)$ is therefore recovered as a category of distinct
framings of one point.
\end{corollary}

\begin{proof}
Multiplication by $Q$ gives an order isomorphism
$H_Q\cong\Z$.  Theorem \ref{thm:rigidified-divisibility} shows that the
framings remain distinct and that their morphisms are exactly divisibility.
\end{proof}

Connes--Consani also realize the full adele class space by framed arithmetic
divisors $(L,\xi,\tau_\infty)$ \cite{CCJacobian}.  The denominator group
has a canonical such frame:
\[
\mathfrak D_Q
:=\left(\frac1Q\Z,\ \xi_Q(x)=Qx\in\widehat\Z,\
\tau_Q(x)=x\in\R\right).
\]

\begin{theorem}[The two denominator laws]
\label{thm:two-denominator-laws}
For positive integers $Q,M$:
\begin{enumerate}[label=\textup{(\roman*)}]
\item tensor product of framed arithmetic divisors gives
\[
\mathfrak D_Q\otimes\mathfrak D_M\cong\mathfrak D_{QM};
\]
\item least common refinement inside $\Q$ gives
\[
H_Q+H_M=H_{\operatorname{lcm}(Q,M)}.
\]
Thus multiplication of denominator indices is Picard multiplication,
whereas least common multiple is join in the embedded-root refinement
poset.
\end{enumerate}
\end{theorem}

\begin{proof}
Multiplication in $\Q$ induces
\[
H_Q\otimes_\Z H_M\xrightarrow{\sim}H_{QM},
\qquad
\frac aQ\otimes\frac bM\longmapsto\frac{ab}{QM}.
\]
Under this map the finite frames multiply as
$(Qa/Q)(Mb/M)=ab=QM(ab/QM)$, and the infinite frames multiply as ordinary
real numbers.  This proves (i).

Put $L=\operatorname{lcm}(Q,M)$.  Both $H_Q$ and $H_M$ lie in $H_L$.
Writing $d_Q=L/Q$ and $d_M=L/M$, one has
$\gcd(d_Q,d_M)=1$; choose $a,b\in\Z$ with $ad_Q+bd_M=1$.  Then
\[
\frac1L=\frac aQ+\frac bM,
\]
so the subgroup generated by $H_Q$ and $H_M$ contains $H_L$.  This proves
(ii).
\end{proof}

\begin{remark}[Why the LCM tower is not Picard multiplication]
The operation $Q\mapsto\Lambda_Q$ adjoins all roots of order at most $Q$
and is therefore a join/perfection operation.  It is not iteration of the
adelic tensor product, which would multiply denominator indices.  Keeping
these laws separate is necessary for the arithmetic-site interpretation of
the Beurling tower.
\end{remark}

For each $Q$, let
\[
A_Q=\Fone[T^{H_Q^+}]
=\Fone[T^{\frac1Q\N}].
\]
There is an abstract isomorphism
\[
\phi_Q:\Fone[T]\xrightarrow{\sim}A_Q,
\qquad
\phi_Q(T^n)=T^{n/Q}.
\]
If $Q\mid M$, inclusion of exponent monoids induces
$j_{Q,M}:A_Q\hookrightarrow A_M$.

\begin{theorem}[Denominator refinement is universal Frobenius]
\label{thm:denominator-frobenius}
If $M=mQ$, then
\[
\boxed{
\phi_M^{-1}\circ j_{Q,M}\circ\phi_Q
=\operatorname{Frob}_m
}
\]
on $\Fone[T]$, where
\[
\operatorname{Frob}_m(T^n)=T^{mn}.
\]
\end{theorem}

\begin{proof}
For $n\in\N$,
\[
T^n\xmapsto{\phi_Q}T^{n/Q}
\xmapsto{j_{Q,M}}T^{n/Q}
=T^{mn/M}
\xmapsto{\phi_M^{-1}}T^{mn}.
\]
This is precisely the $m$-th universal Frobenius.
\end{proof}

\begin{remark}
The finite denominator stages are isomorphic as unframed stalks.  Their
arithmetic content lies in the transition maps, which become Frobenius
endomorphisms after choosing the canonical coordinates $\phi_Q$.  This is
why denominator refinement should be regarded as motion in a framed fibre,
not as motion among ordinary point isomorphism classes.
\end{remark}

\section{LCM perfection, Steinitz points, and emergent places}
\label{sec:lcm-perfection}

\begin{definition}[Frobenius-perfect rigidified point]
\label{def:frobenius-perfect-point}
Let $\Z\subseteq H\subseteq\Q$ be a rigidified arithmetic-site point and
put
\[
A_H=\Fone[T^{H_+}].
\]
For a positive integer $m$, the point is called \emph{$m$-perfect} if
$\Frob_m:A_H\to A_H$ is an automorphism.  It is called \emph{globally
Frobenius-perfect} if it is $m$-perfect for every $m\geq1$.
\end{definition}

\begin{theorem}[Characterization and uniqueness of Frobenius perfection]
\label{thm:unique-frobenius-perfect-point}
Let $\Z\subseteq H\subseteq\Q$.
\begin{enumerate}[label=\textup{(\roman*)}]
\item $\Frob_m$ is an automorphism of $A_H$ if and only if $H$ is
$m$-divisible.
\item The smallest $p$-perfect rigidified point containing $\Z$ is
$\Z[1/p]$.
\item The unique globally Frobenius-perfect rigidified point is $H=\Q$.
\end{enumerate}
\end{theorem}

\begin{proof}
The Frobenius acts on monomials by
\[
T^h\longmapsto T^{mh}.
\]
Multiplication by $m$ is injective on the torsion-free group $H$.
Surjectivity on monomials is equivalent to the condition that for every
$h\in H_+$ there is $h'\in H_+$ with $mh'=h$, which is exactly
$m$-divisibility of $H$.  This proves (i).

The smallest subgroup of $\Q$ containing $\Z$ and closed under division
by $p$ is
\[
\bigcup_{r\geq0}p^{-r}\Z=\Z[1/p],
\]
which proves (ii).  If $H$ is globally Frobenius-perfect, then
$1/m\in H$ for every $m\geq1$, hence $\Q\subseteq H$.  Since
$H\subseteq\Q$, equality follows.  Conversely, $\Q$ is divisible by every
positive integer, proving (iii).
\end{proof}

\begin{warning}[Perfect is not perfectoid]
Theorem \ref{thm:unique-frobenius-perfect-point} concerns bijectivity of
universal Frobenius on an exponent monoid.  It does not provide a complete
nonarchimedean Tate ring, a tilt, or any of the analytic structures required
in perfectoid geometry.
\end{warning}

Let $Q_1,\dots,Q_r$ be positive integers and put
$L=\operatorname{lcm}(Q_1,\dots,Q_r)$.

\begin{theorem}[LCM as categorical join]
\label{thm:lcm-join}
Inside $\Q$ one has
\[
H_L=H_{Q_1}+\cdots+H_{Q_r}.
\]
Equivalently, $\mathfrak x_L$ is the join of the rigidified points
$\mathfrak x_{Q_i}$, and $A_L$ is the least cyclic-root stalk containing
all $A_{Q_i}$.
\end{theorem}

\begin{proof}
Because $Q_i\mid L$, one has $H_{Q_i}\subseteq H_L$.  Put $d_i=L/Q_i$.
The least-common-multiple property implies
$\gcd(d_1,\dots,d_r)=1$.  Hence integers $u_i$ exist with
$\sum_i u_id_i=1$.  Dividing by $L$ gives
\[
\frac1L=\sum_i\frac{u_i}{Q_i},
\]
so the subgroup generated by the $H_{Q_i}$ contains the generator $1/L$
of $H_L$.  This proves equality.  Passing to positive exponent monoids and
then to spherical monoid algebras proves the stalk statement.
\end{proof}

Define
\[
\Lambda_Q=\operatorname{lcm}(1,2,\dots,Q).
\]
Choose $Q_0\geq3$ and recursively set
\[
Q_{j+1}=\Lambda_{Q_j}.
\]

\begin{lemma}[Strict growth]
\label{lem:lcm-growth}
For $Q\geq3$,
\[
\Lambda_Q\geq Q(Q-1)>Q.
\]
\end{lemma}

\begin{proof}
The coprime integers $Q$ and $Q-1$ both occur in the defining least common
multiple, so
\[
\Lambda_Q\geq\operatorname{lcm}(Q,Q-1)=Q(Q-1).
\]
\end{proof}

\begin{theorem}[Global Frobenius perfection]
\label{thm:global-perfection}
The filtered colimits of the LCM tower are
\[
\varinjlim_jH_{Q_j}=\Q,
\qquad
\varinjlim_jA_{Q_j}=\Fone[T^{\Q_+}].
\]
Every universal Frobenius
\[
\operatorname{Frob}_m:\Fone[T^{\Q_+}]
\longrightarrow\Fone[T^{\Q_+}]
\]
is an automorphism.
\end{theorem}

\begin{proof}
The union of the groups $H_{Q_j}$ is contained in $\Q$.  Conversely, take
$a/b\in\Q$.  By Lemma \ref{lem:lcm-growth}, choose $j$ with $Q_j\geq b$.
Then $b\mid Q_{j+1}$, hence $a/b\in H_{Q_{j+1}}$.  This proves the first
identity.  The second follows by taking monoid algebras of the increasing
positive cones.  Multiplication by $m$ is a bijection of $\Q_+$ with
inverse $q\mapsto q/m$, so the induced Frobenius is invertible.
\end{proof}

A supernatural or Steinitz number is a formal product
\[
N=\prod_p p^{n_p},
\qquad
n_p\in\N\cup\{\infty\}.
\]
It determines the ordered subgroup
\[
H_N=
\{x\in\Q:v_p(x)\geq-n_p\text{ for every }p\}.
\]
This is the modern point-theoretic form of Steinitz's $G$-numbers
\cite{Steinitz}.

\begin{corollary}[The maximal supernatural endpoint]
\label{cor:maximal-supernatural}
For
\[
\Omega=\prod_p p^\infty,
\]
one has $H_\Omega=\Q$.  Thus the LCM tower converges to the maximal
supernatural point $x_\Omega$ of the arithmetic site, whose stalk is the
globally Frobenius-perfect algebra $\Fone[T^{\Q_+}]$.
\end{corollary}

\begin{theorem}[Places emerge only at perfection]
\label{thm:places-emerge}
Let $\Sigma_H$ be the Connes--Consani set of places at which $H$ embeds
densely.  Then
\[
\Sigma_{H_{Q_j}}=\varnothing
\quad\text{for every finite stage }j,
\]
whereas
\[
\Sigma_{\Q}=\{\infty\}\cup\{p:p\text{ prime}\}.
\]
In particular,
\[
\bigcup_j\Sigma_{H_{Q_j}}
\neq
\Sigma_{\varinjlim_jH_{Q_j}}.
\]
\end{theorem}

\begin{proof}
Each $H_{Q_j}$ is order-isomorphic to $\Z$, and Connes--Consani prove that
$\Sigma_H$ is empty for $H\cong\Z$.  The group $\Q$ is dense in $\R$ and
is $p$-divisible for every prime $p$, so the classification of places in
\cite{CCAbsolute} gives the second identity.
\end{proof}

\begin{remark}
Theorem \ref{thm:places-emerge} is not a formal paradox.  The functor
$H\mapsto\Sigma_H$ detects density and divisibility properties that need
not occur at any finite cyclic-root stage.  Coherent Frobenius perfection
creates local directions that are absent at every finite level.
\end{remark}

\section{Cyclotomic carry, adelic history, and the perfected denominator anomaly}
\label{sec:perfected-anomaly}

The globally Frobenius-perfect point carries two distinct kinds of data.
The first is the relative denominator quotient produced by the fixed
rigidification $\Z\subseteq\Q$.  The second is the inverse-limit character
data retained by the transition maps of the denominator tower.  They are
related by two canonical nonsplit extensions.  This section computes both
extensions, identifies the resulting profinite symmetry, and compares the
answer with the rooted arithmetic divisors of Connes--Consani.

The word \emph{perfect} in this section always means Frobenius perfection
of an arithmetic-site stalk.  It does not mean that
$\Fone[T^{\Q_+}]$ is a perfectoid Tate ring or that $x_\Omega$ is a
perfectoid space in the sense of Scholze \cite{ScholzePerfectoid}.

\subsection{Finite denominator phases and their extension classes}

\begin{definition}[Relative phase quotient]
\label{def:finite-phase-quotient}
For $Q\geq1$, define
\[
\Phi_Q:=H_Q/\Z=\frac{Q^{-1}\Z}{\Z}.
\]
The map
\[
\eta_Q:\Z/Q\Z\xrightarrow{\sim}\Phi_Q,
\qquad
[a]_Q\longmapsto \frac aQ+\Z,
\]
is a canonical isomorphism.  We write
\[
\mathcal E_Q:
0\longrightarrow\Z\longrightarrow H_Q
\longrightarrow\Phi_Q\longrightarrow0
\]
for the resulting short exact sequence.
\end{definition}

\begin{theorem}[Finite carry class]
\label{thm:finite-carry-class}
Under the standard identification
\[
\Ext^1_{\Z}(\Z/Q\Z,\Z)\cong\Z/Q\Z,
\]
the extension class of $\mathcal E_Q$ is
\[
[\mathcal E_Q]=1\pmod Q.
\]
In particular, $\mathcal E_Q$ is nonsplit for $Q>1$.

If $Q\mid M$, the canonical inclusion
$\iota_{Q,M}:\Phi_Q\hookrightarrow\Phi_M$ satisfies
\[
\iota_{Q,M}^{*}[\mathcal E_M]=[\mathcal E_Q].
\]
Under the displayed Ext-identifications, the induced map is reduction
$\Z/M\Z\to\Z/Q\Z$.
\end{theorem}

\begin{proof}
The projective resolution
\[
0\longrightarrow\Z\xrightarrow{\,Q\,}\Z
\longrightarrow\Z/Q\Z\longrightarrow0
\]
gives
\[
\Ext^1_{\Z}(\Z/Q\Z,\Z)
\cong\operatorname{coker}(Q:\Z\to\Z)
\cong\Z/Q\Z.
\]
Multiplication by $Q$ gives an isomorphism $H_Q\xrightarrow{\sim}\Z$.
Under this isomorphism, the inclusion $\Z\hookrightarrow H_Q$ becomes
multiplication by $Q$.  Thus $\mathcal E_Q$ is the pushout of the displayed
resolution along the identity of its kernel, and therefore represents the
class $1\pmod Q$.

Suppose $M=mQ$.  The inclusion $H_Q\hookrightarrow H_M$ induces
\[
\iota_{Q,M}:\Phi_Q\longrightarrow\Phi_M,
\qquad
\frac aQ+\Z\longmapsto\frac aQ+\Z.
\]
The pullback of $\mathcal E_M$ along this map consists exactly of those
$h\in H_M$ whose class modulo $\Z$ lies in the image of $\Phi_Q$; this
subgroup is $H_Q$.  Hence the pullback extension is $\mathcal E_Q$.
The corresponding contravariant map on Ext sends $1\pmod M$ to
$1\pmod Q$, and is therefore the ordinary reduction map.
\end{proof}

\begin{corollary}[Finite sign and phase shadows]
\label{cor:finite-phase-shadows}
The case $Q=2$ is the binary sign extension, the case $Q=3$ is the
ternary phase extension, and general $Q$ gives a canonical $Q$-phase
extension.  None of these extensions admits a homomorphic choice of a
representative phase.
\end{corollary}

\begin{proof}
Only the terminology is new.  Nonsplitting is Theorem
\ref{thm:finite-carry-class}.
\end{proof}

\begin{definition}[Finite carry cocycle]
\label{def:finite-carry-cocycle}
For $[a]_Q\in\Z/Q\Z$, let $0\leq a<Q$ be its standard representative and
put
\[
s_Q([a]_Q)=\frac aQ\in H_Q.
\]
Define
\[
c_Q([a]_Q,[b]_Q)
:=s_Q([a]_Q)+s_Q([b]_Q)-s_Q([a+b]_Q)
=\left\lfloor\frac{a+b}{Q}\right\rfloor.
\]
\end{definition}

\begin{proposition}[Twisted reconstruction of the denominator group]
\label{prop:finite-twisted-reconstruction}
The function $c_Q$ is a normalized symmetric $2$-cocycle with values in
$\Z$.  On the set $\Z\times\Z/Q\Z$, define
\[
(n,\bar a)+(m,\bar b)
=
\bigl(n+m+c_Q(\bar a,\bar b),\bar a+\bar b\bigr).
\]
Then
\[
\Z\times_{c_Q}\Z/Q\Z\xrightarrow{\sim}H_Q,
\qquad
(n,[a]_Q)\longmapsto n+\frac aQ,
\]
is an isomorphism of abelian groups.  The cocycle $c_Q$ represents
$[\mathcal E_Q]$.
\end{proposition}

\begin{proof}
Normalization and symmetry are immediate.  Associativity of the displayed
operation is equivalent to
\[
c_Q(x,y)+c_Q(x+y,z)=c_Q(y,z)+c_Q(x,y+z),
\]
which follows by expanding each side using the section $s_Q$ and
associativity in $H_Q$.  Euclidean division gives a unique expression
$h=n+a/Q$ with $0\leq a<Q$, so the displayed map is bijective.  Its
compatibility with addition is exactly the definition of $c_Q$.
\end{proof}

\subsection{The global carry class}

The filtered union of the finite phase groups is
\[
\Phi_\Omega:=\varinjlim_Q\Phi_Q=\Q/\Z.
\]
The rigidification $\Z\subseteq\Q$ gives the global extension
\[
\mathcal E_\Omega:
0\longrightarrow\Z\longrightarrow\Q
\longrightarrow\Q/\Z\longrightarrow0.
\]

\begin{theorem}[The global carry class is the profinite unit]
\label{thm:global-carry-class}
There is a canonical isomorphism
\[
\Ext^1_{\Z}(\Q/\Z,\Z)
\xrightarrow{\sim}
\varprojlim_Q\Ext^1_{\Z}(\Phi_Q,\Z)
\xrightarrow{\sim}
\varprojlim_Q\Z/Q\Z
=\widehat\Z.
\]
Under this isomorphism,
\[
\boxed{
[\mathcal E_\Omega]
=(1\bmod Q)_Q
=1_{\widehat\Z}.}
\]
Consequently, $\mathcal E_\Omega$ is nonsplit, and every finite carry
class is its pullback along $\Phi_Q\hookrightarrow\Q/\Z$.
\end{theorem}

\begin{proof}
The system $\{\Phi_Q\}$ is filtered by divisibility and has colimit
$\Q/\Z$.  The standard derived-limit exact sequence for Ext \cite{Weibel} gives
\[
0\longrightarrow
\varprojlim
olimits^{1}_Q\Hom(\Phi_Q,\Z)
\longrightarrow
\Ext^1_{\Z}(\varinjlim_Q\Phi_Q,\Z)
\longrightarrow
\varprojlim_Q\Ext^1_{\Z}(\Phi_Q,\Z)
\longrightarrow0.
\]
Every homomorphism from the finite group $\Phi_Q$ to the torsion-free group
$\Z$ is zero.  Hence the $\varprojlim^1$ term vanishes and the first arrow
displayed in the theorem is an isomorphism.  The second isomorphism is
Theorem \ref{thm:finite-carry-class}; its transition maps are reductions.
The pullback of $\mathcal E_\Omega$ to $\Phi_Q$ is $\mathcal E_Q$, so its
image is the compatible family $(1\bmod Q)_Q$, which is the multiplicative
unit of $\widehat\Z$.
\end{proof}

\begin{definition}[Global carry cocycle]
\label{def:global-carry-cocycle}
Let
\[
s:\Q/\Z\longrightarrow\Q\cap[0,1)
\]
be the standard set-theoretic section.  Define
\[
c_\Omega(x,y)=s(x)+s(y)-s(x+y)\in\Z.
\]
\end{definition}

\begin{proposition}[Explicit global history law]
\label{prop:global-carry-cocycle}
The function $c_\Omega$ is a normalized symmetric $2$-cocycle and
\[
\Q\cong\Z\times_{c_\Omega}(\Q/\Z)
\]
by the same twisted addition as in Proposition
\ref{prop:finite-twisted-reconstruction}.  Its restriction to $\Phi_Q$ is
$c_Q$.  Thus $c_\Omega$ simultaneously records every finite carry law.
\end{proposition}

\begin{proof}
The proof of Proposition \ref{prop:finite-twisted-reconstruction} applies
verbatim.  If $x=[a/Q]$ and $y=[b/Q]$ with $0\leq a,b<Q$, then
\[
c_\Omega(x,y)=\left\lfloor\frac{a+b}{Q}\right\rfloor=c_Q([a]_Q,[b]_Q).
\]
\end{proof}

\begin{corollary}[Classifying-space form]
\label{cor:carry-gerbe}
The extension $\mathcal E_\Omega$ induces a homotopy fibre sequence of
classifying spaces
\[
B\Z\longrightarrow B\Q\longrightarrow B(\Q/\Z).
\]
Equivalently, it determines a principal $B\Z$-bundle over $B(\Q/\Z)$,
classified by the central-extension class represented by $c_\Omega$ in
$H^2(B(\Q/\Z);\Z)$.  Its restriction to $B(\Z/Q\Z)$ is the generator
represented by $c_Q$.
\end{corollary}

\begin{proof}
A central extension of discrete groups gives a fibration of classifying
spaces with fibre the classifying space of the kernel.  The normalized
cocycle is the standard group-cohomological representative of that
extension.
\end{proof}

\subsection{Adelic phase history}

For a discrete abelian group $H$, put
\[
H^\vee:=\Hom(H,\Q/\Z).
\]
Connes--Consani identify the character dual of $\Q$ with the additive group
of finite adeles \cite{CCJacobian}:
\[
\psi:\mathbb A_f\xrightarrow{\sim}\Hom(\Q,\Q/\Z),
\qquad
\psi(a)(q)=\pi(qa),
\]
where $\pi:\mathbb A_f\to\mathbb A_f/\widehat\Z\cong\Q/\Z$ is the
canonical quotient.

\begin{theorem}[The inverse limit of finite phase histories is adelic]
\label{thm:adelic-history-limit}
For every $Q$, evaluation at $1/Q$ gives an isomorphism
\[
\operatorname{ev}_Q:H_Q^\vee\xrightarrow{\sim}\Q/\Z.
\]
If $Q\mid M=mQ$, the restriction map
\[
H_M^\vee\longrightarrow H_Q^\vee
\]
corresponds under these evaluations to multiplication by $m$ on
$\Q/\Z$.  Consequently,
\[
\boxed{
\varprojlim_Q H_Q^\vee
\cong
\Hom\!\left(\varinjlim_QH_Q,\Q/\Z\right)
\cong
\Hom(\Q,\Q/\Z)
\cong
\mathbb A_f.}
\]
A finite adele is therefore exactly a coherent family of denominator-phase
histories.
\end{theorem}

\begin{proof}
The group $H_Q$ is freely generated by $1/Q$, so evaluation at that
generator is an isomorphism.  If $M=mQ$ and $\varphi\in H_M^\vee$, then
\[
\varphi(1/Q)=\varphi(m/M)=m\varphi(1/M),
\]
which proves the transition formula.  Since $\Hom(-,\Q/\Z)$ sends
filtered colimits to inverse limits,
\[
\varprojlim_QH_Q^\vee
\cong\Hom(\varinjlim_QH_Q,\Q/\Z).
\]
Theorem \ref{thm:global-perfection} identifies the colimit with $\Q$, and
the final isomorphism is the Connes--Consani adelic duality theorem.
\end{proof}

\begin{theorem}[The adelic history extension]
\label{thm:adelic-history-extension}
Applying $\Hom(-,\Q/\Z)$ to $\mathcal E_\Omega$ gives the canonical exact
sequence
\[
\boxed{
0\longrightarrow\widehat\Z
\longrightarrow\mathbb A_f
\xrightarrow{\ \pi\ }\Q/\Z
\longrightarrow0.}
\]
Here
\[
\widehat\Z=\End(\Q/\Z)
\]
is the subgroup of histories whose restriction to the rigidifying copy of
$\Z$ is zero, and $\pi(a)=\psi(a)(1)$.  This extension is nonsplit, even as
an abstract extension of abelian groups.
\end{theorem}

\begin{proof}
The group $\Q/\Z$ is divisible and therefore injective in the category of
abelian groups.  Applying $\Hom(-,\Q/\Z)$ to
\[
0\to\Z\to\Q\to\Q/\Z\to0
\]
therefore gives
\[
0\to\Hom(\Q/\Z,\Q/\Z)
\to\Hom(\Q,\Q/\Z)
\to\Hom(\Z,\Q/\Z)
\to0.
\]
The three nonzero terms identify respectively with $\widehat\Z$,
$\mathbb A_f$, and $\Q/\Z$.  Under the Connes--Consani isomorphism, the
last arrow is evaluation at $1$ and is the canonical quotient $\pi$.

The additive group of $\mathbb A_f$ is torsion free.  A splitting would
embed the nonzero torsion group $\Q/\Z$ into $\mathbb A_f$, which is
impossible.
\end{proof}

\begin{figure}[t]
\centering
\begin{tikzcd}[column sep=large,row sep=large]
0 \arrow[r] & \Z \arrow[r] \arrow[d,equal]
& H_Q \arrow[r] \arrow[d,hook]
& \Phi_Q \arrow[r] \arrow[d,hook] & 0 \\
0 \arrow[r] & \Z \arrow[r]
& \Q \arrow[r]
& \Q/\Z \arrow[r] & 0
\end{tikzcd}
\qquad
\begin{tikzcd}[column sep=large,row sep=large]
0 \arrow[r] & \widehat\Z \arrow[r]
& \mathbb A_f \arrow[r,"\pi"]
& \Q/\Z \arrow[r] & 0
\end{tikzcd}
\caption{The finite carry extension embeds in the global carry extension;
character duality converts the global extension into the nonsplit adelic
history extension.}
\label{fig:carry-history-extensions}
\end{figure}

\begin{theorem}[Torsion transmutation at perfection]
\label{thm:torsion-transmutation}
For every finite denominator stage,
\[
\Tor(H_Q^\vee)=H_Q^\vee\cong\Q/\Z.
\]
At the globally perfect point,
\[
\Q^\vee\cong\mathbb A_f,
\qquad
\Tor(\Q^\vee)=0.
\]
Consequently,
\[
\boxed{
\Tor\!\left(\varprojlim_QH_Q^\vee\right)=0,
\qquad
\varprojlim_Q\Tor(H_Q^\vee)\cong\mathbb A_f.}
\]
Thus the torsion functor does not commute with the inverse limit of the
denominator tower: finite cyclotomic phase data is converted into
nontorsion adelic history.
\end{theorem}

\begin{proof}
Each $H_Q$ is infinite cyclic, hence $H_Q^\vee\cong\Q/\Z$ and is entirely
torsion.  Theorem \ref{thm:adelic-history-limit} identifies the inverse
limit with $\mathbb A_f$, whose additive group is torsion free.
\end{proof}

\subsection{Relative phases and rooted phases}

For a set $S$ of rational primes, put
\[
H_S:=\Z[S^{-1}]
=\Z[p^{-1}:p\in S]\subseteq\Q.
\]
Let $C_{p^\infty}=\Q_p/\Z_p$ denote the Pruefer $p$-group.

\begin{theorem}[Phase complementarity for saturated denominator points]
\label{thm:phase-complementarity}
The relative denominator quotient and the torsion of the Connes--Consani
character dual are
\[
H_S/\Z\cong\bigoplus_{p\in S}C_{p^\infty},
\qquad
\Tor(H_S^\vee)
\cong\bigoplus_{p\notin S}C_{p^\infty}.
\]
Hence the canonical primary decomposition of $\Q/\Z$ gives
\[
\boxed{
H_S/\Z\ \oplus\ \Tor(H_S^\vee)
\cong\Q/\Z.}
\]
At $S=\varnothing$, all cyclotomic phase is rooted torsion and no relative
phase is present.  At $S$ equal to the set of all primes, all cyclotomic
phase is relative denominator phase and the rooted torsion vanishes.
\end{theorem}

\begin{proof}
The $p$-primary subgroup of $H_S/\Z$ is zero when $p\notin S$ and is the
union of the cyclic groups $p^{-n}\Z/\Z$ when $p\in S$, hence is
$C_{p^\infty}$.  This proves the first formula.  The singular primes of
$H_S$ in the sense of Connes--Consani are exactly the primes in $S$.
Their character-dual theorem identifies the torsion subgroup with the
roots of unity whose order is prime to every singular prime, namely
$\bigoplus_{p\notin S}C_{p^\infty}$.  The final formula is the canonical
primary decomposition
$\Q/\Z=\bigoplus_pC_{p^\infty}$.
\end{proof}

\begin{corollary}[Root collapse at global perfection]
\label{cor:root-collapse-global-perfection}
For each finite $Q$, the singular-prime set of $H_Q$ is empty and
\[
\Tor(H_Q^\vee)\cong\Q/\Z.
\]
For $H_\Omega=\Q$, every prime is singular and
\[
\Tor(H_\Omega^\vee)=0.
\]
Thus a Connes--Consani root map
\[
\rho:\Q/\Z\longrightarrow H_\Omega^\vee
\]
whose image is required to lie in the torsion subgroup is necessarily
zero.  The finite root data does not survive as torsion at the endpoint;
it survives through the adelic history group of Theorem
\ref{thm:adelic-history-limit}.
\end{corollary}

\begin{proof}
Finite exponents never make $H_Q$ divisible by a prime, whereas $\Q$ is
$p$-divisible for every $p$.  Apply the Connes--Consani torsion theorem and
Theorem \ref{thm:torsion-transmutation}.
\end{proof}

\begin{remark}[The exact location of the history]
The static perfect point does not contain a residual group of torsion roots.
The history is the inverse-limit object $\mathbb A_f$ and, more precisely,
the nonsplit extension of Theorem \ref{thm:adelic-history-extension}.
This distinction prevents the relative quotient $\Q/\Z$ from being
mistaken for the torsion subgroup of $\Q^\vee$.
\end{remark}

\subsection{Cyclotomic symmetry and the Bost--Connes interface}

\begin{theorem}[Cyclotomic symmetry of the anomaly]
\label{thm:cyclotomic-symmetry-anomaly}
There are canonical isomorphisms
\[
\Q/\Z\cong\bigoplus_pC_{p^\infty},
\qquad
\End(\Q/\Z)\cong\prod_p\Z_p=\widehat\Z,
\]
and therefore
\[
\boxed{
\Aut(\Q/\Z)\cong\widehat\Z^{\,\times}.}
\]
The exponential map $x\mapsto e^{2\pi i x}$ identifies $\Q/\Z$ with the
group $\mu_\infty$ of all complex roots of unity.  Multiplication by
$u\in\widehat\Z^{\,\times}$ acts equivariantly on the adelic history
extension:
\[
\begin{tikzcd}[column sep=large]
0 \arrow[r] & \widehat\Z \arrow[r] \arrow[d,"\times u"]
& \mathbb A_f \arrow[r,"\pi"] \arrow[d,"\times u"]
& \Q/\Z \arrow[r] \arrow[d,"\times u"] & 0 \\
0 \arrow[r] & \widehat\Z \arrow[r]
& \mathbb A_f \arrow[r,"\pi"]
& \Q/\Z \arrow[r] & 0.
\end{tikzcd}
\]
\end{theorem}

\begin{proof}
The primary decomposition of a torsion divisible group gives the first
isomorphism.  The endomorphism ring of $C_{p^\infty}$ is $\Z_p$, and
homomorphisms between distinct primary components vanish.  Hence
\[
\End(\Q/\Z)\cong\prod_p\Z_p.
\]
Its invertible elements are exactly $\prod_p\Z_p^\times$.  The
exponential identification is immediate.  Finally, $u$ is a unit of the
maximal compact subring $\widehat\Z\subset\mathbb A_f$, and
$\pi(ua)=u\pi(a)$ under the identification of $\widehat\Z$ with
$\End(\Q/\Z)$.
\end{proof}

\begin{corollary}[Bost--Connes symmetry group]
\label{cor:bc-symmetry-interface}
The automorphism group of the perfected phase quotient is the arithmetic
symmetry group of the Bost--Connes system:
\[
\Aut(\Q/\Z)
\cong\widehat\Z^{\,\times}
\cong\Gal(\Q^{\mathrm{ab}}/\Q).
\]
Under the last isomorphism, the action on $\mu_\infty$ is the usual
cyclotomic Galois action.  This identifies the symmetry group of the
phase-history extension with the symmetry acting on the low-temperature
and ground-state sectors of the Bost--Connes system
\cite{BostConnes,CCBCCyclotomy,CMR}.  It does not assert that the split
semiring itself carries the Bost--Connes $C^*$-dynamics.
\end{corollary}

\begin{proof}
The first isomorphism is Theorem
\ref{thm:cyclotomic-symmetry-anomaly}; the second is the Kronecker--Weber
class-field-theory identification used in the Bost--Connes system.
\end{proof}

\subsection{Steinitz history and Möbius recovery}

\begin{definition}[Steinitz divisor history]
\label{def:steinitz-history}
For a supernatural number $N$, let
\[
\operatorname{Div}_{\mathrm{fin}}(N)
=\{n\in\N^\times:n\mid N\}.
\]
For $\Omega=\prod_pp^\infty$,
\[
\operatorname{Div}_{\mathrm{fin}}(\Omega)=\N^\times.
\]
We call the divisibility category
$(\N^\times,\mid)$ the finite denominator history of $x_\Omega$.
\end{definition}

\begin{proposition}[Cofinality and reconstruction of the anomaly class]
\label{prop:history-cofinality}
Every LCM tower $Q_{j+1}=\Lambda_{Q_j}$ with $Q_0\geq3$ is cofinal in
$(\N^\times,\mid)$.  Consequently,
\[
\widehat\Z
\cong\varprojlim_j\Z/Q_j\Z,
\]
and the class $[\mathcal E_\Omega]$ is recovered from the finite history as
\[
[\mathcal E_\Omega]=(1\bmod Q_j)_j.
\]
\end{proposition}

\begin{proof}
Given $n\geq1$, choose $j$ with $Q_j\geq n$.  Then $n\mid\Lambda_{Q_j}
=Q_{j+1}$.  Thus the tower is cofinal.  Inverse limits are unchanged by
restriction to a cofinal subsystem, and Theorem
\ref{thm:global-carry-class} gives the class formula.
\end{proof}

\begin{theorem}[Incidence inversion on denominator history]
\label{thm:mobius-history}
Let $A$ be an abelian group and let $f:\N^\times\to A$.  Define
\[
g(n)=\sum_{d\mid n}f(d).
\]
Then
\[
\boxed{
f(n)=\sum_{d\mid n}\mu(d)\,g(n/d),}
\]
where $\mu$ is the classical Möbius function.  Equivalently, the zeta
function of the locally finite divisibility poset has convolution inverse
\[
\mu_{\mid}(d,n)=\mu(n/d).
\]
Thus cumulative data on all predecessor denominator stages can be
recovered exactly from the finite Steinitz history.
\end{theorem}

\begin{proof}
Substitute the definition of $g$:
\[
\sum_{d\mid n}\mu(d)g(n/d)
=
\sum_{d\mid n}\mu(d)
\sum_{e\mid n/d}f(e).
\]
The coefficient of $f(e)$ is
\[
\sum_{d\mid n/e}\mu(d),
\]
which equals $1$ if $e=n$ and $0$ otherwise.
\end{proof}

\begin{remark}[Primewise self-similarity]
The exponent-vector identification
\[
\N^\times\cong\bigoplus_p\N
\]
turns denominator history into a restricted product of one-prime chains.
Dually,
\[
\Q/\Z\cong\bigoplus_pC_{p^\infty},
\qquad
\widehat\Z\cong\prod_p\Z_p.
\]
This primewise factorization is the precise self-similar structure behind
the informal language of an infinite or fractal phase history.
\end{remark}

\subsection{Combination with split support and place emergence}

Let $\chi_\C:G(\C)\to\B$ be the Boolean support character.  Under the
exponential identification $\Q/\Z\cong\mu_\infty\subset\C^\times$,
\[
\chi_\C(\zeta)=1
\quad(\zeta\in\mu_\infty),
\qquad
\chi_\C(e)=1,
\qquad
\chi_\C(\tau)=0.
\]
Thus Boolean support remembers neither cyclotomic phase nor the distinction
between nonzero phase and supported zero.  The missing information is not
another scalar in $G(\C)$; it is the relative extension and adelic
inverse-limit data computed above.

\begin{definition}[Perfected denominator anomaly]
\label{def:perfected-denominator-anomaly}
For a cofinal LCM tower $H_{Q_j}\to\Q$, define its perfected denominator
anomaly to be the triple
\[
\mathfrak A_\Omega
:=
\left(
[\mathcal E_\Omega],\
\mathcal E_\Omega^\vee,\
\operatorname{Emerg}_\Sigma
\right),
\]
where
\[
[\mathcal E_\Omega]\in\Ext^1_{\Z}(\Q/\Z,\Z),
\]
\[
\mathcal E_\Omega^\vee:
0\to\widehat\Z\to\mathbb A_f\to\Q/\Z\to0,
\]
and
\[
\operatorname{Emerg}_\Sigma
:=
\Sigma_{\Q}\setminus\bigcup_j\Sigma_{H_{Q_j}}.
\]
\end{definition}

\begin{theorem}[Exact computation of the perfected anomaly]
\label{thm:perfected-anomaly-computation}
The perfected denominator anomaly is
\[
\boxed{
\mathfrak A_\Omega
=
\left(
1_{\widehat\Z},\
0\to\widehat\Z\to\mathbb A_f\to\Q/\Z\to0,\
\{\infty\}\cup\{p:p\text{ prime}\}
\right).}
\]
Both extensions in this formula are nonsplit.  The first component retains
all finite carry classes; the second retains all coherent phase histories;
the third is the complete set of local directions created at Frobenius
perfection.
\end{theorem}

\begin{proof}
The first component is Theorem \ref{thm:global-carry-class}; the second is
Theorem \ref{thm:adelic-history-extension}; and the third is Theorem
\ref{thm:places-emerge}.
\end{proof}

\begin{remark}[Scope of the anomaly terminology]
The term \emph{anomaly} denotes the explicit obstruction classes and
failure-of-continuity statements in Definition
\ref{def:perfected-denominator-anomaly}.  No identification is made with a
particular quantum field theoretic anomaly.  Likewise, the hidden subgroup
$\widehat\Z$ may be viewed as a support-invisible history kernel, but no
identification with the ghost characteristics of nonlocal differential
symmetry calculus is used in any proof.
\end{remark}


\section{Finite quadratic phases, gerbe anomalies, and theta quantization}
\label{sec:gerbe-quantization}

The perfected carry extension determines finite phase groups
\[
\Phi_n=n^{-1}\Z/\Z\cong\Z/n\Z.
\]
A gerbe anomaly requires one further datum: a homogeneous quadratic
refinement.  The rigidified generator $1/n$ supplies such a refinement at
every level.

\begin{definition}[Rigidified cyclic quadratic form]
\label{def:cyclic-quadratic-form}
For $\bar a\in\Phi_n$, represented by $0\leq a<n$, define
\[
q_n(\bar a)=
\begin{cases}
\dfrac{a^2}{2n}\pmod\Z,&n\ \text{even},\\[6pt]
\dfrac{n+1}{2}\dfrac{a^2}{n}\pmod\Z,&n\ \text{odd}.
\end{cases}
\]
\end{definition}

\begin{theorem}[Nondegenerate quadratic refinement]
\label{thm:cyclic-quadratic-refinement}
The function $q_n:\Phi_n\to\Q/\Z$ is well defined and homogeneous:
$q_n(r\bar a)=r^2q_n(\bar a)$ for every integer $r$.  Its polarization is
the perfect pairing
\[
b_n(\bar a,\bar b)
=q_n(\bar a+\bar b)-q_n(\bar a)-q_n(\bar b)
=\frac{ab}{n}\pmod\Z.
\]
\end{theorem}

\begin{proof}
If $n$ is even, replacing $a$ by $a+n$ changes $a^2/(2n)$ by
$a+n/2\in\Z$.  If $n$ is odd, the change is
\[
\frac{n+1}{2n}\bigl((a+n)^2-a^2\bigr)
=(n+1)a+\frac{n(n+1)}2\in\Z.
\]
Thus $q_n$ is well defined, and homogeneity follows from the defining
quadratic formula.  For even $n$ the polarization is $ab/n$.  For odd
$n$ it is $(n+1)ab/n$, which is congruent to $ab/n$ modulo $\Z$.  If
$b_n(a,b)=0$ for every $b$, then $a/n\in\Z$, so $a=0$ in $\Phi_n$.
\end{proof}

\begin{theorem}[Finite gerbe anomaly associated with a denominator level]
\label{thm:finite-gerbe-anomaly}
The quadratic form $q_n$ determines a homotopy class
\[
\widetilde q_n:K(\Phi_n,2)\longrightarrow K(\Q/\Z,4).
\]
For every closed oriented four-manifold $X$, composition with
$\widetilde q_n$ and evaluation on $[X]$ define a quadratic action
\[
q_{n,X}:H^2(X;\Phi_n)\longrightarrow\Q/\Z.
\]
The finite gerbe path integral
\[
\mathcal A_n(X)
=
\sum_{\alpha\in H^2(X;\Phi_n)}
\frac{|H^0(X;\Phi_n)|}{|H^1(X;\Phi_n)|}
\exp\bigl(2\pi i q_{n,X}(\alpha)\bigr)
\]
is the closed-four-manifold value of an invertible four-dimensional TQFT.
In the Freed--Hopkins--Lurie--Teleman framework, an associated
three-dimensional boundary theory is a truncated morphism
$\mathbf1\to\mathcal A_n$.
\end{theorem}

\begin{proof}
Eilenberg--Mac Lane identify homogeneous quadratic maps $A\to B$ with
homotopy classes of maps $K(A,2)\to K(B,4)$.  Applying this to $q_n$
gives $\widetilde q_n$.  The finite path-integral theorem of
\cite{FHLT}, applied to the $2$-groupoid of $\Phi_n$-gerbes and the
nondegenerate form $q_n$, gives the displayed groupoid-cardinality weight,
the invertible four-dimensional theory, and its anomalous boundary
interpretation.
\end{proof}

\subsection{Surface transgression and the finite Heisenberg algebra}

For the oriented torus $T^2$, identify
\[
H^1(T^2;\Phi_n)\cong\Phi_n^2.
\]
The cup product and $b_n$ induce the alternating form
\[
\Omega_n\bigl((a,b),(c,d)\bigr)
=\frac{ad-bc}{n}\pmod\Z.
\]

\begin{definition}[Finite Weyl operators]
Let $\zeta_n=e^{2\pi i/n}$ and let
\[
\mathscr H_n=\C[\Phi_n]
\]
with basis $\delta_0,\dots,\delta_{n-1}$.  Define
\[
P_n\delta_j=\zeta_n^j\delta_j,
\qquad
Q_n\delta_j=\delta_{j+1},
\]
with indices modulo $n$.  Put
\[
\omega_n((a,b),(c,d)):=\zeta_n^{-bc}.
\]
\end{definition}

\begin{theorem}[Finite Stone--von Neumann algebra]
\label{thm:finite-stone-von-neumann}
The Weyl operators satisfy
\[
P_n^n=Q_n^n=1,
\qquad
P_nQ_n=\zeta_n Q_nP_n.
\]
The $n^2$ operators
\[
\{P_n^aQ_n^b:0\leq a,b<n\}
\]
form a basis of $\End_\C(\mathscr H_n)$.  Consequently,
\[
\C_{\omega_n}[\Phi_n^2]
\cong
\End_\C(\mathscr H_n)
\cong M_n(\C),
\]
and $\mathscr H_n$ is the irreducible Schr\"odinger module.  The
commutator of the Weyl operators is
\[
W_xW_yW_x^{-1}W_y^{-1}
=\exp\bigl(2\pi i\Omega_n(x,y)\bigr).
\]
\end{theorem}

\begin{proof}
The commutation relation follows on each basis vector.  It implies
\[
(P_n^aQ_n^b)(P_n^cQ_n^d)
=\zeta_n^{-bc}P_n^{a+c}Q_n^{b+d},
\]
which is the cocycle $\omega_n$ and gives the stated commutator.  Suppose
\[
\sum_{a,b}c_{a,b}P_n^aQ_n^b=0.
\]
Applying this operator to $\delta_j$ and projecting to
$\delta_{j+b}$ shows, for every fixed $b$, that
\[
\sum_a c_{a,b}\zeta_n^{aj}=0
\qquad\text{for all }j.
\]
The finite Fourier matrix is invertible, so every $c_{a,b}=0$.  Thus the
operators are linearly independent.  Their number equals
$\dim_\C M_n(\C)=n^2$, so they form a basis.  The image algebra is all of
$\End_\C(\mathscr H_n)$, which proves irreducibility.
\end{proof}

\subsection{The explicit bundle-gerbe realization}

Let $L_N\to T^2$ be the degree-$N$ Hermitian line bundle and let
$\mathcal G_N\to T^3$ be the explicit bundle gerbe constructed in
\cite{GerbeTheta}.  Its Dixmier--Douady class and curvature are
\[
DD(\mathcal G_N)
=N[dt]\smile[dx\wedge dy],
\qquad
H_N=2\pi iN\,dt\wedge dx\wedge dy.
\]

\begin{theorem}[Winding-sector realization of the arithmetic phase]
\label{thm:winding-theta-realization}
On the winding-$r$ loop sector,
\[
\Trg(\mathcal G_N)|_{\mathcal L_r}
\cong L_N^{\otimes r}.
\]
Put $d=rN>0$.  For each flat parameter $(\phi,\psi)$, the holomorphic
theta space has dimension $d$ and carries normalized operators
$P_{r,N},Q_{r,N}$ satisfying
\[
P_{r,N}Q_{r,N}=e^{2\pi i/d}Q_{r,N}P_{r,N}.
\]
For $N=1$ and $r=n$, there is an isomorphism of Heisenberg modules
\[
\boxed{
\mathscr H_n
\cong
H^0(E_\tau,L_1^{\otimes n}\otimes F_{\phi,\psi}).
}
\]
It is unique up to a nonzero scalar.
\end{theorem}

\begin{proof}
The transgression, dimension, theta-basis, and magnetic-translation
statements are proved in \cite{GerbeTheta}.  After the scalar
normalization which removes the flat holonomies, its theta basis
$\theta_0,\dots,\theta_{n-1}$ satisfies
\[
P\theta_j=\zeta_n^j\theta_j,
\qquad
Q\theta_j=\theta_{j+1}.
\]
Thus $\delta_j\mapsto\theta_j$ intertwines the generators in
Theorem~\ref{thm:finite-stone-von-neumann}.  Uniqueness up to scalar is
Schur's lemma for the irreducible finite Heisenberg representation.
\end{proof}

\begin{remark}[Exact level of identification]
Theorem~\ref{thm:winding-theta-realization} identifies the transgressed
torus boundary algebra and module.  It does not identify the entire
four-dimensional finite gerbe theory $\mathcal A_n$ with the bundle gerbe
$\mathcal G_1$ as bulk theories.
\end{remark}

\begin{theorem}[Frobenius as winding dilation]
\label{thm:frobenius-winding-dilation}
Let $\rho_m:S^1\to S^1$ be the degree-$m$ map, and let
$c_m:\mathcal L_n\to\mathcal L_{mn}$ be precomposition of loops by
$\rho_m$.  Then
\[
c_m^*\bigl(\Trg(\mathcal G_N)|_{\mathcal L_{mn}}\bigr)
\cong
\bigl(\Trg(\mathcal G_N)|_{\mathcal L_n}\bigr)^{\otimes m}.
\]
Thus multiplication by $m$ on arithmetic Frobenius degree is realized by
multiplication of loop winding.
\end{theorem}

\begin{proof}
A winding-$n$ loop composed with the degree-$m$ circle map has winding
$mn$.  The two line bundles in the displayed formula identify respectively
with $L_N^{\otimes mn}$ and
$(L_N^{\otimes n})^{\otimes m}$, which are canonically isomorphic.
\end{proof}

\subsection{Divisibility as a system of Lagrangian correspondences}

Let
\[
\Phi_n^\vee=\Hom(\Phi_n,\Q/\Z),
\qquad
D_n=\Phi_n\oplus\Phi_n^\vee,
\]
and equip $D_n$ with the hyperbolic quadratic form
\[
q_n^{\mathrm{hyp}}(a,\chi)=\chi(a).
\]
For $n\mid m$, let $i_{n,m}:\Phi_n\hookrightarrow\Phi_m$ be the
canonical inclusion and let
$r_{m,n}:\Phi_m^\vee\to\Phi_n^\vee$ be restriction.

\begin{definition}[Denominator correspondence]
Define
\[
\mathcal C_{n,m}
=
\left\{
\bigl((a,\chi),(i_{n,m}a,\psi)\bigr):
\chi=r_{m,n}\psi
\right\}
\subset D_n^-\oplus D_m.
\]
\end{definition}

\begin{theorem}[Lagrangian denominator history]
\label{thm:lagrangian-denominator-history}
The subgroup $\mathcal C_{n,m}$ is Lagrangian.  If
$n\mid m\mid\ell$, then relational composition satisfies
\[
\mathcal C_{m,\ell}\circ\mathcal C_{n,m}
=\mathcal C_{n,\ell}.
\]
Hence the divisibility category maps functorially to the category of finite
quadratic groups and Lagrangian correspondences.

Moreover, $\widehat\Z^\times$ acts by
\[
(a,\chi)\longmapsto(ua,u^{-1}\chi),
\]
preserves $q_n^{\mathrm{hyp}}$, and preserves every
$\mathcal C_{n,m}$.  The Lagrangian history is therefore canonically
cyclotomic-equivariant.
\end{theorem}

\begin{proof}
On $\mathcal C_{n,m}$ one has
\[
-q_n^{\mathrm{hyp}}(a,\chi)
+q_m^{\mathrm{hyp}}(i_{n,m}a,\psi)
=-\chi(a)+\psi(i_{n,m}a)=0.
\]
Thus the subgroup is isotropic.  Its cardinality is
$|\Phi_n||\Phi_m^\vee|=nm$, which is the square root of
$|D_n^-\oplus D_m|=n^2m^2$; hence it is maximal isotropic.
Composition follows from functoriality of inclusion and restriction.
The unit action preserves the hyperbolic form because
$(u^{-1}\chi)(ua)=\chi(a)$, and inclusion and restriction commute with
the reductions of a profinite unit, so the correspondences are invariant.
\end{proof}

\section{Quantized denominator history and the universal UHF limit}
\label{sec:uhf-history}

The Lagrangian correspondence system is canonical at the Morita level.  A
strict inductive system of matrix algebras requires an additional choice.
For $n\mid m$, write $m=kn$ and define
\[
\iota_{n,m}:M_n(\C)\longrightarrow M_m(\C)
\]
on matrix units by
\[
\boxed{
\iota_{n,m}(E_{ab}^{(n)})
=\sum_{r=0}^{k-1}E_{a+nr,b+nr}^{(m)}.
}
\]
These embeddings are unital and satisfy
$\iota_{m,\ell}\iota_{n,m}=\iota_{n,\ell}$.

\begin{definition}[Perfected Heisenberg algebra]
\label{def:perfected-uhf}
Define
\[
\mathcal Q_\Omega
:=\varinjlim_{n\mid m}(M_n(\C),\iota_{n,m}).
\]
\end{definition}

\begin{theorem}[UHF algebra of the globally perfect denominator]
\label{thm:perfected-uhf}
The algebra $\mathcal Q_\Omega$ is the universal UHF algebra of
supernatural type $\Omega=\prod_pp^\infty$.  Its ordered $K$-theory is
\[
K_0(\mathcal Q_\Omega)
\cong(\Q,\Q_{\geq0},1),
\qquad
K_1(\mathcal Q_\Omega)=0.
\]
\end{theorem}

\begin{proof}
At stage $n$, identify $K_0(M_n(\C))$ with $\Z$ using a rank-one
projection as generator.  Under $\iota_{n,m}$ a rank-one projection maps
to rank $m/n$, so the connecting map is multiplication by $m/n$.  The
map
\[
\Z\ni r\text{ at stage }n
\longmapsto\frac rn\in\Q
\]
is compatible with all connecting maps and identifies the direct limit
with $\Q$.  Positive ranks give $\Q_{\geq0}$ and the identity projection
gives the order unit $1$.  Since $K_1(M_n(\C))=0$ at every stage and
$K$-theory commutes with inductive limits, $K_1$ vanishes.  The directed
set contains arbitrarily large powers of every prime, so the UHF
supernatural type is $\Omega$.
\end{proof}

\begin{theorem}[The profinite diagonal]
\label{thm:profinite-diagonal}
Let $\mathcal D_n\subset M_n(\C)$ be the diagonal algebra, identified with
$C(\Z/n\Z)$.  The embeddings restrict to pullback along reduction
$\Z/m\Z\to\Z/n\Z$.  Therefore
\[
\mathcal D_\Omega
:=\varinjlim_n\mathcal D_n
\cong C(\widehat\Z).
\]
By Pontryagin duality,
\[
\boxed{
\mathcal D_\Omega
\cong C^*(\Q/\Z)
\subset\mathcal Q_\Omega.
}
\]
The compact space $\widehat\Z$ is a Cantor space.
\end{theorem}

\begin{proof}
Continuous functions on
$\widehat\Z=\varprojlim_n\Z/n\Z$ are the norm completion of locally
constant functions, and those are exactly the direct limit of the finite
diagonal algebras.  Pontryagin duality identifies the group $C^*$-algebra
of the discrete group $\Q/\Z$ with continuous functions on its compact
dual $\widehat\Z$.  Finally, $\widehat\Z$ is compact, metrizable, totally
disconnected, and has no isolated points; hence it is a Cantor space.
\end{proof}

\subsection{The strict cyclotomic carry obstruction}

For an integer $x$ and $q\geq1$, let $\langle x\rangle_q$ denote the
representative of $x\bmod q$ in $\{0,\dots,q-1\}$.  Let $m=kn$ and let
$u\in(\Z/m\Z)^\times$.  Define
\[
u_n(a):=\langle ua\rangle_n,
\qquad
c_{u;n,m}(a)
:=\frac{\langle ua\rangle_m-\langle ua\rangle_n}{n}
\in\Z/k\Z.
\]
The quotient is integral because the two residues agree modulo $n$.

\begin{theorem}[Cyclotomic carry cocycle and failure of strict equivariance]
\label{thm:cyclotomic-carry-obstruction}
Let $U_q(u)$ be the permutation unitary on $\C[\Z/q\Z]$ defined by
$U_q(u)e_a=e_{\langle ua\rangle_q}$.  Then:
\begin{enumerate}[label=\textup{(\arabic*)}]
\item for $a\in\Z/n\Z$ and $r\in\Z/k\Z$,
\[
\langle u(a+nr)\rangle_m
=u_n(a)+n\langle ur+c_{u;n,m}(a)\rangle_k;
\]
\item the carry functions satisfy the crossed cocycle law
\[
c_{uv;n,m}(a)
=u\,c_{v;n,m}(a)+c_{u;n,m}(v_n(a))
\pmod k;
\]
\item on matrix units,
\[
\Ad(U_m(u))\iota_{n,m}(E_{ab})
=
\sum_{t\in\Z/k\Z}
E_{u_n(a)+nt,\,
  u_n(b)+n(t+c_{u;n,m}(b)-c_{u;n,m}(a))};
\]
\item the square
\[
\begin{tikzcd}[column sep=large,row sep=large]
M_n(\C) \arrow[r,"\iota_{n,m}"]
  \arrow[d,"\Ad(U_n(u))"']
& M_m(\C) \arrow[d,"\Ad(U_m(u))"]\\
M_n(\C) \arrow[r,"\iota_{n,m}"']
& M_m(\C)
\end{tikzcd}
\]
commutes on $E_{ab}$ exactly when
$c_{u;n,m}(a)=c_{u;n,m}(b)$.  It commutes on the entire diagonal for every
$u$, and on the entire matrix algebra exactly when
$c_{u;n,m}$ is identically zero.
\end{enumerate}
Consequently, $\widehat\Z^\times$ acts canonically on the profinite
diagonal $C(\widehat\Z)$, but the naive finite permutation operators do
not define a strict action on the chosen UHF inductive system.  A strict
lift requires coherent trivializations of the carry cocycle.
\end{theorem}

\begin{proof}
Write the residue of $ua$ modulo $m$ uniquely as
$u_n(a)+nc_{u;n,m}(a)$.  Adding $nur$ and reducing modulo $m=nk$ gives
(1).  Applying (1) successively to $v$ and then $u$, and comparing with
$uv$, gives (2).

Conjugating a summand of the matrix embedding gives
\[
U_m(u)E_{a+nr,b+nr}U_m(u)^{-1}
=E_{u_n(a)+n(ur+c_u(a)),\,
    u_n(b)+n(ur+c_u(b))},
\]
where fibre coordinates are taken modulo $k$.  Since $u$ is invertible
modulo $k$, reindex by $t=ur+c_u(a)$ to obtain (3).  The lower-left then
has zero relative fibre displacement, so equality on $E_{ab}$ is
precisely $c_u(a)=c_u(b)$.  For $a=b$ this always holds, proving diagonal
equivariance.  Equality for all matrix units forces $c_u$ to be constant;
since $c_u(0)=0$, it is identically zero.  Passing to the direct limit of
the diagonals gives the canonical action on $C(\widehat\Z)$.
\end{proof}

\begin{remark}[Morita equivariance versus strict equivariance]
The Lagrangian correspondence system of
Theorem~\ref{thm:lagrangian-denominator-history} is strictly
$\widehat\Z^\times$-equivariant.  The obstruction appears only after one
chooses the matrix embeddings $\iota_{n,m}$.  This is the precise sense in
which denominator history is canonical at the correspondence/Morita level
but carries a nontrivial arithmetic correction at the strict operator
level.
\end{remark}

\begin{remark}[Exact relation to the Bost--Connes coefficient algebra]
The equality
\[
C^*(\Q/\Z)=C(\widehat\Z)=\mathcal D_\Omega
\]
is literal.  The full Bost--Connes system additionally uses the
$\N^\times$ endomorphisms and a crossed-product construction.  Thus
$\mathcal Q_\Omega$ is a UHF thickening of the cyclotomic coefficient
algebra, not a replacement for the complete BC dynamical system.
\end{remark}

\subsection{Split support for anomaly lines}

For a complex vector space $V$, put
\[
\Spl(V)=V\sqcup\{\tau_V\},
\]
where $\tau_V$ is external absence and $0_V\in V$ is the supported zero.
A linear map $f:V\to W$ extends by
\[
\Spl(f)(\tau_V)=\tau_W,
\qquad
\Spl(f)(v)=f(v).
\]

\begin{proposition}[Defined vanishing versus absent amplitude]
\label{prop:split-anomaly-line}
The assignment $V\mapsto\Spl(V)$ is a functor from complex vector spaces
to pointed sets.  If $v$ is supplied and $f(v)=0_W$, then
\[
\Spl(f)(v)=0_W\neq\tau_W.
\]
Consequently, for an anomalous field theory whose value on a closed
$d$-manifold is a possibly zero vector in an anomaly line, split support
distinguishes
\[
\begin{array}{c|c}
\tau&\text{the tier or field assignment is absent},\\
0&\text{the path integral is defined and vanishes},\\
z\neq0&\text{a nonzero anomalous amplitude}.
\end{array}
\]
The ordinary projection $\Spl(V)\to V$ identifies the first two states,
whereas the support character separates them.
\end{proposition}

\begin{proof}
Identity maps and composition are preserved by the definition.  A supplied
vector always remains in the supported copy of the target, even when its
ordinary linear image is zero.  The final statement applies this observation
to the anomaly line assigned by a truncated morphism
$\mathbf1\to\mathcal A$.
\end{proof}

\subsection{Dual anomalies and Morita-level identity}

\begin{proposition}[Anomaly cancellation]
\label{prop:anomaly-cancellation}
The inverse of the finite quadratic theory associated with $q_n$ is the
theory associated with $-q_n$, and
\[
\mathcal A(q_n)\otimes\mathcal A(-q_n)\simeq\mathbf1.
\]
For the explicit bundle gerbe,
\[
DD(\mathcal G_N^\vee)=-DD(\mathcal G_N),
\]
so $\mathcal G_N\otimes\mathcal G_N^\vee$ is stably trivial.  On winding
sectors,
\[
L_N^{\otimes r}\otimes L_N^{\otimes(-r)}\cong\mathbf1.
\]
\end{proposition}

\begin{proof}
The finite theory is invertible and changing $q_n$ to $-q_n$ complex
conjugates every Gauss phase, hence gives the tensor inverse.  Dualizing a
bundle gerbe negates its Dixmier--Douady class.  The last isomorphism is the
evaluation pairing between a line bundle and its dual.
\end{proof}

\begin{theorem}[Morita identification of the finite models]
\label{thm:morita-univalent-finite}
After choosing the standard Lagrangian polarization of $\Phi_n^2$, there
are algebra isomorphisms
\[
\C_{\omega_n}[\Phi_n^2]
\cong
\End_\C(\mathscr H_n)
\cong
\End_\C H^0(E_\tau,L_1^{\otimes n}\otimes F_{\phi,\psi})
\cong M_n(\C).
\]
The matrix algebra is Morita equivalent to $\C$ through the bimodule
$\C^n$.  Therefore the finite phase, Heisenberg, theta, and matrix
presentations determine the same boundary object in the Morita bicategory.
In a univalent completion of that bicategory, this equivalence determines
a path between the represented objects; it is not an elementwise equality
of the original presentations.
\end{theorem}

\begin{proof}
The first and last algebra isomorphisms are
Theorem~\ref{thm:finite-stone-von-neumann}; the middle one is
Theorem~\ref{thm:winding-theta-realization}.  The standard
$(M_n(\C),\C)$-bimodule $\C^n$ and its dual implement the Morita
equivalence.  The final sentence is the defining interpretation of
equivalence in a univalent or Rezk completion.
\end{proof}

\begin{figure}[htbp]
\centering
\begin{tikzcd}[column sep=large,row sep=large]
(\Phi_n,q_n) \arrow[r,"\mathrm{deloop}"]
  \arrow[d,"\mathrm{transgress}_{T^2}"']
& \widetilde q_n:K(\Phi_n,2)\to K(\Q/\Z,4)
  \arrow[r,"\mathrm{finite\ path\ integral}"]
& \mathcal A_n \\
(\Phi_n^2,\Omega_n) \arrow[r,"\omega_n"']
& \C_{\omega_n}[\Phi_n^2]
  \arrow[r,"\rho_n"']
& \End_\C(\mathscr H_n)\cong M_n(\C)
\end{tikzcd}
\caption{Finite denominator phase, Eilenberg--Mac Lane delooping,
gerbe integration, surface transgression, and Heisenberg quantization.}
\label{fig:phase-gerbe-quantization}
\end{figure}

\section{Beurling realization of the reciprocal-root skeleton}
\label{sec:beurling-realization}

We now connect the perfect stalk $\Fone[T^{\Q_+}]$ with the reciprocal
Nyman--Beurling system studied in \cite{Schultze}.  Let
\[
\cH=L^2((1,\infty),y^{-2}\,dy)
\]
and, for $0<\alpha<1$, define
\[
F_\alpha(y)=\alpha\lfloor y\rfloor-\lfloor\alpha y\rfloor.
\]
Since
\[
F_\alpha(y)=\{\alpha y\}-\alpha\{y\},
\]
these functions are bounded and therefore belong to $\cH$.

For $Q\geq2$, put
\[
V_Q=\operatorname{span}\{F_{a/Q}:1\leq a\leq Q-1\},
\]
let $P_Q$ be orthogonal projection onto $V_Q$, and define
\[
r_Q=(I-P_Q)1,
\qquad
x_Q=\|r_Q\|^2
=\operatorname{dist}_{\cH}(1,V_Q)^2.
\]
Define the reciprocal subspace and distance
\[
\cR_Q=\operatorname{span}\{F_{1/n}:2\leq n\leq Q\},
\qquad
 d_Q=\operatorname{dist}_{\cH}(1,\cR_Q).
\]

\begin{definition}[Beurling realization of the perfect stalk]
\label{def:beurling-realization}
On the complex span of monomials $T^q$, $q\in\Q\cap(0,1)$, define
\[
\mathfrak B_\Omega(T^q)=F_q.
\]
The reciprocal-root skeleton is
\[
\mathscr R_\Omega=\{T^{1/n}:n\geq2\}
\subset\Fone[T^{\Q_+}].
\]
This is a linear analytic realization of stalk monomials; it is not an
$\Fone$-algebra morphism.
\end{definition}

\begin{theorem}[Direct Frobenius-refinement domination]
\label{thm:direct-beurling-domination}
For every $Q\geq2$,
\[
\cR_Q\subseteq V_{\Lambda_Q}
\]
and consequently
\[
\boxed{x_{\Lambda_Q}\leq d_Q^2.}
\]
\end{theorem}

\begin{proof}
For $2\leq n\leq Q$,
\[
\frac1n=\frac{\Lambda_Q/n}{\Lambda_Q},
\]
where $1\leq\Lambda_Q/n\leq\Lambda_Q-1$.  Hence
$F_{1/n}\in V_{\Lambda_Q}$, proving the subspace inclusion.  Distance to
a larger closed subspace cannot increase, so
\[
\operatorname{dist}(1,V_{\Lambda_Q})
\leq\operatorname{dist}(1,\cR_Q)=d_Q.
\]
Squaring gives the result.
\end{proof}

\begin{corollary}[Strengthened refinement alternative]
\label{cor:strong-refinement-alternative}
For each $Q\geq2$, either
\[
x_Q\leq4d_Q^2,
\]
or
\[
x_{\Lambda_Q}<\frac14x_Q
\qquad\text{and}\qquad
x_Q-x_{\Lambda_Q}>\frac34x_Q.
\]
\end{corollary}

\begin{proof}
If the first inequality fails, then $d_Q^2<x_Q/4$.  Apply Theorem
\ref{thm:direct-beurling-domination}.
\end{proof}

The estimate is stronger than the quadratic decrement proved in
\cite{Schultze}; it follows from the exact root inclusion and does not
require the unit-step visibility lemma.

\begin{lemma}[Reciprocal atoms in Nyman form]
\label{lem:nyman-form}
The map
\[
U:\cH\longrightarrow L^2(0,1),
\qquad
(Uf)(x)=f(1/x),
\]
is an isometry and satisfies
\[
UF_{1/n}(x)
=
\rho\!\left(\frac1{nx}\right)
-\frac1n\rho\!\left(\frac1x\right),
\]
where $\rho(t)=t-\lfloor t\rfloor$.
\end{lemma}

\begin{proof}
The substitution $y=1/x$ gives $y^{-2}dy=-dx$, proving the isometry.
Furthermore,
\[
\begin{aligned}
F_{1/n}(1/x)
&=\frac1n\left(\frac1x-\rho(1/x)\right)
 -\left(\frac1{nx}-\rho(1/(nx))\right)\\
&=\rho(1/(nx))-\frac1n\rho(1/x).
\end{aligned}
\]
\end{proof}

\begin{theorem}[Arithmetic-site form of the integer
Nyman--B\'aez-Duarte criterion]
\label{thm:arithmetic-site-nyman}
Let $Q_{j+1}=\Lambda_{Q_j}$ with $Q_0\geq3$.  The following are equivalent:
\begin{enumerate}[label=\textup{(\roman*)}]
\item the Riemann hypothesis;
\item $d_Q\to0$ as $Q\to\infty$;
\item $x_{Q_j}\to0$ as $j\to\infty$;
\item
\[
1\in
\overline{
\mathfrak B_\Omega
\bigl(\operatorname{span}_{\C}\mathscr R_\Omega\bigr)
}^{\,\cH}.
\]
\end{enumerate}
\end{theorem}

\begin{proof}
By the integer-dilation form of the Nyman--Beurling criterion due to
B\'aez-Duarte \cite{Nyman,Beurling,BaezDuarte}, the Riemann hypothesis is
equivalent to density of the functions in Lemma \ref{lem:nyman-form} at
the target $1$.  This is exactly the equivalence of (i), (ii), and (iv).

Assume (ii).  By Theorem \ref{thm:direct-beurling-domination},
\[
x_{Q_{j+1}}\leq d_{Q_j}^2\longrightarrow0,
\]
so (iii) follows.

Conversely, every $V_{Q_j}$ is contained in the full Nyman--Beurling
space generated by the functions
$\rho(\alpha/x)-\alpha\rho(1/x)$, $0<\alpha<1$.  If (iii) holds, then $1$
lies in the closure of that space, and the Nyman--Beurling theorem implies
(i).
\end{proof}

\begin{warning}[Claim boundary]
Theorem \ref{thm:arithmetic-site-nyman} is an equivalence, not a proof of
RH.  The unproved endpoint is exactly the density assertion in (ii) or
(iv).  The split-support and arithmetic-site reformulation locates the
obstruction at the globally Frobenius-perfect point; it does not remove
that obstruction.
\end{warning}

\begin{theorem}[Mellin transform and Frobenius characters]
\label{thm:mellin-frobenius-characters}
For $0<\alpha<1$ and $\Re(s)>1$,
\[
\int_1^\infty F_\alpha(y)y^{-s-1}\,dy
=
\frac{\zeta(s)}s\bigl(\alpha-\alpha^s\bigr).
\]
In particular,
\[
\mathcal MF_{1/n}(s)
=
\frac{\zeta(s)}s
\bigl(\chi_1(n)-\chi_s(n)\bigr),
\qquad
\chi_s(n)=n^{-s}.
\]
\end{theorem}

\begin{proof}
For $\Re(s)>1$,
\[
\lfloor y\rfloor=\sum_{m\geq1}\mathbf1_{[m,\infty)}(y)
\]
with absolute convergence after integration, whence
\[
\int_1^\infty\lfloor y\rfloor y^{-s-1}\,dy
=\frac1s\sum_{m\geq1}m^{-s}
=\frac{\zeta(s)}s.
\]
Similarly,
\[
\int_1^\infty\lfloor\alpha y\rfloor y^{-s-1}\,dy
=\frac1s\sum_{m\geq1}(m/\alpha)^{-s}
=\frac{\alpha^s\zeta(s)}s.
\]
Subtracting gives the first identity.  The second is the specialization
$\alpha=1/n$.
\end{proof}

\section{Frobenius local systems and the helix multiplier}
\label{sec:helix-local-system}

The helix construction of \cite{Lavery} is most naturally expressed on the
Connes--Consani periodic orbit.  Fix a prime $p$ and write
\[
C_p=\R_+^\times/p^\Z
\cong\R/(\log p)\Z.
\]
Let $u=\log\lambda$ be the coordinate on the universal cover.

\begin{definition}[Frobenius phase line]
\label{def:frobenius-phase-line}
For $t\in\R$, let $L_{p,t}\to C_p$ be the complex line bundle
\[
L_{p,t}=(\R\times\C)/\sim,
\]
where
\[
(u,z)\sim(u+k\log p,p^{-ikt}z),
\qquad k\in\Z.
\]
Let $\Delta_p^{1/2}$ denote the real one-dimensional factor of automorphy
$p^{1/2}$.  Define
\[
\cL^{\mathrm{hel}}_{p,t}=\Delta_p^{1/2}\otimes L_{p,t},
\qquad
\cL^{\mathrm{crit}}_{p,t}=\Delta_p^{-1/2}\otimes L_{p,t}.
\]
\end{definition}

\begin{theorem}[Helix phase as Frobenius holonomy]
\label{thm:helix-holonomy}
The function
\[
\psi_t(u)=e^{-itu}
\]
is an equivariant lift of a section of $L_{p,t}$ and satisfies
\[
\psi_t(u+\log p)=p^{-it}\psi_t(u).
\]
Consequently,
\[
\operatorname{Hol}(L_{p,t})=p^{-it},
\]
while the two twists have Frobenius factors
\[
p^{1/2-it}
\qquad\text{and}\qquad
p^{-1/2-it},
\]
respectively.
\end{theorem}

\begin{proof}
Direct calculation gives
\[
\psi_t(u+k\log p)
=e^{-it(u+k\log p)}
=p^{-ikt}\psi_t(u),
\]
which is exactly the factor-of-automorphy condition.  Tensoring with the
half-density factors multiplies the holonomy by $p^{\pm1/2}$.
\end{proof}

\begin{corollary}[Dual chirality and determinant one]
\label{cor:dual-chirality}
Complex conjugation gives
\[
L_{p,-t}\cong L_{p,t}^{\vee},
\]
and therefore
\[
\det(L_{p,t}\oplus L_{p,-t})\cong\mathbf1.
\]
In a local trivialization, Frobenius acts by
\[
\begin{pmatrix}
p^{-it}&0\\0&p^{it}
\end{pmatrix},
\]
which is unitary of determinant one.
\end{corollary}

\begin{proof}
The inverse of the character $p^{-it}$ is $p^{it}$, which is also its
complex conjugate.  Hence the second line is the dual of the first, and
the determinant line is their tensor product.
\end{proof}

\begin{proposition}[Descent and winding closure]
\label{prop:helix-descent}
The equivariant section $\psi_t$ descends to a scalar-valued function on
$C_p$ if and only if
\[
p^{-it}=1,
\]
equivalently
\[
t\in\frac{2\pi}{\log p}\Z.
\]
On the rectangular Tate torus
\[
E_p\cong C_p\times S^1
\]
from \cite{CCAbsolute}, the trajectory
\[
u\longmapsto([u],[-tu])
\]
is closed if $t\log p/(2\pi)\in\Q$ and dense if this quotient is
irrational.
\end{proposition}

\begin{proof}
Scalar descent requires invariance under the deck translation
$u\mapsto u+\log p$, which is equivalent to $p^{-it}=1$.  This gives the
stated integral condition.  The final assertion is the standard
classification of linear flows on a two-dimensional torus: the orbit is
closed precisely when the two periods are rationally dependent.
\end{proof}

\begin{corollary}[Split fixed-locus jump for Frobenius phase]
\label{cor:helix-fixed-locus-jump}
Let $u_{p,t}=p^{-it}\in\C^\times$.  Multiplication by $u_{p,t}$ on
$G(\C)$ has fixed locus
\[
\operatorname{Fix}_{G(\C)}(u_{p,t})=
\begin{cases}
G(\C),&t\log p\in2\pi\Z,\\
\{\tau,e\}\cong\B,&t\log p\notin2\pi\Z.
\end{cases}
\]
When the holonomy is nontrivial, the two chiral phase lines have fixed locus $\B^2$.
Thus scalar descent is exactly the parameter locus at which invariant
amplitude jumps from Boolean support to the full split complex fibre.
\end{corollary}

\begin{proof}
Combine Proposition~\ref{prop:helix-descent} with
Theorem~\ref{thm:fixed-locus-unit}.  The two-chirality statement is the
Cartesian product of the fixed loci for $u_{p,t}$ and $u_{p,t}^{-1}$.
\end{proof}

\begin{proposition}[Supernatural endpoint obstruction]
\label{prop:supernatural-holonomy-obstruction}
Consider the $p$-power ray
\[
1,p,p^2,\dots,p^\infty
\]
in the supernatural number space, and the unitary Frobenius character
\[
\chi_t(p^k)=p^{-ikt}.
\]
The sequence $\chi_t(p^k)$ has a limit in $S^1$ if and only if
$p^{-it}=1$.  Thus the same scalar is simultaneously:
\begin{enumerate}[label=\textup{(\roman*)}]
\item the holonomy of $L_{p,t}$;
\item the obstruction to scalar descent on $C_p$;
\item the obstruction to continuous extension of the phase along the
$p$-power ray to $p^\infty$.
\end{enumerate}
The full helix multiplier $p^{k(1/2-it)}$ never converges in $\C^\times$.
\end{proposition}

\begin{proof}
Put $z=p^{-it}\in S^1$.  If $z^k\to L$, then also
$z^{k+1}\to L$, while $z^{k+1}=zz^k\to zL$.  Since $L\neq0$, one has
$z=1$.  The converse is immediate.  Finally,
$|p^{k(1/2-it)}|=p^{k/2}\to\infty$.
\end{proof}

\begin{remark}
The three-dimensional growing helix is an extrinsic metric realization.
The invariant arithmetic-geometric object is the Frobenius-equivariant
line with factor $n^{1/2-it}$, together with its dual chirality.  This
reformulation preserves the Lean-verified determinant and multiplicativity
statements of \cite{Lavery} while separating them from any claim about the
location of all analytic zeros.
\end{remark}

\section{Split local moduli and Boolean place cubes}
\label{sec:split-local-moduli}

Connes--Consani show that, for $H_p=\Z[1/p]$ and each
$v\in\{p,\infty\}$, the nontrivial local points form a principal
homogeneous space $\cM_p^v$ under the local Weil group
\[
W_p=\Q_p^\times,
\qquad
W_\infty=\C^\times.
\]
The split-support refinement distinguishes absence from the neutral local
morphism before taking the Weil quotient.

\begin{definition}[Split completion of a torsor]
\label{def:split-torsor-completion}
Let a group $W$ act simply transitively on a nonempty set $M$.  Define
\[
\widetilde M=\{\tau_M\}\sqcup\{\mathbf1_M\}\sqcup M,
\]
where $W$ fixes $\tau_M$ and $\mathbf1_M$ and retains its torsor action on
$M$.  Define the support map
\[
\sigma_M:\widetilde M\to\B
\]
by
\[
\sigma_M(\tau_M)=0,
\qquad
\sigma_M(\mathbf1_M)=1,
\qquad
\sigma_M(m)=1\quad(m\in M).
\]
\end{definition}

\begin{theorem}[Univalent quotient of a split torsor]
\label{thm:split-torsor-groupoid}
For the action groupoid one has an equivalence
\[
\boxed{
[\widetilde M/W]\simeq BW\sqcup BW\sqcup *.
}
\]
The two copies of $BW$ correspond to absence and the supported neutral
point; the terminal component corresponds to the free torsor orbit.
\end{theorem}

\begin{proof}
The action preserves the disjoint decomposition
$\widetilde M=\{\tau_M\}\sqcup\{\mathbf1_M\}\sqcup M$.  A fixed singleton
has action groupoid $BW$.  Since the action on $M$ is simply transitive,
its action groupoid is equivalent to the terminal groupoid.  Taking the
disjoint union gives the result.
\end{proof}

\begin{corollary}[The split $p$--$\infty$ square]
\label{cor:p-infinity-square}
For $H_p=\Z[1/p]$, define
\[
\widetilde\cM_{p,\infty}
=\widetilde\cM_p^p\times\widetilde\cM_p^\infty.
\]
Its coarse orbit stratification has nine states
\[
\{\text{absent},\text{neutral},\text{nontrivial}\}^{\{p,\infty\}},
\]
and the support map
\[
\widetilde\cM_{p,\infty}\longrightarrow\B^2
\]
forgets only the distinction between neutral and nontrivial supported
states.  At groupoid level,
\[
[\widetilde\cM_{p,\infty}/(W_p\times W_\infty)]
\simeq
(BW_p\sqcup BW_p\sqcup *)
\times
(BW_\infty\sqcup BW_\infty\sqcup *).
\]
\end{corollary}

\begin{proof}
Apply Theorem \ref{thm:split-torsor-groupoid} independently at the two
places and take products.  The coarse states are the nine products of the
three one-place states, while the Boolean support quotient records only
whether each coordinate is absent or supported.
\end{proof}

\begin{theorem}[Finite Boolean place cubes]
\label{thm:boolean-place-cubes}
Let $H$ be an arithmetic-site point and let
$J\subseteq\Sigma_H$ be finite.  Suppose each $v\in J$ carries a
nontrivial local torsor $\cM_H^v$ under $W_v$.  Then the split local moduli
object
\[
\widetilde\cM_H^J
=\prod_{v\in J}\widetilde\cM_H^v
\]
has a canonical support map to $\B^J$.  The fibre over a mask
$A\subseteq J$ is the disjoint union, over all choices of neutral versus
nontrivial state on $A$, of products of the corresponding local torsors.
The quotient groupoid retains the stabilizer $BW_v$ at every absent or
neutral coordinate.
\end{theorem}

\begin{proof}
Take the product of the one-place support maps.  A coordinate with Boolean
value $0$ is forced to be absent.  A coordinate with Boolean value $1$ may
be neutral or belong to the nontrivial torsor.  The action-groupoid
statement follows from Theorem \ref{thm:split-torsor-groupoid} by finite
products.
\end{proof}

\begin{remark}[Why orbit sets are insufficient]
Passing directly to coarse orbit sets turns the three one-place components
into three points.  The action groupoid retains $W_v$ as the identity type
of the absent and neutral states.  This is the precise univalent content of
the split completion: the refinement is stacky, not merely set-theoretic.
\end{remark}

\section{Sheafwise split support, the split arithmetic site, and the split scaling site}
\label{sec:split-scaling-site}

The correct sheafwise construction is internal.  Let $\mathscr E$ be a
Grothendieck topos and let $A$ be a commutative semiring object of
$\mathscr E$.

\begin{definition}[Internal split globalization]
\label{def:internal-split}
Define the object
\[
\Spl_{\mathscr E}(A):=1\amalg A.
\]
On the $A$-summand retain the operations of $A$; let the point of the
$1$-summand be the new additive identity and multiplicative absorber.
\end{definition}

Because a topos is extensive, maps out of finite coproducts are determined
by their restrictions to the summands.  The semiring axioms can therefore
be checked on the finitely many support cases.

\begin{theorem}[Internal functoriality and universal maps]
\label{thm:internal-split-functor}
The assignment $A\mapsto\Spl_{\mathscr E}(A)$ is functorial on commutative
semiring objects.  It carries natural morphisms
\[
p_A:\Spl_{\mathscr E}(A)\to A,
\qquad
\chi_A:\Spl_{\mathscr E}(A)\to\B_{\mathscr E},
\]
where $p_A$ sends the new point to $0_A$ and $\chi_A$ records the coproduct
summand.  At every point $x$ of a topos with enough points,
\[
x^*\Spl_{\mathscr E}(A)\cong G(x^*A).
\]
\end{theorem}

\begin{proof}
For a semiring morphism $f:A\to A'$, define
$\Spl(f)$ to be the identity on the $1$-summand and $f$ on the supported
summand.  The case-by-case definitions make this a semiring morphism and
make composition immediate.  The maps $p_A$ and $\chi_A$ are likewise
specified on the two summands and preserve both operations.  An inverse
image functor of a geometric morphism preserves finite coproducts and
finite limits, hence transports the construction to
$1\amalg x^*A=G(x^*A)$.
\end{proof}

For sheaves on a topological space this internal coproduct has the explicit
form previously obtained by sheafification.

\begin{corollary}[Clopen support formula]
\label{cor:clopen-support-formula}
If $A=\cA$ is a sheaf of semirings on a topological space $X$, then
\[
\Spl(\cA)(U)
\cong
\coprod_{V\in\Clop(U)}\cA(V).
\]
A pair $(V,a)$ is the section equal to $a$ on the clopen support locus $V$
and absent on $U\setminus V$.  On connected $U$ this reduces to
$G(\cA(U))$.
\end{corollary}

\begin{proof}
A map $U\to1\amalg\cA$ has a locally constant support map
$U\to\{0,1\}$; its supported inverse image is clopen, and on that locus it
is exactly a section of $\cA$.  This correspondence is reversible and
compatible with the semiring operations.
\end{proof}

\subsection{The split arithmetic site}

The original arithmetic site is
\[
\mathfrak A=(\widehat{\N^\times},\Nbar),
\qquad
\Nbar=(\N\cup\{\infty\},\min,+),
\]
with the $n$-th Frobenius acting by multiplication by $n$ on finite
exponents and fixing $\infty$ \cite{CCArithmeticSite}.  Here $\infty$ is
the additive zero of the tropical semiring.

\begin{definition}[Split arithmetic site]
\label{def:split-arithmetic-site}
Define
\[
\mathfrak A^{\mathrm{sp}}
:=\bigl(\widehat{\N^\times},\Spl(\Nbar)\bigr).
\]
Write $\tau_{\mathrm{ar}}$ for the new zero and
$e_{\mathrm{ar}}:=\infty\in\Nbar$ for the supported tropical zero.
\end{definition}

\begin{theorem}[Frobenius-equivariant two-zero arithmetic site]
\label{thm:split-arithmetic-site}
Every universal Frobenius of $\Nbar$ extends uniquely to
$\Spl(\Nbar)$ by fixing $\tau_{\mathrm{ar}}$.  The support character
\[
\chi:\Spl(\Nbar)\to\B
\]
is $\N^\times$-equivariant.  At every point of the arithmetic topos the
stalk separates
\[
\tau_{\mathrm{ar}},
\qquad
e_{\mathrm{ar}}=\infty,
\qquad
n\in\N.
\]
Localizing at $e_{\mathrm{ar}}$ gives the Boolean support residue, whereas
localizing at $\tau_{\mathrm{ar}}$ gives the terminal semiring.
\end{theorem}

\begin{proof}
Functoriality in Theorem~\ref{thm:internal-split-functor} extends each
Frobenius and makes $\chi$ equivariant.  The stalk statement follows from
inverse-image compatibility.  The localization statement is the
idempotent division trichotomy applied stalkwise; the old tropical zero is
idempotent for multiplication because $\infty+\infty=\infty$.
\end{proof}

\begin{remark}[Synchronized and independent squares]
Theorem~\ref{thm:boolean-fibre-product} internalizes to semiring objects:
\[
\Spl(A\times B)\cong\Spl(A)\times_{\B}\Spl(B).
\]
This is the synchronized-support square.  The Cartesian product
$\Spl(A)\times\Spl(B)$ is larger and carries the full independent Boolean
square.  This distinction is separate from the Connes--Consani reduced and
unreduced tensor squares of the arithmetic site, which use
$\Nbar\otimes_{\B}\Nbar$ and its Newton-polygon reduction rather than a
Cartesian product.
\end{remark}

\subsection{The split scaling site}

Let
\[
\mathscr S=([0,\infty)\rtimes\N^\times,\cO)
\]
be the Connes--Consani scaling site, whose characteristic-one structure
sheaf consists of convex piecewise affine functions with integral slopes
\cite{CCScalingSite,CCAbsolute}.

\begin{definition}[Split scaling site]
\label{def:split-scaling-site}
Define
\[
\mathscr S^{\mathrm{sp}}
:=([0,\infty)\rtimes\N^\times,\Spl(\cO)).
\]
\end{definition}

\begin{theorem}[Three local scaling states]
\label{thm:three-scaling-states}
At every point $x$ of the scaling site, the split stalk separates
\[
\tau_x,
\qquad
e_x=0_{\cO_x},
\qquad
f\in\cO_x\setminus\{0\}.
\]
These are, respectively, absence of a local valuation profile, a present
profile equal to tropical zero, and a nonzero profile.  Boolean
localization at $e_x$ forgets all slope and amplitude data and retains only
support.
\end{theorem}

\begin{proof}
Apply Theorem~\ref{thm:internal-split-functor} at the point $x$ and then
the scalar localization theorem.
\end{proof}

\begin{warning}[The spherical coefficient boundary]
The construction above applies to ordinary semiring objects such as
$\Nbar$ and $\cO$.  The absolute structure sheaf
$\cO_{\Fone}=\Fone[T]$ is a spherical algebra or $\Gamma$-ring.  A split
construction at that coefficient level requires an explicit coproduct and
multiplication in the category of spherical algebras; it is not obtained by
silently replacing $\Fone$ with $\B$.  The present paper therefore uses
$\Fone[T^{H_+}]$ for Frobenius-root stalks and $\Spl(\Nbar)$ for the exact
semiringed split arithmetic site.
\end{warning}

\section{Logical dependencies and scope of the claims}
\label{sec:claim-boundaries}

The results of this paper separate into six logical layers.

\begin{enumerate}[label=\textup{(\arabic*)},leftmargin=3.2em]
\item The split-zero algebra, Boolean fibre-product theorem, fixed-locus
formula, localization results, projective decompositions, projective
semimodule constraints, conservative ideal arithmetic, monad formulas,
surcomplex valuation results, and internal split-site construction are
proved from the definitions and the two split-zero source notes
\cite{Globalization,Secondary}.

\item The distributive-lattice globalization, its universal ring reflection,
principal-interval localizations, chain and nilpotent-thickening models,
and the support-associator separation theorem are proved directly in
Sections~\ref{sec:lattice-support}--\ref{sec:nonassociative-support}.
The Cayley--Dickson threshold uses the classical division-algebra and
zero-divisor results in
\cite{BaezOctonions,MorenoCD,ImaedaCD,WilmotCD}.  The tensor-category
interpretation of a $3$-cocycle associator follows the standard framework
of \cite{BHMTensor}.

\item A stronger classification of all finitely generated projective
$G(R)$-semimodules by independently split decorated posets would require a
coherent simultaneous splitting of all transport maps.  No such
classification is asserted.  Theorem
\ref{thm:projective-support-structure} records the consequences supplied by
the retract argument: finite distributive support, projective fibres, and
split-monic transport.

\item The absolute $\Fone$-site, its points and stalks, the map $\Theta$,
the place sets $\Sigma_H$, local torsors, periodic orbits, and scaling-site
geometry are imported from Connes--Consani
\cite{CCAbsolute,CCArithmeticSite,CCScalingSite}.  Their arithmetic Picard,
Jacobian, framed-divisor, rooted-divisor, and character-dual theorems are
imported from \cite{CCJacobian}.  The split-product, quotient, denominator,
and phase-complementarity statements are formal consequences proved here
from those theorems.

\item The finite and global carry extensions, their Ext classes, the carry
cocycles, the cofinality theorem, and the failure of torsion to commute with
the inverse limit are proved in Sections
\ref{sec:lcm-perfection}--\ref{sec:perfected-anomaly}.  The identification
$\Hom(\Q,\Q/\Z)\cong\mathbb A_f$ and the rooted-torsion formula are imported
from \cite{CCJacobian}.  The identification of
$\widehat\Z^{\times}$ with the Bost--Connes symmetry group and with
$\Gal(\Q^{\mathrm{ab}}/\Q)$ is classical and is used in the form recorded
in \cite{BostConnes,CCBCCyclotomy,CMR}.  No new $C^*$-dynamical system is
constructed.

\item The helix paper and Lean development \cite{Lavery} supply a formalized
extrinsic carrier and its multiplier, unitary block, determinant, and
analytic phasor statements.  Section~\ref{sec:helix-local-system} extracts
the factor-of-automorphy content and combines it with the internally proved
fixed-locus theorem.  Neither source proves RH.

\item The finite quadratic-gerbe theory in
Section~\ref{sec:gerbe-quantization} is the FHLT finite path-integral
theorem \cite{FHLT} applied to the explicitly proved forms $q_n$.  The
bundle-gerbe transgression and theta-module realization are imported from
\cite{GerbeTheta}; their identification with the finite Schr\"odinger
module is proved here.  The Lagrangian denominator correspondences, UHF
limit, ordered $K$-theory, profinite diagonal, and strict carry obstruction
are proved from the displayed constructions.  The UHF classification uses
the standard theory of \cite{GlimmUHF}.  No fully extended
infinite-level TQFT is asserted.

\item The dyadic-to-integer note \cite{Schultze} supplies the reciprocal
Beurling refinement problem.  Theorem
\ref{thm:direct-beurling-domination} strengthens its decrement estimate by
the exact inclusion $\cR_Q\subseteq V_{\Lambda_Q}$, but the remaining
density endpoint is precisely the classical RH-equivalent
Nyman--B\'aez-Duarte statement.
\end{enumerate}

The Connes--Consani spectral-triple and real-zero constructions provide a
natural analytic interface, but no mixed-support spectral convergence
theorem is asserted without an explicit realization map into their
quadratic-form spaces.  The perfected carry extension is not identified
with a QFT anomaly by raw equality.  Instead, Sections
\ref{sec:gerbe-quantization}--\ref{sec:uhf-history} construct the actual
morphism chain: rigidified phase group, quadratic refinement,
Eilenberg--Mac Lane delooping, finite gerbe theory, surface transgression,
Heisenberg quantization, theta realization, Lagrangian correspondence, and
UHF strictification.  The strictification exposes, rather than suppresses,
the arithmetic carry obstruction to cyclotomic equivariance.  Julia-line
value distribution and ghost-characteristic calculi remain additional
analytic models; a further identification with either would require an
equivariant realization of its dynamics or differential complex.

\section{Final synthesis}
\label{sec:final-synthesis}

\begin{theorem}[Split-support arithmetic-geometric synthesis]
\label{thm:final-synthesis-v5}
Let $R$ be a nonzero commutative ring, let $S=G(R)$, and let
$e=0_R\neq\tau=0_S$.  The following statements hold.
\begin{enumerate}[label=\textup{(\arabic*)},leftmargin=3.2em]
\item The globalization of a product is the synchronized-support fibre
product
\[
G(A\times B)\cong G(A)\times_{\B}G(B),
\]
while the full Cartesian product contains two additional mixed faces.
For $u\in R^\times$,
$\operatorname{Fix}_{G(R)}(u)=\{\tau\}\sqcup\Ann_R(u-1)$.

\item Localization at ordinary nonzero amplitudes preserves split
arithmetic; localization at $e$ gives $\B$; localization at $\tau$ gives
$\zero$.  Mixed localization preserves, Booleanizes, or erases a coordinate
according as its denominator is nonzero supported, supported zero, or
absent.

\item For every bounded distributive lattice $L$, the semiring $G_L(R)$ has
universal ring reflection $p_L$, ring-null fibre $L$, and localization
\[
G_L(R)[z_\lambda^{-1}]\cong\downarrow\lambda.
\]
Boolean lattices give mixed-support cubes, chains give iterated sub-zero
towers, and the same chains occur as the support semirings of
$\Z/p^n\Z$ and $k[\epsilon]/(\epsilon^n)$.  Their ordinary ideal products
instead give the truncated valuation semiring $V_n$.

\item For a nonassociative coefficient algebra, lattice support preserves
the associator and turns an intrinsic zero-divisor product into supported
zero.  Signed graded multiplication has basis associator
$\phi_F=\delta F$.  In the real Cayley--Dickson tower, nonzero supported
zero-divisor products occur first at the sedenion stage.

\item Every finitely generated projective $G(R)$-semimodule has a finite
distributive support lattice, finitely generated projective $R$-module
fibres, and split-monic transport maps.

\item For a Dedekind domain, supported ideal multiplication, the normed
ideal class group, and the Dedekind zeta function are classical.  The
support-separated quotient-size weight $N(\mathfrak a)+1$ is
nonmultiplicative and its Dirichlet series has no ordinary Euler product.

\item For a field $K$, split projective space is the disjoint union of the
classical projective spaces over its nonempty support faces, and its linear
automorphism group is monomial.

\item Residue maps produce Boolean cubes of ordinary Leibniz monads.  In the
surcomplex realization the local invariants are support, valuation order,
and complex leading phase.  Projective inversion identifies the structured
zero and infinity monads without identifying their centres arithmetically.

\item The split point chain
\[
P_\tau\subset P_0\subset P_p
\]
resolves the absolute point chain
\[
\{0\}\hookrightarrow\Z\hookrightarrow\Z[1/p],
\]
while the arithmetic Abel--Jacobi map identifies the first two terms at the
finite divisor-class level.

\item The completed arithmetic Jacobian has the exact split-support form
\[
\Jacar\cong\Picf\times\B,
\qquad
\Idem(\Jacar)\cong
\mathcal P(\mathbb P_f\sqcup\{\infty\}).
\]
Its quotient obtained by collapsing all zero-seminorm points is the pointed
split monoid $\Picf^\tau$.  The Boolean coordinate is the support
truncation of the archimedean seminorm and therefore records degeneracy,
not its metric magnitude.

\item Finite denominator groups become distinct after rigidification.
Tensor product obeys $H_Q\otimes H_M\cong H_{QM}$, whereas root refinement
obeys $H_Q+H_M=H_{\operatorname{lcm}(Q,M)}$.  Stalk transitions along the
latter divisibility system are universal Frobenius maps.

\item The unique globally Frobenius-perfect rigidified point is $\Q$.
The LCM tower converges to its stalk $\Fone[T^{\Q_+}]$, and all finite and
archimedean places emerge only at this limit.

\item The rigidification $\Z\subset\Q$ carries the nonzero class
\[
[\mathcal E_\Omega]=1_{\widehat\Z}
\in\Ext^1_{\Z}(\Q/\Z,\Z),
\]
represented by the global carry cocycle.  Its character dual is the
nonsplit adelic history extension
\[
0\to\widehat\Z\to\mathbb A_f\to\Q/\Z\to0.
\]
Finite character duals are entirely torsion, while their inverse limit is
the torsion-free group $\mathbb A_f$.  For $H_S=\Z[S^{-1}]$, relative
phase and rooted phase are complementary primary summands of $\Q/\Z$.
The phase automorphism group is $\widehat\Z^\times$.

\item For each denominator $n$, the explicit nondegenerate quadratic form
$q_n$ on $\Phi_n$ defines an invertible FHLT gerbe anomaly.  Torus
transgression gives a finite Heisenberg algebra
\[
\C_{\omega_n}[\Phi_n^2]\cong M_n(\C),
\]
and the winding-$n$ theta sector of the level-one $T^3$ gerbe is its
irreducible module.  Frobenius multiplication is loop-winding dilation.

\item Denominator divisibility is functorial through cyclotomic-equivariant
Lagrangian correspondences.  The strict matrix tower has limit
$\mathcal Q_\Omega$ with
\[
K_0(\mathcal Q_\Omega)\cong(\Q,\Q_{\geq0},1),
\qquad
\mathcal D_\Omega\cong C(\widehat\Z)\cong C^*(\Q/\Z).
\]
The cyclotomic action is strict on the correspondence system and the
diagonal.  Its failure to commute with the standard full-matrix embeddings
is exactly the crossed carry cocycle $c_{u;n,m}$.  Split support distinguishes
an absent anomalous tier from a defined vanishing section.

\item The reciprocal-root skeleton of the perfect stalk has Beurling
realization $T^{1/n}\mapsto F_{1/n}$ and satisfies
$x_{\Lambda_Q}\leq d_Q^2$.  Its density at $1$ is equivalent to RH by the
classical Nyman--B\'aez-Duarte theorem.

\item The helix multiplier is the factor of automorphy of a Frobenius line
on $C_p$.  Its invariant split fibre is Boolean unless the holonomy is
trivial, exactly when it jumps to the full $G(\C)$ fibre.

\item A split local Weil torsor has quotient groupoid
$BW_v\sqcup BW_v\sqcup *$.  Finite sets of places therefore carry Boolean
support cubes refined by neutral and nontrivial local states.

\item The internal split construction produces the exact semiringed
arithmetic site
\[
(\widehat{\N^\times},\Spl(\Nbar))
\]
and the split scaling site.  It does not identify the spherical coefficient
$\Fone$ with $\B$.
\end{enumerate}
\end{theorem}

\begin{proof}
Item (1) is Theorems~\ref{thm:boolean-fibre-product} and
\ref{thm:fixed-locus-unit}.  Item (2) is
Theorems~\ref{thm:division-trichotomy} and
\ref{thm:mixed-localization}.  Item (3) is
Theorems~\ref{thm:lattice-normal-form},
\ref{thm:lattice-localization}, and
\ref{thm:nilpotent-chain-realization}, together with
Proposition~\ref{prop:truncated-valuation-semiring}.  Item (4) is
Theorems~\ref{thm:support-associator-separation},
\ref{thm:associator-cocycle}, and
\ref{thm:cayley-dickson-support-threshold}.

Item (5) is Theorem~\ref{thm:projective-support-structure}.  Item (6) is
Theorem~\ref{thm:ideal-multiplication-classical},
Corollary~\ref{cor:dedekind-preservation}, and
Theorem~\ref{thm:separated-weight}.  Item (7) is
Theorems~\ref{thm:projective-face-decomposition} and
\ref{thm:monomial-gl}.  Item (8) is
Theorems~\ref{thm:cubical-monad},
\ref{thm:surcomplex-standard-part},
\ref{thm:complex-graded-layers}, and
\ref{thm:univalent-interpretation}.

Item (9) is Theorem~\ref{thm:abel-jacobi-resolution}.  Item (10) is
Theorem~\ref{thm:jacobian-boolean-product},
Proposition~\ref{prop:jacobian-reduction}, and
Proposition~\ref{prop:arakelov-support-truncation}.  Item (11) is
Theorems~\ref{thm:rigidified-divisibility},
\ref{thm:two-denominator-laws}, and
\ref{thm:denominator-frobenius}.  Item (12) is
Theorems~\ref{thm:unique-frobenius-perfect-point},
\ref{thm:global-perfection}, and \ref{thm:places-emerge}.  Item (13) is
Theorems~\ref{thm:global-carry-class},
\ref{thm:adelic-history-extension},
\ref{thm:torsion-transmutation},
\ref{thm:phase-complementarity}, and
\ref{thm:cyclotomic-symmetry-anomaly}.

Item (14) is Theorems~\ref{thm:cyclic-quadratic-refinement},
\ref{thm:finite-gerbe-anomaly},
\ref{thm:finite-stone-von-neumann},
\ref{thm:winding-theta-realization}, and
\ref{thm:frobenius-winding-dilation}.  Item (15) is
Theorems~\ref{thm:lagrangian-denominator-history},
\ref{thm:perfected-uhf}, \ref{thm:profinite-diagonal}, and
\ref{thm:cyclotomic-carry-obstruction}, together with
Proposition~\ref{prop:split-anomaly-line}.  Item (16) is
Theorems~\ref{thm:direct-beurling-domination} and
\ref{thm:arithmetic-site-nyman}.  Item (17) is
Theorem~\ref{thm:helix-holonomy},
Proposition~\ref{prop:helix-descent}, and
Corollary~\ref{cor:helix-fixed-locus-jump}.  Item (18) is
Theorems~\ref{thm:split-torsor-groupoid} and
\ref{thm:boolean-place-cubes}.  Item (19) is
Theorems~\ref{thm:internal-split-functor},
\ref{thm:split-arithmetic-site}, and
\ref{thm:three-scaling-states}.
\end{proof}

\begin{remark}[Geometric meaning]
The full local invariant is a hierarchy rather than a single Boolean tag:
\[
\begin{gathered}
\text{prime-configuration support}
\times
\text{archimedean support}
\times
\text{relative denominator phase}
\\
\times
\text{rooted cyclotomic phase}
\times
\text{adelic phase history}
\times
\text{local amplitude or valuation data}.
\end{gathered}
\]
The first two factors occur in the idempotent stratum of the arithmetic
Jacobian.  Relative and rooted phases are related by Theorem
\ref{thm:phase-complementarity}; their coherent transition history is
$\mathbb A_f$.  Split zero is the reduced algebraic shadow obtained when
unsupported or degenerate labels are collapsed to one absorber and phase
history is forgotten.  The unreduced arithmetic and rooted-divisor
geometries retain which object carried the degeneration and how it arose
through the denominator tower.
\end{remark}

\begin{thebibliography}{99}

\bibitem{Globalization}
ChatGPT 5.4 and Gemini 2.5 DeepThink,
\emph{A Note on the Globalization Semiring $G(\Z)$},
Zenodo, 2026.
DOI: \href{https://doi.org/10.5281/zenodo.18976982}{10.5281/zenodo.18976982}.

\bibitem{Secondary}
The Clankers,
\emph{A Secondary Note on Split-Zero Globalization: Semimodule
Classification, Multiplicative Reconstruction, and All-Orders Zeta
Deformation},
Zenodo, 2026.
DOI: \href{https://doi.org/10.5281/zenodo.19120905}{10.5281/zenodo.19120905}.

\bibitem{Homogenization}
The Clankers,
\emph{Support-Idempotent Homogenization for a Split-Zero Collatz F-Series
Frame},
Zenodo, 2026.
DOI: \href{https://doi.org/10.5281/zenodo.19900461}{10.5281/zenodo.19900461}.

\bibitem{CCAbsolute}
A.~Connes and C.~Consani,
\emph{On the Absolute Geometry of $\operatorname{Spec}\Z$},
arXiv:2606.06604 [math.AG], 2026.
DOI: \href{https://doi.org/10.48550/arXiv.2606.06604}{10.48550/arXiv.2606.06604}.

\bibitem{CCArithmeticSite}
A.~Connes and C.~Consani,
\emph{Geometry of the Arithmetic Site},
Adv. Math. \textbf{291} (2016), 274--329.

\bibitem{CCScalingSite}
A.~Connes and C.~Consani,
\emph{Geometry of the Scaling Site},
Selecta Math. (N.S.) \textbf{23} (2017), no.~3, 1803--1850.


\bibitem{CCJacobian}
A.~Connes and C.~Consani,
\emph{On the Jacobian of $\operatorname{Spec}\Z$},
arXiv:2602.15941 [math.NT], 2026.
DOI: \href{https://doi.org/10.48550/arXiv.2602.15941}{10.48550/arXiv.2602.15941}.

\bibitem{CCAbsoluteAlgebra}
A.~Connes and C.~Consani,
\emph{Absolute algebra and Segal's $\Gamma$-rings: au-dessous de
$\operatorname{Spec}(\Z)$},
J. Number Theory \textbf{162} (2016), 518--551.

\bibitem{CCAffineAbsolute}
A.~Connes and C.~Consani,
\emph{On absolute algebraic geometry: the affine case},
Adv. Math. \textbf{390} (2021), Article 107909.

\bibitem{CCGammaUniversal}
A.~Connes and C.~Consani,
\emph{Segal's gamma rings and universal arithmetic},
Q. J. Math. \textbf{72} (2021), 1--29.


\bibitem{BostConnes}
J.-B.~Bost and A.~Connes,
\emph{Hecke algebras, type III factors and phase transitions with
spontaneous symmetry breaking in number theory},
Selecta Math. (N.S.) \textbf{1} (1995), no.~3, 411--457.

\bibitem{CCBCCyclotomy}
A.~Connes and C.~Consani,
\emph{BC-system, absolute cyclotomy and the quantized calculus},
EMS Surv. Math. Sci. \textbf{9} (2022), 447--475.
DOI: \href{https://doi.org/10.4171/EMSS/64}{10.4171/EMSS/64}.

\bibitem{CMR}
A.~Connes, M.~Marcolli, and N.~Ramachandran,
\emph{KMS states and complex multiplication},
Selecta Math. (N.S.) \textbf{11} (2005), no.~3--4, 325--347.

\bibitem{ScholzePerfectoid}
P.~Scholze,
\emph{Perfectoid spaces},
Publ. Math. Inst. Hautes Etudes Sci. \textbf{116} (2012), 245--313.

\bibitem{Weibel}
C.~A.~Weibel,
\emph{An Introduction to Homological Algebra},
Cambridge Studies in Advanced Mathematics \textbf{38},
Cambridge University Press, Cambridge, 1994.

\bibitem{FarguesFontaine}
L.~Fargues and J.-M.~Fontaine,
\emph{Courbes et fibr\'es vectoriels en th\'eorie de Hodge $p$-adique},
Ast\'erisque \textbf{406} (2018), xiii+382 pp.

\bibitem{Lavery}
S.~Lavery,
\emph{The Double-Ended Helix: An Intrinsically On-Axis Geometric State
Space for Dirichlet $L$-Function Phasors},
manuscript with Lean~4 formalization, 2026,
\url{https://github.com/samlavery/helix_frobenius}.

\bibitem{Schultze}
K.~Schultze,
\emph{A Dyadic-to-Integer Beurling Reduction for a Reciprocal
Nyman--Beurling Refinement},
preliminary research note, v0.1.2, 2026,
\url{https://github.com/KarlSchultze/dyadic-beurling-reduction}.

\bibitem{Nyman}
B.~Nyman,
\emph{On the One-Dimensional Translation Group and Semi-Group in Certain
Function Spaces},
Ph.D. thesis, Uppsala University, 1950.

\bibitem{Beurling}
A.~Beurling,
\emph{A closure problem related to the Riemann zeta-function},
Proc. Nat. Acad. Sci. U.S.A. \textbf{41} (1955), 312--314.

\bibitem{BaezDuarte}
L.~B\'aez-Duarte,
\emph{A strengthening of the Nyman--Beurling criterion for the Riemann
hypothesis},
Rend. Mat. Acc. Lincei \textbf{14} (2003), 5--11.

\bibitem{Steinitz}
E.~Steinitz,
\emph{Algebraische Theorie der K\"orper},
J. Reine Angew. Math. \textbf{137} (1910), 167--309.

\bibitem{BerarducciMantova}
A.~Berarducci and V.~Mantova,
\emph{Surreal numbers, derivations and transseries},
J. Eur. Math. Soc. \textbf{20} (2018), no.~2, 339--390.
DOI: \href{https://doi.org/10.4171/JEMS/769}{10.4171/JEMS/769}.

\bibitem{Conway}
J.~H.~Conway,
\emph{On Numbers and Games},
2nd ed., A K Peters, Natick, MA, 2001.

\bibitem{Deitmar}
A.~Deitmar,
\emph{Schemes over $\Fone$},
in \emph{Number Fields and Function Fields---Two Parallel Worlds},
Progr. Math. \textbf{239}, Birkh\"auser, Boston, 2005, 87--100.

\bibitem{Golan}
J.~S.~Golan,
\emph{Semirings and Their Applications},
Kluwer Academic Publishers, Dordrecht, 1999.

\bibitem{Goldblatt}
R.~Goldblatt,
\emph{Lectures on the Hyperreals: An Introduction to Nonstandard Analysis},
Graduate Texts in Mathematics \textbf{188}, Springer, New York, 1998.

\bibitem{GrothendieckEGAII}
A.~Grothendieck,
\emph{\'El\'ements de g\'eom\'etrie alg\'ebrique II: \'Etude globale
\'el\'ementaire de quelques classes de morphismes},
Publ. Math. IH\'ES \textbf{8} (1961), 5--222.

\bibitem{HoTT}
The Univalent Foundations Program,
\emph{Homotopy Type Theory: Univalent Foundations of Mathematics},
Institute for Advanced Study, 2013,
\url{https://homotopytypetheory.org/book/}.

\bibitem{Lorscheid}
O.~Lorscheid,
\emph{The geometry of blueprints. Part I: Algebraic background and scheme
theory},
Adv. Math. \textbf{229} (2012), no.~3, 1804--1846.

\bibitem{Robinson}
A.~Robinson,
\emph{Non-standard Analysis},
North-Holland, Amsterdam, 1966.

\bibitem{Simons}
J.~Simons,
\emph{Meet the Surreal Numbers},
expository notes, 4 April 2017.


\bibitem{FHLT}
D.~S.~Freed, M.~J.~Hopkins, J.~Lurie, and C.~Teleman,
\emph{Topological quantum field theories from compact Lie groups},
arXiv:0905.0731v2, 2009.

\bibitem{GerbeTheta}
The Clankers,
\emph{Coupled Gerbes, Transgression, Explicit Theta Sectors, Hesse Cubics,
and Landau/BCM Structures on $T^3$},
pre-publication manuscript, March 2026.

\bibitem{BaezOctonions}
J.~C.~Baez,
\emph{The octonions},
Bull. Amer. Math. Soc. (N.S.) \textbf{39} (2002), 145--205.

\bibitem{MorenoCD}
G.~Moreno,
\emph{The zero divisors of the Cayley--Dickson algebras over the real
numbers},
Bol. Soc. Mat. Mexicana (3) \textbf{4} (1998), 13--28.

\bibitem{ImaedaCD}
K.~Imaeda and M.~Imaeda,
\emph{Sedenions: algebra and analysis},
Appl. Math. Comput. \textbf{115} (2000), 77--88.

\bibitem{WilmotCD}
G.~P.~Wilmot,
\emph{Structure of the Cayley--Dickson algebras},
arXiv:2505.11747v1, 2025.

\bibitem{BHMTensor}
P.~Bouwknegt, K.~C.~Hannabuss, and V.~Mathai,
\emph{$C^*$-algebras in tensor categories},
arXiv:math/0702802v2, 2009.

\bibitem{GlimmUHF}
J.~Glimm,
\emph{On a certain class of operator algebras},
Trans. Amer. Math. Soc. \textbf{95} (1960), 318--340.

\end{thebibliography}

\end{document}
